\documentclass{amsart}
\usepackage{amsmath,amssymb,mathtools,array,longtable,booktabs,geometry,multicol}
\usepackage[T1]{fontenc}
\usepackage[hidelinks]{hyperref}
\newenvironment{VChainTable}[6]{%
\begingroup\normalsize\normalfont%
\begin{longtable}{@{}rrrrrrrr@{}}%
\multicolumn{8}{c}{$q=#1$; $n=#3$; $#2$; $start=#4$; $corr=(#5,#6)$}\\
\hline
$\ell_i$ & $\Delta m_i$ & $\ell_i$ & $\Delta m_i$ & $\ell_i$ & $\Delta m_i$ & $\ell_i$ & $\Delta m_i$\\
\hline\endfirsthead
\multicolumn{8}{c}{$q=#1$ (continued)}\\
\hline
$\ell_i$ & $\Delta m_i$ & $\ell_i$ & $\Delta m_i$ & $\ell_i$ & $\Delta m_i$ & $\ell_i$ & $\Delta m_i$\\
\hline\endhead
}{\end{longtable}\endgroup}
\newenvironment{VPoissonTable}[3]{%
\begingroup\normalsize\normalfont%
\begin{longtable}{@{}rrrrrrrrr@{}}%
\multicolumn{9}{c}{Poisson $q=#1$; #2 rows, #3 prefix triples}\\
\hline
$\Delta x_i$ & $\Delta y_i$ & $\Delta m_i$ & $\Delta x_i$ & $\Delta y_i$ & $\Delta m_i$ & $\Delta x_i$ & $\Delta y_i$ & $\Delta m_i$\\
\hline\endfirsthead
\multicolumn{9}{c}{Poisson $q=#1$ (continued)}\\
\hline
$\Delta x_i$ & $\Delta y_i$ & $\Delta m_i$ & $\Delta x_i$ & $\Delta y_i$ & $\Delta m_i$ & $\Delta x_i$ & $\Delta y_i$ & $\Delta m_i$\\
\hline\endhead
}{\end{longtable}\endgroup}
\geometry{paperwidth=46pc,paperheight=62pc,textwidth=38pc,textheight=54pc}
\setlength{\emergencystretch}{3em}
\newtheorem{theorem}{Theorem}[section]
\newtheorem{lemma}[theorem]{Lemma}
\newtheorem{proposition}[theorem]{Proposition}
\newtheorem{remark}[theorem]{Remark}
\numberwithin{equation}{section}
\numberwithin{table}{section}
\newcommand{\Li}{\operatorname{Li}}
\newcommand{\Z}{\mathbb Z}
\newcommand{\Q}{\mathbb Q}
\newcommand{\R}{\mathbb R}
\newcommand{\N}{\mathbb N}
\newcommand{\eps}{\varepsilon}

\begin{document}
\title[Reciprocal dilogarithms]{Irrationality of the dilogarithm at reciprocals of integers}
\author{Ruiyi Zhang}
\address{Independent Researcher, Columbia, Maryland, USA}
\email{ryzfsu@gmail.com}
\subjclass[2020]{Primary 11J72; Secondary 33B30}
\keywords{dilogarithm, determinant, logarithmic energy, Hermite--Pad\'e, Frobenius}

\begin{abstract}
We prove that \(\Li_2(1/n)\) is irrational for every nonzero integer \(n\).
The proof combines a determinant criterion with exact fixed-base
integerization, logarithmic-potential estimates, Frobenius nonvanishing, and a
three-form Hermite--Pad\'e construction.  The finite positive range is treated
by a common variational problem and two endpoint certificates; the negative
middle range is handled by signed contour integrals at the six bases
\(-2,\ldots,-7\); the two remaining tails follow from an explicit
Hermite--Pad\'e estimate.  Every numerical constant used in the argument is
given as a rational bound or an explicitly directed interval.
\end{abstract}
\maketitle

\section{Introduction}
For an integer $s$ with $|s|\ge2$, irrationality of $\Li_2(1/s)$ was proved
by Hata~\cite{Hata} for all $s\le-5$ and $s\ge7$, and by Rhin and
Viola~\cite{RhinViola} for $s=6$.  A recent preprint of
Prellberg~\cite{Prellberg} treats $s=-4,-3$ and $5$.  These works are cited
for history only: no irrationality statement from them is used below, and
the proof is self-contained.  The cases $s=-2,2,3,4$ were open; among them
is $\Li_2(1/2)=\pi^2/12-(\log2)^2/2$.  For $s=\pm1$ the values are rational
multiples of $\pi^2$, which is irrational~\cite{Legendre}.

The general tools are the prime number theorem~\cite{Davenport}; Kummer's
carry theorem~\cite{Kummer} and Lucas' congruence~\cite{Lucas}; the
Cauchy--Binet formula~\cite{Gantmacher} and its integral form due to
Andr\'eief~\cite{Andreief,Forrester}; Cauchy's double-alternant determinant
~\cite{Krattenthaler}; logarithmic potential theory~\cite{SaffTotik}; and
Stirling's formula~\cite{WhittakerWatson}.

We call the positive determinant P0 and the signed contour
determinant V.  These names refer to the two explicit constructions below.

The endpoint energy lower bounds and fixed-base field bounds are
computer-assisted: directed interval arithmetic verifies them from the
printed tables, with the algorithm and code supplied in the ancillary files.

The proof applies the criterion below to one defined determinant at a time.
Sections~\ref{sec:p0-arithmetic}--\ref{sec:p0-analytic} treat positive bases;
Sections~\ref{sec:v-arithmetic}--\ref{sec:v-analytic} treat the negative
bases; Section~\ref{sec:hp} treats both tails, and
Section~\ref{sec:nonvanishing} supplies eventual nonvanishing sequences.  The final
assembly checks the ranges and their boundary points.

\section{Main theorem}
For \(|z|\le1\), put
\[
  \Li_2(z)=\sum_{r\ge1}\frac{z^r}{r^2}.
\]
\begin{theorem}\label{thm:main}
For every nonzero integer $n$, the number $\Li_2(1/n)$ is irrational.
\end{theorem}

For $n=1$ and $n=-1$ this follows from
\(\Li_2(1)=\pi^2/6\) and \(\Li_2(-1)=-\pi^2/12\), together with the
irrationality of $\pi^2$~\cite{Legendre}; the determinant argument covers all
$|n|\ge2$.

\subsection{Coverage map}
The ranges below are disjoint and cover all remaining nonzero integers.
\begin{center}\normalsize
\begin{tabular}{@{}p{0.20\linewidth}p{0.28\linewidth}p{0.40\linewidth}@{}}
\toprule
base range & construction & ingredients\\
\midrule
\(\pm1\) & classical evaluations & $\pi^2$\\
\(2\le a\le11\) & positive determinant & arithmetic, energy, nonvanishing\\
\(-7\le a\le-2\) & signed determinant & arithmetic, contour energy, nonvanishing\\
\(a\ge12\) or \(a\le-8\) & three-form Hermite--Pad\'e & residues, integerization, saddle estimates\\
\bottomrule
\end{tabular}
\end{center}
In each application of the criterion, the determinant, integerizer, analytic
estimate, and nonvanishing sequence belong to the same polynomial family at the
indicated base.

\section{The determinant criterion}\label{sec:criterion}
Let \((m_k)_{k\ge1}\) be a strictly increasing sequence of positive integers,
and let \(G_{m_k,a}(T)\in\Q[T]\) have degree at most \(m_k\).  Suppose that,
for every fixed integer \(a\) under consideration, there are positive
rational numbers \(D_{m_k,a}\)
 and a fixed real target \(t\) such that
\begin{align}
 D_{m_k,a}G_{m_k,a}(T)&\in\Z[T], &
 \limsup_{k\to\infty}m_k^{-2}\log D_{m_k,a}&\le\tau(a), \label{eq:criterion-int}\\
 \log|G_{m_k,a}(t)|&\le-\alpha(a)m_k^2+O(m_k\log m_k), \label{eq:criterion-small}
\end{align}
and, for every rational \(\xi\), there is \(k_0=k_0(a,\xi)\) such that
\[
 G_{m_k,a}(\xi)\ne0\qquad(k\ge k_0).
\]
If
\(\alpha(a)>\tau(a)\), then \(t\) is irrational.

Indeed, if \(t=\xi=u/v\in\Q\) in lowest terms, then for every sufficiently
large \(k\),
\(v^{m_k}D_{m_k,a}G_{m_k,a}(\xi)\) is a nonzero integer, while
\eqref{eq:criterion-int} and \eqref{eq:criterion-small} make its absolute
value tend to zero.  The determinant, integerizer and nonvanishing statements
must refer to the same polynomial.

\section{Positive determinant}\label{sec:p0}
For \(g=m/2\), let \(I_m=\{m/4,\ldots,5m/4-1\}\) and
\(J_m=\{0,\ldots,2m-1\}\).  Set
\[
 U_i(a)=a^i\binom{i+g}{g},\quad
 t_i(a)=\sum_{r=1}^i\frac{U_{i-r}(a)}r
 +a^i\sum_{s=1}^g\frac{(-1)^{s-1}}s
       \binom{i+g}{g-s}\left(\left(\frac a{a-1}\right)^s-1\right).
\]
For \(i\ne j\), define
\[
 C_{ij}(a)=\frac{U_i(a)-U_j(a)}{j-i},\qquad
 B_{ij}(a)=\frac{t_i(a)-t_j(a)}{j-i},
\]
and put \(A_{ij}(a)=0\).  On the diagonal,
\[
 A_{ii}(a)=U_i(a),\quad C_{ii}(a)=-U_i(a)\sum_{s=1}^g\frac1{i+s},
\]
\[
 B_{ii}(a)=\sum_{r=1}^i\frac{U_{i-r}(a)}{r^2}
 -a^i\sum_{s=2}^g\frac{(-1)^s}{s}\binom{i+g}{g-s}
       \sum_{r=1}^{s-1}\frac{(a/(a-1))^r-1}{r}.
\]
The P0 determinant is
\begin{equation}\label{eq:p0det}
 G_{m,a}(T)=\det\begin{bmatrix}C(a)_{I_m,J_m}\\
                      B(a)_{I_m,J_m}-T A(a)_{I_m,J_m}\end{bmatrix}.
\end{equation}
\section{Arithmetic for the positive determinant}\label{sec:p0-arithmetic}

This section proves the fixed-base denominator estimate for the P0
pencil.  The base is fixed before the limit in $m$ is taken.  No estimate
uniform in unbounded $|a|$ at finite $m$ is asserted.  The proof has an
algebraic part, an exact sum over primes, and a prime-number-theorem limit.
The algebraic and limiting arguments below are all-$m$ arguments.  The
large-ratio contribution is specified by the rational table
\ref{tab:p0-high-potentials} and its type-generation rule.

Write $m=4h$, $g=2h$, $I=\{h,\ldots,5h-1\}$,
$J=\{0,\ldots,8h-1\}$, and $U=8h-1=2m-1$.  Let
\[
 \begin{aligned}
 d_U&=\operatorname{lcm}(1,2,\ldots,U),\\
 N_m&=13h^2-4h=\frac{13}{16}m^2-m,\\
 H_m&=h^2-3h=\frac1{16}m^2-\frac34m.
 \end{aligned}
\]
For a prime $p$, put $d_p=\lfloor\log_p U\rfloor$ and use $v_p$ for the
usual valuation, including its coefficientwise extension to polynomials.
The parameter-dependent entries $A(a),B(a),C(a)$ and determinant
$G_{m,a}(T)$ are those of \eqref{eq:p0det}; in particular the diagonal
conventions in Section~\ref{sec:p0} are retained.

\subsection{The source Mellin identity}
For real $a$ with $|a|>1$ put
\[
 t_i^0(a)=\sum_{r=1}^{i}\frac{a^{i-r}}r,
 \quad b_i^0(a)=\sum_{r=1}^{i}\frac{a^{i-r}}{r^2},
 \quad L(a)=-\sum_{r\ge1}\frac{a^{-r}}r,
 \quad \beta(a)=\sum_{r\ge1}\frac{a^{-r}}{r^2}.
\]
Thus $L(a)=\log(1-1/a)$ and $\beta(a)=\Li_2(1/a)$ with their real
values on both components of $\{|a|>1\}$.  Define
\[
 A^0_{ij}(a)=a^i\delta_{ij},\qquad
 C^0_{ij}(a)=\begin{cases}(a^i-a^j)/(j-i),&i\ne j,\\0,&i=j,\end{cases}
\]
and
\[
 B^0_{ij}(a)=\begin{cases}(t_i^0(a)-t_j^0(a))/(j-i),&i\ne j,\\
 b_i^0(a),&i=j.\end{cases}
\]

\begin{lemma}[source Mellin identity]
\label{lem:p0-source-mellin}
For $|a|>1$ and all $i,j\ge0$,
\begin{equation}\label{eq:p0-source-mellin}
 M_{ij}(a):=B^0_{ij}(a)+L(a)C^0_{ij}(a)-\beta(a)A^0_{ij}(a)
 =-\sum_{r\ge1}\frac{a^{-r}}{(i+r)(j+r)}
 =-\int_0^1\!\!\int_0^1\frac{x^iy^j}{a-xy}\,dx\,dy.
\end{equation}
\end{lemma}
\begin{proof}
The two finite identities
\[
 t_i^0(a)+a^iL(a)=-\sum_{r\ge1}\frac{a^{-r}}{i+r},\qquad
 b_i^0(a)-a^i\beta(a)=-\sum_{r\ge1}\frac{a^{-r}}{(i+r)^2}
\]
follow by separating the first $i$ terms of the defining series.  Taking the
divided difference of the first identity and using
\[
 \frac{(i+r)^{-1}-(j+r)^{-1}}{j-i}
 =\frac1{(i+r)(j+r)}
\]
proves the off-diagonal case of \eqref{eq:p0-source-mellin}; the second
identity proves the diagonal case.  Since $|a|>1$, the geometric series
$(a-xy)^{-1}=\sum_{r\ge1}a^{-r}(xy)^{r-1}$ converges uniformly on
$[0,1]^2$, so termwise integration gives the last equality.
\end{proof}

Put
\[
 D_g f(a)=\frac1{g!}\frac{d^g}{da^g}\bigl(a^gf(a)\bigr).
\]
On every compact subset of $\{|a|>1\}$ the series in
\eqref{eq:p0-source-mellin}, after multiplication by $a^g$, and all its
derivatives through order $g$, are dominated by $C(1+r)^g\rho^r$ with
$\rho<1$.  Termwise differentiation is therefore valid and
\begin{equation}\label{eq:p0-source-Dg}
 D_gM_{ij}(a)=(-1)^{g+1}\sum_{r>g}
 \frac{\binom{r-1}{g}a^{-r}}{(i+r)(j+r)}.
\end{equation}
Indeed $D_ga^{-r}=0$ for $r\le g$ and equals
$(-1)^g\binom{r-1}{g}a^{-r}$ for $r>g$.

For completeness, the right side is exactly the finite rational table used in
the determinant.  If $U_i(a)=a^i\binom{i+g}{g}$ and
$Q_s(a)=(a/(a-1))^s-1$, the finite definitions of $A,B,C$ in
Section~\ref{sec:p0} already use $a$ throughout.  Leibniz's rule gives
$D_ga^i=U_i(a)$ and, for $s\ge1$,
\[
 \frac{a^sL^{(s)}(a)}{s!}=\frac{(-1)^{s-1}}sQ_s(a),\qquad
 \frac{a^s\beta^{(s)}(a)}{s!}=\frac{(-1)^{s-1}}s
 \left(L(a)-\sum_{r=1}^{s-1}\frac{Q_r(a)}r\right).
\]
The latter follows from $\beta'(a)=L(a)/a$.  Substitution in
\eqref{eq:p0-source-mellin} proves
\begin{equation}\label{eq:p0-source-identity}
 D_gM_{ij}(a)=B^{(g)}_{ij}(a)+L(a)C^{(g)}_{ij}(a)
                         -\beta(a)A^{(g)}_{ij}(a).
\end{equation}
For $i=j$, the coefficient of $L(a)$ is
$-U_i(a)\sum_{s=1}^g(i+s)^{-1}$, because
\[
 \sum_{s=1}^g\frac{(-1)^{s-1}}s\binom{i+g}{g-s}
 =\binom{i+g}{g}\sum_{s=1}^g\frac1{i+s}.
\]
This identity is the derivative at $x=0$ of
$\binom{i+g+x}{g}=\sum_{s=0}^g\binom{i+g}{g-s}\binom{x}{s}$;
the coefficient of $-\beta$ is $U_i(a)$.  Hence the diagonal conventions,
and not only the off-diagonal divided differences, are covered by
\eqref{eq:p0-source-identity}.

\subsection{The polynomial divisor}

We first record the two polynomial divisibilities that account for the
exceptional primes.  All statements in this subsection are identities for
every $h\ge1$.

\begin{lemma}[formal divisor at $a=0$]\label{lem:p0-a-divisor}
With $b=a-1$,
\begin{equation}\label{eq:p0-denom}
 d_U^{3m}b^{gm}G_{m,a}(T)\in\mathbb Z[a,T],
 \qquad
 d_U^{3m}b^{gm}G_{m,a}(T)\in a^{N_m}\mathbb Z[a,T].
\end{equation}
\end{lemma}

\begin{proof}
The only parameter denominators in the finite source formulas are powers of
$b$ in
\[
 Q_s(a)=\left(\frac{a}{a-1}\right)^s-1.
\]
After multiplying each lower row by $b^g$, its occurrence becomes the integer
polynomial
\[
 R_s(a)=b^gQ_s(a)=a^s b^{g-s}-b^g\in\mathbb Z[a]
 \quad(1\le s\le g).
\]
Its constant coefficient is $-(-1)^g$, a unit.  The other source
denominators are harmonic indices, summation indices, and divided-difference
denominators.  Every such positive index is at most $U$: for example
$i+g\le7h-1<U$, and $|j-i|\le7h-1<U$.  Multiplication of each top row by
$d_U$ and each lower row by $d_U^2$ clears them.  Thus
\[
 P_m(a,T):=d_U^{3m}b^{gm}G_{m,a}(T)\in\mathbb Z[a,T].
\]

To find its order at $a=0$, work in $\mathbb Q[[a]]$.  Set
\[
 c_i=\binom{i+g}{g},\quad u_i=c_i a^i,\quad
 h_i=\sum_{s=1}^g\frac1{i+s},\quad
 \gamma_i=\frac12\left(h_i^2-\sum_{s=1}^g\frac1{(i+s)^2}\right),
\]
and $L_1=\sum_{r\ge1}a^r/r$, $L_2=\sum_{r\ge1}a^r/r^2$.  The source
identities, coefficientwise in $a$, are
\begin{align}
 t_i-L_1u_i&=\sum_{k\ne i}\frac{c_ka^k}{i-k}-u_i h_i,
 \label{eq:p0-formal-t}\\
 b_i+L_1u_i h_i+L_2u_i
 &=\sum_{k\ne i}\frac{c_ka^k}{(i-k)^2}+\gamma_i u_i.
 \label{eq:p0-formal-b}
\end{align}
Here the $b_i$ on the left is the diagonal source entry, not the parameter
$b=a-1$.  We use a direct coefficient check, so no auxiliary logarithmic
symbol is being differentiated.  For $r\ge0$, expand every finite source
entry in $\mathbb Q[[a]]$, replace
\[
 [a^r]L_1=\frac1r\quad(r\ge1),\qquad
 [a^r]L_2=\frac1{r^2},
\]
and use
\[
 [a^r]u_i=\binom{i+g}{g}[r=i],\qquad
 \gamma_i=\sum_{1\le s<t\le g}\frac1{(i+s)(i+t)}.
\]
The coefficient of the first source identity is then
\[
 [a^r]t_i-[a^r](L_1u_i)=
 \begin{cases}c_r/(i-r),&r\ne i,\\-c_i h_i,&r=i,\end{cases}
\]
which is exactly the right side of \eqref{eq:p0-formal-t}.  For the second
identity, the coefficient is
\[
 [a^r]\bigl(b_i+L_1u_i h_i+L_2u_i\bigr)=
 \begin{cases}c_r/(i-r)^2,&r\ne i,\\c_i\gamma_i,&r=i.\end{cases}
\]
This is the coefficient of its right side, including the diagonal correction.
All sums are finite at a fixed coefficient, so these are identities in
$\mathbb Q[[a]]$ coefficient by coefficient.  The only binomial identity
used in collecting the harmonic terms is the differentiated finite identity
\[
 \binom{i+g+x}{g}=\sum_{s=0}^g\binom{i+g}{g-s}\binom{x}{s}
\]
at $x=0$; hence the argument does not assume a formal primitive for $1/a$.

For an independent indeterminate $X$, form
\[
 \mathcal H(a,X)=
 \begin{bmatrix}
 C(a)\\ B(a)-L_1(a)C(a)+L_2(a)A(a)-XA(a)
 \end{bmatrix}.
\]
For a fixed $k\ge0$, let $e_I(i)=[i=k]$,
$v_I(i)=(i-k)^{-1}$ for $i\ne k$ and $v_I(k)=0$, and
$w_I=v_I^2$; define $e_J,v_J,w_J$ in the same way.  Equations
\eqref{eq:p0-formal-t}--\eqref{eq:p0-formal-b} give the exact rank-one
expansion
\begin{equation}\label{eq:p0-rank-one}
 [a^k]\mathcal H=c_k\left(
 \binom{e_I}{v_I-h_ke_I}v_J^{\mathsf T}
 +\binom{0}{-e_I}w_J^{\mathsf T}
 +\binom{v_I-h_ke_I}{-w_I-h_kv_I+(\gamma_k-X)e_I}e_J^{\mathsf T}
 \right).
\end{equation}
In entries, the lower block of this coefficient is
\[
 [a^k]\bigl(B-L_1C+L_2A-XA\bigr)_{ij}
 =c_k\begin{cases}
 ((i-k)(j-k))^{-1},&i\ne k,\ j\ne k,\\
 -h_k(j-k)^{-1}-(j-k)^{-2},&i=k,\ j\ne k,\\
 -(i-k)^{-2}-h_k(i-k)^{-1},&i\ne k,\ j=k,\\
 \gamma_k-X,&i=j=k.
 \end{cases}
\]
The first feature has weight $k$; the second is present only when $k\in I$,
and the third only when $k\in J$.  A nonzero $2m$-row determinant term must
choose $2m$ distinct rank-one features.  Since $I=[h,5h)$ and $J=[0,8h)$,
the least possible total feature weight is obtained from two copies of
$0,\ldots,h-1$ and three copies of $h,\ldots,3h-1$.  Its value is
\[
 2\sum_{k=0}^{h-1}k+3\sum_{k=h}^{3h-1}k=13h^2-4h=N_m.
\]
Finite multilinearity is legitimate coefficientwise: a coefficient of $a^r$
can involve only features of weight at most $r$.  Thus
$a^{N_m}\mid\det\mathcal H(a,X)$ in $\mathbb Q[[a]][X]$.  The determinant-one
row shear gives
$\det\mathcal H(a,X)=G_{m,a}(X-L_2(a))$; since $L_2\in a\mathbb Q[[a]]$,
translation preserves the order at zero.  Therefore $a^{N_m}$ divides
$P_m$ in $\mathbb Q[a,T]$.  As $P_m\in\mathbb Z[a,T]$ and the divisor is
the monomial $a^{N_m}$, the quotient is also in $\mathbb Z[a,T]$.
\end{proof}

\begin{lemma}[endpoint pole divisor]\label{lem:p0-pole}
Coefficientwise in $T$,
\[
 \operatorname{ord}_{a=1}G_{4h,a}(T)\ge3h-h^2.
\]
Consequently
\begin{equation}\label{eq:p0-refined-divisor}
 d_U^{3m}(a-1)^{H_m}G_{m,a}(T)
       \in a^{N_m}\mathbb Z[a,T].
\end{equation}
The exponent $3h-h^2$ is a sufficient lower bound; equality for every $h$ is
not asserted.
\end{lemma}

\begin{proof}
Put $b=a-1$ and $f(x)=\binom{x+g}{g}$.  At $a=1$, the full top block,
including its diagonal, is the divided-difference polynomial
\[
 C(x,y)=\frac{f(x)-f(y)}{y-x},\qquad C(x,x)=-f'(x),
\]
of total degree at most $g-1$.  For the lower block, the source Mellin
identity gives
\[
 (D_gM_{ij})(a)=(-1)^{g+1}\sum_{r\ge1}
   \binom{r-1}{g}\frac{a^{-r}}{(r+i)(r+j)}
 =B_{ij}+L(a)C_{ij}-\Li_2(1/a)A_{ij}.
\]
Polynomial division in $r$ yields
\[
 \frac{\binom{r-1}{g}}{(r+i)(r+j)}
 =\sum_{d=0}^{g-2}q_d(i,j)r^d+O(r^{-1}),
 \qquad \deg_{i,j}q_d\le g-2-d.
\]
The power sum $\sum_{r\ge1}r^d(1+b)^{-r}$ has pole order at most $d+1$;
the remainder is $O(1+|\log b|)$.  Since $B$ is rational in $b$, while
$LC-\Li_2(1/a)A=O(1+|\log b|)$, this identifies its entire principal
part, including diagonal entries:
\[
 [b^{-s}]B_{ij}(1+b)=P_s(i,j),\qquad
 \deg P_s\le g-1-s\quad(1\le s\le g-1).
\]
Apply forward differences of orders $u=0,\ldots,m-1$ to the rows in each
block and $v=0,\ldots,2m-1$ to columns.  These are unitriangular transforms,
so they preserve the determinant.  The degree bounds imply
\[
 \operatorname{ord}_b\widehat C_{uv}\ge[ u+v\ge g],\qquad
 \operatorname{ord}_b(\widehat B-T\widehat A)_{uv}
       \ge\min(0,u+v-g+1).
\]
In a determinant term, let $k$ be the number of lower entries with a negative
bound.  Then $0\le k\le g-1$, and their distinct row and column indices
give a total order at least $k(k-g)$.  Their columns lie among
$0,\ldots,g-2$; hence at most $g-k$ top entries can have order-zero bounds,
so at least $m-g+k$ top entries contribute order at least one.  Every
determinant term has order at least
\[
 m-g+k+k(k-g)=m-g+k(k-g+1).
\]
For $m=4h$, $g=2h$, this is
\[
 2h+(k-h)(k-h+1)-h(h-1)\ge3h-h^2,
\]
because $(k-h)(k-h+1)\ge0$ for integral $k$.  This proves the order bound,
including all diagonal entries.

By Lemma~\ref{lem:p0-a-divisor},
$P_m=d_U^{3m}b^{8h^2}G=a^{N_m}S$ with $S\in\mathbb Z[a,T]$.  The pole
bound shows that $b^{7h^2+3h}$ divides $P_m$ in $\mathbb Q[a,T]$.
This divisor is monic, so division preserves integrality; it is coprime to
$a^{N_m}$.  Cancelling it yields
\[
 d_U^{3m}b^{8h^2-(7h^2+3h)}G
 =d_U^{3m}b^{H_m}G\in a^{N_m}\mathbb Z[a,T],
\]
as claimed.  The argument also covers $h=1,2$, when $H_m=-2$; in those
cases the displayed expression is an identity of rational functions whose
right side is the stated integer polynomial.
\end{proof}

\subsection{Primewise valuation bounds}

For a fixed integer $a$ with $|a|\ge2$, define
\[
 F_{m,p}(a)=3m d_p+H_m v_p(a-1)-N_m v_p(a).
\]
We next define a second, moving-prime bound.  For an odd prime satisfying
$p^2>U+g$, put
\[
 r=g\bmod p,\qquad
 \kappa_n=[n\bmod p\ge p-r],\qquad
 \theta_n=(\kappa_n-1)[n\ge p\text{ or }p\le g],
\]
where $\kappa_n=0$ if $r=0$.  Define the integer $W_m(p)$ by the following
deterministic rule (it is the chosen upper bound, not an undefined optimum):
\begin{itemize}
\item if $p^2\le U+g$, set $W_m(p)=3m d_p$;
\item if $p^2>U+g$ and $p$ lies on a positive-ratio wall, set $W_m(p)=3m$;
\item otherwise, for $p/m<1/2$, set $W_m(p)$ to the smaller of the two
  integer right sides in \eqref{eq:p0-low-W}; for
  $1/2<p/m<5/2$, generate the groups from the affine endpoints below,
  sum the displayed row and column potentials in Table~\ref{tab:p0-high-potentials}
  over those groups, and take that exact integer objective;
\item if $p>U+g$, set $W_m(p)=0$.
\end{itemize}
At the low-ratio endpoint $p/m=1/2$, the wall convention $3m$ applies.
The group widths are integers because $4\mid m$ and all endpoint intercepts
have denominator at most four; consequently the high-ratio objective is an
integer, not a real relaxation.  These clauses define $W_m(p)$ for every
prime before it is used in \ref{lem:p0-entry-bounds}.
The finite collection of walls in $[1/128,5/2]$ is defined below.  These
rules are compatible at overlapping cases: the generic denominator bound is
used first, and $p>U+g$ is outside every positive-ratio cell.

\begin{lemma}[coefficientwise entry bounds]\label{lem:p0-entry-bounds}
For every integer specialization $a$ and every coefficient in $T$, the
valuation of the determinant satisfies
\[
 -v_p(G_{m,a})\le W_m(p)+gm\,v_p(a-1).
\]
Also, for every prime $p$,
\[
 -v_p(G_{m,a})\le F_{m,p}(a).
\]
\end{lemma}

\begin{proof}
The second inequality follows directly from
\eqref{eq:p0-refined-divisor}.  For the first, multiply every lower row by
$b^g$, so that
$\det[C;b^g(B-TA)]=b^{gm}G_{m,a}(T)$.  When $p^2\le U+g$, the factors
$d_U$ and $d_U^2$ clear top and lower rows, giving the claimed $3md_p$
bound.  When $p>U+g$, every source denominator is a $p$-unit and the scaled
entries are $p$-integral.

It remains to consider odd $p$ with $p^2>U+g$.  Every relevant upper index
is less than $p^2$.  Kummer's carry rule therefore gives
\[
 v_p\binom{n+g}{g}=\kappa_n.
\]
If $p\mid s$, then $g-s$ has the same units digit as $g$, so
\[
 v_p\binom{n+g}{g-s}\ge\kappa_n;
\]
 and, if $p\mid s$ and $s\le n$,
$v_p\binom{n-s+g}{g}\ge\kappa_n$.  No second carry is possible because
all sums are below $p^2$.  After multiplying by $b^g$, each $Q_s$ is the
primitive polynomial $R_s$ above, whose Gauss valuation is zero.  In these
sums, terms with $p\nmid s$ are $p$-integral; in the mixed
$B_{ii}$ terms, a single $p$-denominator contributes at worst
$-1\ge\kappa_i-2$, since $\kappa_i\le1$.  Every nonzero source denominator
and every row or column difference is less than
$p^2$, hence has $p$-valuation at most one.  The ultrametric inequality now
gives the following lower bounds for scaled entries:
\begin{center}
\renewcommand{\arraystretch}{1.12}
\normalsize
\begin{tabular}{@{}p{0.43\linewidth}p{0.46\linewidth}@{}}
\toprule
Entry & Lower bound for $v_p$ \\
\midrule
$C_{ij}$, $i\ne j$ & $\min(\kappa_i,\kappa_j)-[i\equiv j\pmod p]$ \\
$C_{ii}$ & $(\kappa_i-1)[p\le g]$ \\
$(b^gB)_{ij}$, $i\ne j$ & $\min(\theta_i,\theta_j)-[i\equiv j\pmod p]$ \\
$(b^gB)_{ii}$, $i\ge p$ or $2p\le g$ & $\kappa_i-2$ \\
$(b^gB)_{ii}$, $i<p$ and $p\le g<2p$ & $-1$ \\
$(b^gB)_{ii}$, $i<p$ and $p>g$ & $0$ \\
\bottomrule
\end{tabular}
\end{center}
For $C_{ii}$, the case $p\le g$ uses the general bound
$v_p(U_i/(i+s))\ge\kappa_i-1$, since $i+s<p^2$.  If $p>g$ and
$p\mid i+s$, then $1\le s\le g<p$, so $i\bmod p=p-s$; adding $g$ to
$i$ produces exactly one base-$p$ carry and hence $\kappa_i=1$.  Thus every
term with a $p$-divisible denominator is still $p$-integral, and all other
terms are $p$-integral because $v_p(U_i)=\kappa_i\ge0$.  For the lower diagonal, a squared denominator
in the first sum requires $i\ge p$; two $p$-denominators in the correction
sum require $2p\le g$.  In either case the relevant binomial coefficient
has valuation at least $\kappa_i$.  If $p\le g<2p$ and $i<p$, at most one
denominator is divisible by $p$; if $p>g$ and $i<p$, none is.  The term
$-TA$ is integral after the row scaling and does not weaken these bounds.

By the Leibniz expansion, minus the determinant valuation is at most the
maximum assignment weight formed by negating these entry lower bounds.
The exact profile $W_m(p)$ bounds that assignment in each range.  The
certificate details follow.
\end{proof}

\subsection{Exact moving-prime profile}

Let $\rho=p/m$.  For a residue $v\in\{0,\ldots,p-1\}$, the copy counts are
\[
 c_I(v)=\#\{n\in\mathbb Z:1/4\le v/m+n\rho<5/4\},\qquad
 c_J(v)=\#\{n\in\mathbb Z:0\le v/m+n\rho<2\}.
\]
The carry flag is constant on each residue group.  For $p<g$, let $K_I,K_J$
count the killed row and column copies ($\kappa_v=1$), and put
\[
 S_U=\sum_{\kappa_v=0}\min(2c_I(v),c_J(v)),\qquad
 S_K=\sum_{\kappa_v=1}\min(2c_I(v),c_J(v)).
\]
The two feasible assignment bounds are
\begin{equation}\label{eq:p0-low-W}
 W_m(p)\le\min\{m+S_U,\;3m-2K_I-K_J+S_K\}.
\end{equation}
Indeed, after subtracting the lower-row baseline $m$, the assignment score is
$x-z$, where $x$ counts unkilled same-residue matches and $z$ killed
off-residue matches.  If $y$ is the number of killed same-residue matches,
then $x\le S_U$, $y\le S_K$, and counting killed row and column copies gives
$z+y\ge2K_I+K_J-2m+x$.  These yield the two inequalities.

All group endpoints are generated by a finite exact rule.  For a fixed open
wall cell and any interior rational sample $\rho_0$, set
$q=\lfloor1/(2\rho_0)\rfloor$ and order the affine endpoints
\[
 (0,0),\ (0,1),\ (-1/2,q+1),\quad
 (x,-\lfloor x/\rho_0\rfloor)\ (x=1/4,5/4,2),
\]
by their values $A+B\rho_0$.  Consecutive endpoints give a group of width
$A+B\rho$.  At its midpoint $t$, the $I$- and $J$-copy lists are respectively
\[
 \left\lceil\frac{1/4-t}{\rho_0}\right\rceil\le n<
 \left\lceil\frac{5/4-t}{\rho_0}\right\rceil,
 \qquad
 \left\lceil\frac{-t}{\rho_0}\right\rceil\le n<
 \left\lceil\frac{2-t}{\rho_0}\right\rceil,
\]
and the group is killed precisely when
$t\ge(q+1)\rho_0-1/2$.  These data remain constant throughout the open cell.
The full wall set in a compact ratio interval is generated by
\[
\mathcal W([r_0,r_1])=\{r_0,r_1\}\cup
\left\{\frac{x-y}{n}:\begin{aligned}[t]
&x,y\in\{0,1/4,5/4,2,-1/2\},\quad x>y,\\
&n\in\mathbb N,\quad r_0<\frac{x-y}{n}<r_1
\end{aligned}\right\}.
\]
It is finite because $x-y\le5/2$ and $n<(x-y)/r_0$.

For $\rho\in[\delta,1/2]$, $\delta=1/128$, the rule gives 756 lower-envelope
affine pieces.  Their exact integral is
\begin{equation}\label{eq:p0-low-integral}
 L_*=
 \frac{\begin{array}{c}
 \texttt{1448276314740194964855260100800462127260298}\\
 \texttt{03104279674952972562737001039673489563365828533335}\\
 \texttt{0450905704443029571002838710237}
 \end{array}}
 {\begin{array}{c}
 \texttt{140409888158080315725937691852162573583946167555}\\
 \texttt{38734362972176526622161627829570435939660600468548}\\
 \texttt{26150978176673015236224000}
 \end{array}}.
\end{equation}
In each array, concatenate its three rows to obtain the integer.  This value
is obtained by integrating the minimum in \eqref{eq:p0-low-W}
on each rational group cell, splitting at its rational affine crossings.
The formula above is the exact reconstruction rule; no optimizer is used.

For $0<\rho<\delta$, zero column potentials give
$W_m(p)\le3m-2K_I$.  Since an interval of $m$ consecutive integers contains
at least $\lfloor m/p\rfloor(g\bmod p)$ killed indices, the profile
$f_{\rm tail}(\rho)=3-2\{1/(2\rho)\}+2\rho$ satisfies
$W_m(p)/m\le f_{\rm tail}(\rho)$.
Writing $\delta=1/(2Q)$ with $Q=64$ and substituting
$x=1/(2\rho)$ gives
\[
 \int_0^\delta\left(3-2\left\{\frac1{2\rho}\right\}\right)d\rho
 =\frac1{2Q}+\sum_{q\ge Q}\left(\frac1q-\log(1+1/q)\right).
\]
For each $q\ge Q$, the summand differs from $1/(2q(q+1))$ by
\[
 \int_0^1\frac{s(1-s)}{q(q+s)(q+1)}\,ds
 \le\frac1{6q^2(q+1)}.
\]
Summing this bound and integrating the added $2\rho$ term yields
\begin{equation}\label{eq:p0-small-tail}
 \int_0^\delta f_{\rm tail}(\rho)\,d\rho\le
 2\delta+2\delta^2=\frac{129}{8192}.
\end{equation}

For $1/2<\rho<5/2$, $p>g$, so each group has a carry bit $k\in\{0,1\}$.
If $n$ denotes a copy index, put
$\theta=(k-1)[n\ge1]$.  A row type is $(j,n,e,k)$, where $j$ is its group,
$n$ is an $I$-copy and $e\in\{0,1\}$ is the top or lower block; a column
type is $(j,n,k)$ with $n$ a $J$-copy.  Order row types by increasing group,
then copy, then $e=0,1$; order column types by increasing group then copy.
For a row type $r$ and column type $c$, the required potential inequality is
\[
 u_r+v_c\ge w_{rc}.
\]
For different groups the right side is $-\min(k_r,k_c)$ in the top block
and $-\min(\theta_r,\theta_c)$ in the lower block.  For the same group but
distinct copies add $1$.  For a same-index top diagonal use $0$; for a
same-index lower diagonal use $(2-k_r)$ if $n\ge1$, and $0$ if $n=0$.
These are exactly the negated entry bounds of
Lemma~\ref{lem:p0-entry-bounds}, including both diagonal cases.  The
following integer vectors give feasible potentials on every large-ratio
interval; spaces separate successive coordinates and the semicolon separates
the row and column vectors.  The group and copy ordering is the rule just
given.  Each group contains exactly $m(A_j+B_j\rho)$ residues; this is an
integer because $m$ is divisible by four and all endpoint intercepts have
denominator at most four.  Thus the potential objective is an integer upper
bound at every prime ratio, not merely a real-valued relaxation.

\renewcommand{\arraystretch}{1.15}
\normalsize
\begin{longtable}{@{}c p{0.33\linewidth}p{0.34\linewidth}c@{}}
\caption{Exact large-ratio assignment potentials.  The last column is
  $\sum_r u_r(A_{j(r)}+B_{j(r)}\rho)+
    \sum_c v_c(A_{j(c)}+B_{j(c)}\rho)$.}\label{tab:p0-high-potentials}\\
\toprule
Ratio cell & Row potentials $u$ & Column potentials $v$ & Objective \\
\midrule
\endfirsthead
\toprule
Ratio cell & Row potentials $u$ & Column potentials $v$ & Objective \\
\midrule
\endhead
$(1/2,7/12)$ &
\texttt{1 2 1 2 0 1 0 1 1 2 0 1 0 1 0 1 0 1} &
\texttt{0 0 0 0 0 0 0 0 -1 -1 -1 -1 0 0 0 0 0 0 0} & $1$ \\
$(7/12,5/8)$ &
\texttt{1 2 1 2 1 2 1 2 1 2 0 1 0 1} &
\texttt{0 0 0 0 0 0 0 0 -1 -1 -1 -1 -1 -1 -1 0 0 0} & $3\rho-3/4$ \\
$(5/8,2/3)$ &
\texttt{1 2 1 2 1 2 0 1 0 1 1 2} &
\texttt{0 0 0 0 0 0 0 -1 -1 -1 0 0 0 -1 -1 -1} & $\rho+1/2$ \\
$(2/3,3/4)$ &
\texttt{1 2 1 2 0 1 0 1 1 2 0 1} &
\texttt{0 0 0 -1 -1 -1 0 0 0 -1 -1 -1 0 0} & $4\rho-3/2$ \\
$(3/4,5/6)$ &
\texttt{1 2 0 1 0 1 0 1 0 1 0 1 0 1 0 1} &
\texttt{0 0 0 1 1 1 0 0 0 0 0 0 0} & $3\rho-3/4$ \\
$(5/6,7/8)$ &
\texttt{1 2 0 1 0 1 0 1 0 1 0 1 0 1 0 1} &
\texttt{0 0 0 1 1 1 1 1 0 0 0 0} & $7/4$ \\
$(7/8,1)$ &
\texttt{1 2 0 1 0 1 0 1 0 1 0 1} &
\texttt{0 0 0 1 1 1 1 0 1 0 0} & $7/4$ \\
$(1,5/4)$ &
\texttt{1 2 1 1 1 1 1 1} &
\texttt{0 0 -1 0 -1 1 -1 0 -1} & $15/4-2\rho$ \\
$(5/4,7/4)$ &
\texttt{0 0 0 0 0 0} &
\texttt{0 1 0 2 0 0 0} & $15/4-2\rho$ \\
$(7/4,2)$ &
\texttt{0 0} &
\texttt{0 1 0 0 0 0} & $2-\rho$ \\
$(2,5/2)$ & $\texttt{0 0}$ & $\texttt{0 0 0 0}$ & $0$ \\
\bottomrule
\end{longtable}

The eleven large-ratio intervals are exactly the walls generated by
$\mathcal W([1/2,5/2])$.  For each interval, the group formula above
determines the row and column types, and the displayed vectors satisfy all
the pairwise rational inequalities.  The edge, diagonal, and wall comparisons
are generated by the same displayed rule.  Summing the
type potentials times their affine group widths
gives the objective column; taking its exact integral over the high cells
yields
\[
 \int_{1/2}^{5/2}f_{\rm high}(\rho)\,d\rho=\frac{287}{192}.
\]
At a wall use $W_m(p)=3m$; there are finitely many walls, so this changes
the prime sum by $O(m\log m)$ for each $m$.

The exact profile constant is
\begin{align}
 K&=L_*+\frac{129}{8192}+\frac{287}{192}\\
 &=\frac{\begin{array}{c}
 \texttt{28553776527336974955476146348812001767941}\\
 \texttt{86855224421109769015232357578975297803041084792395}\\
 \texttt{8183617568292693602659280722163521}
 \end{array}}
 {\begin{array}{c}
 \texttt{112327910526464252580750153481730058867156934044}\\
 \texttt{30987490377741221297729302263656348751728480374838}\\
 \texttt{609207825413384121889792000}
 \end{array}}\\
 &=2.542001929307655\ldots<2.542002.
\label{eq:p0-K}
\end{align}
The displayed rational table and the type-generation rule are the complete
finite input for this part of the estimate.

\subsection{Exceptional primes and the PNT limit}

Define the rational scalar
\begin{equation}\label{eq:p0-delta}
 \Delta_{m,a}=\prod_p p^{E_{m,p}(a)},\qquad
 E_{m,p}(a)=\min\{W_m(p)+gm\,v_p(a-1),\ F_{m,p}(a)\}.
\end{equation}
Negative exponents are allowed; $\Delta_{m,a}>0$ is a rational number.
The product is finite: if $p>U+g$ and $p\nmid a(a-1)$, both entries in
the minimum are zero.  Lemma~\ref{lem:p0-entry-bounds} implies
coefficientwise that
\begin{equation}\label{eq:p0-integerization}
 \Delta_{m,a}G_{m,a}(T)\in\mathbb Z[T].
\end{equation}
This is a specialization-wise p-adic argument; it does not claim that the
sharp scalar $\Delta_{m,a}$ is the specialization of an integer polynomial
in $a$.

Since $\gcd(a,a-1)=1$, the exceptional sets $p\mid a$ and $p\mid a-1$
are disjoint.  For a fixed $p\mid a$,
\[
 F_{m,p}=O(m\log m)-N_m v_p(a),
 \qquad W_m(p)+gmv_p(a-1)=W_m(p),
\]
and the quadratic negative term makes $F_{m,p}$ the minimum for all
sufficiently large $m$.  For a fixed $p\mid a-1$,
\[
 F_{m,p}=O(m\log m)+H_m v_p(a-1),\qquad
 W_m(p)+gmv_p(a-1),
\]
and $gm-H_m=7h^2+3h$ makes $F_{m,p}$ eventually the minimum.  These
comparisons apply also to $p=2$, which divides exactly one of $a,a-1$.
For every other fixed prime the moving profile applies, with the
$p^2\le U+g$ range contributing $o(m^2)$ in aggregate:
Put $X=U+g$.  The squarefull part is bounded without counting a
prime more than once by
\[
 \sum_{p\le\sqrt X}3m d_p\log p
 \le 3m\sum_{j=1}^{\lfloor\log_2 U\rfloor}
    \vartheta\!\left(\min\{\sqrt X,U^{1/j}\}\right)
 =O(m^{3/2}\log m)=o(m^2),
\]
where $\vartheta(x)=\sum_{p\le x}\log p$.  The displayed equality is an
upper-bound comparison (the left side is bounded by the right side),
since $d_p$ is the number of admissible exponent layers and each layer is
charged to $p^j\le U$; the extra $\log m$ is harmless after normalization.

On any fixed positive ratio interval, $W_m(p)\le m f(p/m)$ with $f$ bounded
and piecewise affine, apart from the finite wall values.  The weighted prime
number theorem, obtained from the ordinary PNT by partial summation, gives
\[
 \frac1m\sum_p f(p/m)\log p\longrightarrow\int f(\rho)\,d\rho.
\]
For $0<\rho<\eps$, first remove the already negligible range
$p\le\sqrt{U+g}$ and then omit $\sqrt{U+g}<p\le\eps m$; there $d_p=1$ and
the crude bound $3m\vartheta(\eps m)$ has normalized limsup at most
$3\eps$.  Apply the weighted PNT on $[\eps,\delta]$ and then let
$\eps\downarrow0$.  The exact
tail bound \eqref{eq:p0-small-tail}, the low integral \eqref{eq:p0-low-integral},
and the high integral therefore imply
\[
 \limsup_{\substack{m\to\infty\\4\mid m}}
 \frac1{m^2}\sum_{p\nmid a(a-1)}E_{m,p}(a)\log p\le K.
\]
The fixed exceptional primes contribute
\[
 \frac{H_m}{m^2}\log|a-1|
 -\frac{N_m}{m^2}\log|a|
 \longrightarrow
 \frac1{16}\log|a-1|-\frac{13}{16}\log|a|.
\]
We have proved the following fixed-base estimate.

\begin{proposition}[P0 fixed-base denominator rate]\label{prop:p0-sharp-arith}
For every fixed integer $|a|\ge2$, the rational number $\Delta_{m,a}$ in
\eqref{eq:p0-delta} satisfies \eqref{eq:p0-integerization} and
\[
 \limsup_{\substack{m\to\infty\\4\mid m}}
 \frac{\log\Delta_{m,a}}{m^2}
 \le K+\frac1{16}\log|a-1|-
             \frac{13}{16}\log|a|,
 \qquad K<2.542002.
\]
The statement is for each fixed base.  In particular, its finite-$m$
remainder is not uniform as $|a|\to\infty$.
\end{proposition}

\paragraph{Completion of the estimate.}
The formal divisor and pole lemmas, Kummer bounds, residue-cell formulas,
exceptional-prime comparisons, and PNT passage are all-$m$ derivations.
The low profile is reconstructed by the rational wall-and-crossing rule
above.  The high profile is hand-checkable from the printed potential vectors
and the stated type-generation rule: each vector pair reduces the full
assignment inequality to finitely many integer comparisons.  The low profile
has 756 affine pieces and the high profile eleven cells.  The all-$m$
conclusion uses the displayed identities, the finite comparison count above,
and the weighted prime number theorem.


\section{Positive logarithmic energy on an interval}\label{sec:p0-analytic}

This section records the analytic passage for the positive P0 determinant,
including the dependence on the base.  The finite determinant statement is
for integers $q$ with $2\le q\le11$; a real parameter is used only in the
continuum interpolation.  All measures, intervals, and the contour below are
fixed as $q$ varies.  The two endpoint reference measures are the exact data
in Tables~\ref{tab:q2data} and~\ref{tab:q11data}.

Put $m\in4\mathbb N$, $m\ge8$, $g=m/2$,
\[
 I_m=\{m/4,\ldots,5m/4-1\},\qquad
 J_m=\{0,\ldots,2m-1\},\qquad E=[1/4,5/4],
\]
and fix
\[
 \Gamma=\{\xi+A\cos u+iB\sin u:0\le u<2\pi\},\quad
 (\xi,A,B)=\left(\frac{11281}{10000},\frac{12783}{10000},
                                      \frac{229}{200}\right).
\]
Write $\Delta(X)=\prod_{a<b}(x_b-x_a)$ and $K(x,y)=(x+y)^{-1}$.
The source function and the frequency weight at base $q>1$ are
\[
 F_{g,q}(z)=q^z\binom{z+g}{g},\qquad
 U_i(q)=q^i\binom{i+g}{g},\qquad
 w_{r,q}=\binom{r-1}{g}q^{-r}\quad(r>g).
\]
For integer $q$, let $G_{m,q}$ be the P0 determinant of
Section~\ref{sec:p0}, with $a=q$.  The diagonal contour identity is
\begin{equation}\label{eq:p0-contour-q}
 C^{(g)}_{ij}(q)=-\oint_{m\Gamma}
       \frac{F_{g,q}(z)}{(z-i)(z-j)}\frac{dz}{2\pi i}
             +(\log q)U_i(q)\,\mathbf1_{i=j}.
\end{equation}
For $i\ne j$ this is the divided-difference residue identity; for $i=j$ it
follows by differentiating $F_{g,q}$, since
\[
 F_{g,q}'(i)=U_i(q)\left(\log q+\sum_{s=1}^g\frac1{i+s}\right).
\]
Expanding the rational moment expression in the frequency basis and applying
Cauchy--Binet~\cite{Gantmacher} gives an absolutely convergent sum over
$R\subset\{g+1,g+2,\ldots\}$, $|R|=m$, with weight
$\prod_{r\in R}w_{r,q}$.  After choosing $k$ diagonal rows $S\subset I_m$
and putting $d=m-k$, its absolute integrand is
\begin{equation}\label{eq:p0-q-integrand}
 \begin{aligned}
 Q_{m,q}(S,R,Z)={}&\left(\prod_{i\in S}(\log q)U_i(q)\right)
   \left(\prod_{r\in R}w_{r,q}\right)
   \left(\prod_{z\in Z}|F_{g,q}(z)|\right)\\
 &\times\left|\det K_{I_m\setminus S,-Z}\det K_{I_m,R}
                \det K_{(-Z,R),J_m\setminus S}\right|.
 \end{aligned}
\end{equation}
where $Z=(z_1,\ldots,z_d)\in(m\Gamma)^d$.  The integration measure is
$\prod |dz|/(2\pi)$ and the symmetrization factor is $1/d!$.  The three
Cauchy factors have, after restoring the deleted fixed-grid factors, the
same three Vandermonde interactions as at $q=2$: contour-contour, pinned-
pinned, and frequency-frequency each have coefficient one; pinned-contour
has coefficient two; pinned-frequency and contour-frequency each have
 coefficient one.  In particular, no factor depending on $q$ enters the
 quadratic interaction energy.

\begin{proof}[Cauchy--Binet energy expansion~\cite{Gantmacher}]
The rational moment expression is obtained by expanding the two Cauchy kernels
geometrically and using the exact generating identity
\[
 \sum_{r\ge g+1}\binom{r-1}{g}(t/q)^r
       =\left(\frac{t}{q-t}\right)^{g+1},\qquad 0\le t<q.
\]
Inserting these moments and applying
\[
 \det(XY)=\sum_{|R|=m}\det X_{:,R}\det Y_{R,:}
\]
gives the sum over the full frequency sets $R$ stated above.  Splitting the
diagonal term in \eqref{eq:p0-contour-q} row by row chooses
$S\subset I_m$; the remaining $d=m-|S|$ rows are represented by contour
variables $Z$, with the factor $1/d!$ from their antisymmetrization.
For any two equal-size sets, the Cauchy determinant identity~\cite{Krattenthaler}
\[
 \det\left(\frac1{x_i+y_j}\right)_{i,j=1}^r
 =\frac{\Delta(x)\Delta(y)}{\prod_{i,j}(x_i+y_j)}
\]
converts each of the three kernel minors into a Vandermonde quotient.
Restoring the fixed-grid Vandermonde factors leaves contour-contour,
pinned-pinned, and frequency-frequency interactions with coefficient one,
pinned-contour with coefficient two, and the other two mixed interactions
with coefficient one.  Taking absolute values gives exactly
\eqref{eq:p0-q-integrand}; all signs and $q$-powers are one-body terms.
For fixed $m$, the Cauchy minors have at most polynomial growth in the
frequencies and $w_{r,q}=O_q(r^gq^{-r})$.  The absolute frequency sum
therefore converges, and increasing finite truncations converge to its
full absolute envelope.
\end{proof}

For the normalization, the signed Cauchy--Binet expansion~\cite{Gantmacher}
can be written
explicitly. Put $n=m/4$, $Y=S\uplus Z$, and
\[
 \begin{aligned}
 D_m(a)&=(a-n)!(n+m-1-a)!a!(2m-1-a)!,\\
 P_m(z)&=\prod_{a\in I_m\uplus J_m}(z-a),\qquad
 Q_m^*(r)=\prod_{a\in I_m\uplus J_m}(r+a).
 \end{aligned}
\]
Then
\begin{align*}
G_{m,q}={}&\Delta(I_m)^2\Delta(J_m)
 \sum_{S\subset I_m}\frac{(-1)^{|S|(n+1)}(\log q)^{|S|}}{d!}
 \sum_{|R|=m}
 \left(\prod_{a\in S}\frac{U_a(q)}{D_m(a)}\right)
 \left(\prod_{r\in R}w_{r,q}\right)\\
&\quad\times\oint_{(m\Gamma)^d}
 \frac{\Delta(Y)^2\Delta(R)^2\prod_{y\in Y,r\in R}(y+r)}
      {\prod_{z\in Z}P_m(z)\prod_{r\in R}Q_m^*(r)}
      \prod_{z\in Z}\left(F_{g,q}(z)\frac{dz}{2\pi i}\right).
\end{align*}
This follows by applying Cauchy--Binet~\cite{Gantmacher} to the two Cauchy
moment factors,
then expanding the diagonal summands in the contour formula row by row.
For every equal-size pair, use the Cauchy determinant identity~\cite{Krattenthaler}
\[
 \det\left((x_i+y_j)^{-1}\right)_{i,j=1}^r
 =\frac{\Delta(x)\Delta(y)}{\prod_{i,j}(x_i+y_j)}.
\]
Restoring the deleted fixed-grid factors uses
\[
 \Delta(I_m\setminus S)\Delta(J_m\setminus S)
 =\frac{\Delta(I_m)\Delta(J_m)\Delta(S)^2}
        {\prod_{a\in S}D_m(a)}.
\]
Each pinned index is therefore counted twice in a contour denominator and
once in a frequency denominator. The three Cauchy factors give
contour-contour, pinned-pinned, and frequency-frequency interactions with
coefficient one, pinned-contour with coefficient two, and each remaining
mixed interaction with coefficient one. Taking absolute values and scaling
all positions by $m$ gives \eqref{eq:p0-q-integrand} and
\eqref{eq:p0-q-energy}. The factors $q^a$, $q^{-r}$, and $q^z$ contribute
only the one-body fields $\lambda\Re z$ and $-\lambda t$ in
\eqref{eq:p0-q-fields}; the quadratic interaction coefficients are
independent of $q$.

The fixed-grid normalization is also explicit.  Consecutive integer lists
satisfy $\Delta(I_m)=\prod_{j=1}^{m-1}j!$ and
$\Delta(J_m)=\prod_{j=1}^{2m-1}j!$.  Stirling's formula~\cite{WhittakerWatson}
summed over these
products gives
\[
 \log(\Delta(I_m)^2\Delta(J_m))
 =3m^2\log m+(2\log2-9/2)m^2+O(m\log m).
\]
Equivalently, the fixed Vandermonde contributes $C_0=9/2-2\log2$ to
the negative logarithmic energy.  The scaled Cauchy quotients cancel all
$m^2\log m$ terms; their remaining homogeneous factor is
$m^{-2m-2d}$ before the contour differentials and $m^{-2m-d}$ after them.
This accounts for the full quadratic normalization in
\eqref{eq:p0-q-energy}.

\subsection{The fixed variational problem and its affine parameter}

Let $\eta=\kappa+\nu$, where
\[
 0\le\kappa\le1_E\,dx,\quad \nu\ge0,\quad
 \operatorname{supp}\nu\subset\Gamma,\quad |\kappa|+|\nu|=1,
\qquad 0\le\mu\le1_{[1/2,20]}dt,\quad |\mu|=1.
\]
Write $\mu^-=(-\operatorname{id})_\#\mu$ for reflection onto the negative
real axis.
The set $\mathcal A$ of such pairs $(\eta,\mu)$ is independent of $q$.
Define
\[
 A_0(z)=\int_E\log|z-x|\,dx+\int_0^2\log|z-x|\,dx,
 \qquad
 B_0(t)=\int_E\log(t+x)\,dx+\int_0^2\log(t+x)\,dx,
\]
\begin{align*}
 \phi_2(z)&=(\log2)\Re z+\int_0^{1/2}\log|z+s|\,ds
                              +\frac12+\frac12\log2,\\
 \psi_2(t)&=t\log t-(t-1/2)\log(t-1/2)
                         +\frac12\log2-t\log2.
\end{align*}
For measures on the corresponding reflected supports write
\[
 I(\sigma,\tau)=-\iint\log|x-y|\,d\sigma(x)d\tau(y).
\]
The energy is
\begin{align}
 \mathcal E_{q}(\eta,\mu)={}&C_0+
     \int(A_0-\phi_q)\,d\eta+\int(B_0-\psi_q)\,d\mu \\
 &+I(\eta)+I(\mu)+I(\eta,\mu^-),
 \label{eq:p0-q-energy}
\end{align}
Here $I$ is positive on zero-mass variations, while $Q_{\log}$ below uses
the opposite, positive-logarithm convention.  Also
$C_0=9/2-2\log2$ and, on writing $q=2e^\lambda$,
\begin{equation}\label{eq:p0-q-fields}
 \phi_q(z)=\phi_2(z)+\lambda\Re z,
 \qquad \psi_q(t)=\psi_2(t)-\lambda t.
\end{equation}
The logarithmic integrals are understood in the following regularized sense:
truncate at the logarithmic singularity, integrate the resulting bounded
functions, and take the limit after replacing the coincident-cell term by
the continuous primitive \(u(\log|u|-1)\), with value \(0\) at \(u=0\).
The measures used below have bounded densities and compact support, so this
limit is finite.  Substitution into
\eqref{eq:p0-q-energy} shows, for every fixed admissible pair,
\begin{equation}\label{eq:p0-q-affine-energy}
 \mathcal E_\lambda(\eta,\mu)=\mathcal E_0(\eta,\mu)
       +\lambda\left(-\int\Re z\,d\eta(z)+\int t\,d\mu(t)\right).
\end{equation}
Consequently
\begin{equation}\label{eq:p0-q-alpha}
 \alpha(\lambda)=\inf_{(\eta,\mu)\in\mathcal A}
                    \mathcal E_\lambda(\eta,\mu)
\end{equation}
is concave.  Indeed, for $0\le\theta\le1$, each fixed-pair energy at
$\theta\lambda_0+(1-\theta)\lambda_1$ is the corresponding convex
combination of its endpoint energies, and taking an infimum gives
\[
 \alpha(\theta\lambda_0+(1-\theta)\lambda_1)
       \ge\theta\alpha(\lambda_0)+(1-\theta)\alpha(\lambda_1).
\]
Apply the shared quadratic-duality identity of
Section~\ref{sec:certificate-algorithm} with quadratic part
$Q_{\rm P0}(\eta,\mu)=I(\eta)+I(\mu)+I(\eta,\mu^-)$.  Fix an admissible
reference $(\eta_0,\mu_0)$ and put
\begin{align*}
 Q_{\log}={}&\iint\log|z-w|\,d\eta_0(z)d\eta_0(w)
 +\iint\log|t-s|\,d\mu_0(t)d\mu_0(s)\\
 &+\iint\log|z+t|\,d\eta_0(z)d\mu_0(t).
\end{align*}
\[
Q_{\rm P0}(\eta_0,\mu_0)=-Q_{\log},\qquad
W_\eta=A_0-\phi_q-2L_{\eta_0}-L_{\mu_0^-},\qquad
 W_\mu=B_0-\psi_q-2L_{\mu_0}-L_{\eta_0}(-t).
\]
If $W_\eta\ge h_\eta$ on $\Gamma$, set
\[
 P=\int_E(h_\eta-W_\eta)_+\,dx,\qquad
 T=\int_{1/2}^{20}(h_\mu-W_\mu)_+\,dt.
\]
Then
\begin{equation}\label{eq:p0-q-dual}
 \alpha(\lambda)\ge C_0+Q_{\log}+h_\eta+h_\mu-P-T.
\end{equation}
For $\sigma=\eta-\eta_0$ and $\tau=\mu^- -\mu_0^-$, the shared identity
has remainder
\[
 I(\sigma)+I(\tau)+I(\sigma,\tau)
 =\tfrac12 I(\sigma)+\tfrac12 I(\tau)
                         +\tfrac12 I(\sigma+\tau)\ge0,
\]
where $I(\rho)=-\iint\log|x-y|\,d\rho(x)d\rho(y)$ is nonnegative for
zero-mass signed measures.  The density constraints then give
$\int W_\eta d\eta\ge h_\eta-P$ and
$\int W_\mu d\mu\ge h_\mu-T$, proving \eqref{eq:p0-q-dual}.

\subsection{Finite configurations and cell estimates}

For a selected pinned set $S\subset I_m$ and a frequency set $R$, define
\begin{align*}
 \kappa_S&=\sum_{i\in S}\mathbf1_{[i/m,(i+1)/m)}dx,
 &\mu_R&=\sum_{r\in R}\mathbf1_{[(r-1)/m,r/m)}dt,\\
 \nu_\delta&=\frac1m\sum_{a=1}^{d}\sigma_{z_a/m,\delta},
 &\delta&=m^{-2}.
\end{align*}
Here $\sigma_{z,\delta}$ is probability arclength on the radius-$\delta$
circle centered at $z$.
Put $\zeta_a=z_a/m\in\Gamma$.  Thus all logarithms in the continuum
energy are in scaled coordinates; the powers of $m$ arising from the change
from $z_a$ to $\zeta_a$ have already been included in the fixed
Vandermonde normalization in the cell estimate.
The first two measures have density at most one,
masses $|S|/m$ and one, and the first frequency cell begins at $1/2$.
Circle averaging gives
\begin{equation}\label{eq:p0-circle}
 \sum_{a\ne b}\log|\zeta_a-\zeta_b|
 \le m^2\iint\log|z-w|\,d\nu_\delta(z)d\nu_\delta(w)
                                          +2d\log m.
\end{equation}
The fixed-grid and reflected-frequency logarithmic potentials are harmonic
on the smoothing disks.  The contour field is subharmonic, including across
the zero segment of $F_{g,q}$.  Thus smoothing has no additional field
error; the explicit $2d\log m$ is the full smoothing loss.

The exact real-cell estimates used in the envelope are listed here to make
the finite-size accounting reproducible.  For $n$ distinct lattice sites,
the cell self-energy inequality is
\[
 \sum_{a\ne b}\log|x_a-x_b|
 \le m^2\iint\log|x-y|\,d\rho(x)d\rho(y)
                                  +n(\log m+\log(2\pi)),
\]
where $\rho$ is the union of their length-$1/m$ cells.  To see the constant,
write
\[
 e_d=\log d-\int_0^1\!\int_0^1\log|d+u-v|\,du\,dv.
\]
Then $e_1=3/2-2\log2$ and for $d\ge2$,
\[
 e_d=\sum_{j\ge1}\frac1{j(2j+1)(2j+2)d^{2j}}
       \le\frac1{12d(d-1)},\qquad
 \sum_{d\ge1}e_d=\tfrac12\log(2\pi)-\tfrac34.
\]
Indeed, for $d\ge2$ expand
$\log(1+(u-v)/d)$ in its absolutely convergent power series and integrate
term by term; the even moments
$\int_0^1\!\int_0^1(u-v)^{2j}\,du\,dv=2/((2j+1)(2j+2))$
give the displayed series, while the $d=1$ integral is elementary.  Summing
the same identity from $d=1$ to $M$, exchanging the positive series, and
using Stirling's Wallis limit~\cite{WhittakerWatson}
$\log M!=(M+\tfrac12)\log M-M+\tfrac12\log(2\pi)+o(1)$
gives the displayed infinite sum.  Thus no numerical cell constant is being
imported.
There are at most two sites at each distance from a fixed site.  This proves
the displayed inequality by summing the self-cell and pair-cell errors.
Applied to the selected pinned and frequency sites and to the fixed
Vandermonde, it contributes $4m(\log m+\log(2\pi))$.

The remaining one-body and mixed-cell bounds are as follows.  Binomial
entropy gives
\[
 \log U_i(q)\le m\phi_q(i/m),\qquad
 \log w_{r,q}\le m\psi_q(r/m).
\]
The exact cell corrections from the affine parameter are
$-\lambda m/2$ for frequencies and $-\lambda|S|/2$ for pinned cells.
For the base-$2$ frequency field,
\[
 \psi_2'(t)=\log\frac{t}{t-1/2}-\log2,\qquad
 \int_{1/2}^{1}(\psi_2'(t))_+\,dt=\tfrac12\log2.
\]
Summing endpoint-to-cell differences over disjoint frequency cells costs
at most $(m/2)\log2$.  Monotone quadrature for the fixed-grid field gives,
for $t\ge1/2$,
\[
 0\le mB_0(t)-\sum_{i\in I_m}\log(t+i/m)
                    -\sum_{j\in J_m}\log(t+j/m)
 \le\log(35/3).
\]
Thus the frequency one-body cost is at most
$m(\log(35/3)+\tfrac12\log2)-\lambda m/2$.

For the pinned denominator, take $N$ consecutive sites starting at $a$ and
containing $i$, with $n_L=i-a$ and $n_R=a+N-1-i$.  Comparing the omitted
site product on the two sides of $i$ with its cell integral gives
\[
 \log(n_L!n_R!)-(N-1)\log m
 -m\int_{a/m}^{(a+N)/m}\log|i/m-y|\,dy
 \ge\log\frac{m}{n_R+1}.
\]
This follows from monotone logarithmic sums and
$n_R\log(1+1/n_R)\le1$; the zero case is direct.  The $I_m$ contribution
is nonnegative and the $J_m$ contribution is at least $-\log2$.
Direct integration of the explicit derivatives gives
$\int_E| (\phi_2-A_0)'|\le\log6+2\log2+\log7=\log168$.
Consequently pinned quadrature costs at most $m\log336$.  Its affine
correction is $-\lambda|S|/2$, since the linear-field integral on each
cell is centered one half-cell to the right of the selected lattice point.
For an exact bookkeeping check, $q/2-\log q$ is increasing and positive on
$[2,\infty)$, hence $(\log\log q)_+\le\lambda$.  Therefore the diagonal
and cell-shift contribution satisfies
\[
 |S|(\log\log q)_+-\frac\lambda2(|S|+m)
 \le-\frac\lambda2(m-|S|)\le0.
\]
The displayed pointwise estimate keeps the positive term and drops these
favorable shifts, so it remains valid without relying on this cancellation.

The ellipse clearances are
\[
 d_+=\operatorname{dist}(\Gamma,[0,2])=751/5000,\qquad
 d_-=\operatorname{dist}(\Gamma,(-\infty,-1/2])=1749/5000.
\]
They follow by minimizing the squared distance as a quadratic in
$\cos u$ on the ellipse parametrization.  The pinned-contour displacement
costs $m/d_+$, the fixed contour denominators cost $3m/(2d_+)$, and the
contour-frequency displacement costs $m/(2d_-)$.  The pinned-frequency
cell displacement is centered; Taylor's inequality, with
$i/m+r/m\ge3/4$ and cell width $1/m\le1/4$, bounds its total loss by $1/3$.
Together these contributions are at most
\[
 m\left(\frac5{2d_+}+\frac1{2d_-}\right)+\frac13.
\]
Set $T_\delta=\{z:\operatorname{dist}(z,\Gamma)\le\delta\}$, with the tube
radius $\delta$ used in the transfer argument below.
For the finite contour product the required inequality is one-sided.  Fix
$z\in T_\delta$, put $\zeta=mz$, $r=-\Re\zeta$, $b=\Im\zeta$, and
$f(u)=\log|\zeta+u|$.  Then
\[
 f'(u)=\frac{u-r}{(u-r)^2+b^2},
\]
so $f$ decreases to $r$ and increases after $r$.  Put $H=\max(f,0)$.
The negative part of $f$ is supported near $r$; after translating $r$ and
using $\int_{-1}^{1}(\log|u|)_-\,du=2$, its integral over every unit cell
meeting $r$ is at most $2$.  On each of the two monotone pieces the
  right-endpoint quadrature error for $H$ is nonpositive on the decreasing
  piece and nonnegative on the increasing piece.  The positive contribution
  is bounded by the variation on the increasing piece.  The sampled endpoint
  values are bounded by
$\log(4(g+1))$ near the turning cell and by
$\log(6(g+1))$ at the right endpoint, so this variation is at most
$\log(6(g+1))$.  Thus the complete discrepancy satisfies
\[
 \sum_{k=0}^{g-1}f(k+1)-\int_0^g f(u)\,du
 \le 2+\log(6(g+1)).
\]
This is the bounded-variation estimate for the truncated function $H$; it
also covers the cases where the zero of $f$ lies outside $[0,g]$.  We retain
the looser uniform quantity
\[
 \beta_m=4\log(g+1)+9+10\log2
\]
below.  Restoring the factors $\log m$ gives the cell integral in the
contour-product estimate.
Thus, for every $z\in T_\delta$,
\begin{equation}\label{eq:p0-q-contour-product}
 \log|F_{g,q}(mz)|\le m\phi_q(z)+\beta_m.
\end{equation}
The scalar factor $q^{mz}$ is exactly the affine field term and adds no
discrepancy.  No lower bound for this discrepancy is used.
Since $g+1\le m$ for $m\ge8$, the contour
and circle errors together contribute $2m\log m+m\beta_m$.

Combining the preceding terms with the fixed Vandermonde normalization and
discarding favorable powers of $m$ gives, for every configuration,
\[
 \log Q_{m,q}(S,R,Z)
 \le -m^2\mathcal E_q(\kappa_S+\nu_\delta,\mu_R)
   +10m\log m+c_*m+\frac13
   +|S|\log\log q,
\]
where the real-cell constant is
\begin{align*}
c_*={}&4\log(2\pi)+9+\frac{21}{2}\log2+\log336+\log(35/3)
                         +\frac5{2d_+}+\frac1{2d_-}\\
 <{}&49.978.
\end{align*}
The displayed logarithm bounds can be checked with the rational series
\[
 \log x=2\sum_{j=0}^{N}\frac{u^{2j+1}}{2j+1}+R_N,\qquad
 0<R_N<\frac{2u^{2N+3}}{(2N+3)(1-u^2)},\quad
 u=\frac{x-1}{x+1}.
\]
For $x=710/113$ use $N=50$; for
$x=2,21/16,35/24,39/32$ use $N=16$.  The rational series gives
\[
 \log(710/113)<1.83787716,\quad \log2<0.69314719,
 \quad \log(21/16)<0.27193372,
\]
\[
 \log(35/24)<0.37729424,\qquad \log(39/32)<0.19782575.
\]
The classical rational bound $\pi<355/113$ therefore gives
$\log(2\pi)<1.83787716$; hence
$\log336<8(0.69314719)+0.27193372$ and
$\log(35/3)<3(0.69314719)+0.37729424$.
Substitution of the exact rational $d_+=751/5000$ and
$d_-=1749/5000$ gives $c_*<49.977264<49.978$ and
$c_*+\log39+1+A<55.920<56$.
This is the base-$2$ cell constant; the $q$-dependent cell shifts above are
favorable.  The only $q$-dependent finite factor not absorbed by the affine
energy is the diagonal multiplier $(\log q)^{|S|}$.  In the pointwise
envelope it appears as the exact term $|S|\log\log q$.  At $q=2$ this term
is negative, so the base-$2$ estimate omits that favorable factor; whenever an
upper bound uniform in $q$ is taken below, it is replaced by
$|S|(\log\log q)_+$.  Uniformly
for $2\le q\le11$,
\[
 \log q<3,\qquad 0<\log q<3<5,\qquad \log(\log q)<\log3<2,
\]
so $|S|(\log\log q)_+<2m$.  For all bases we therefore replace $56m$ by at
most $58m$; this
finite correction is $O(m)$ and does not change the limiting energy rate.

\subsection{Tube transfer, frequency tail, and summation}

Let $D$ be the interior of the ellipse.  For $z$ inside $D$ use its interior
harmonic measure $\omega_z$ on $\Gamma$; outside $D$ use exterior harmonic
measure of total mass one; on $\Gamma$ use $\delta_z$.
Let $W$ be the exterior inverse of
$\mathcal J(w)=\xi+cw+d/w$, where $c=(A+B)/2$ and $d=(A-B)/2$.  Define
$e(z)=0$ in $D$ and $e(z)=\log|W(z)|$ outside.  The factorization
$\mathcal J(w_1)-\mathcal J(w_2)=(w_1-w_2)(c-d/(w_1w_2))$ and the Green
identity imply
\[
 \int K_0(x,y)\,d\omega_z(y)\le K_0(x,z)+e(z),\qquad
 I(\omega_z,\omega_w)\le K_0(z,w)+e(z)+e(w),
\]
where $K_0=-\log|\cdot|$ and $I$ denotes its mutual energy.  The interior
deficit is the Green function of $D$; the exterior deficit is the
nonnegative disk Green function; the cross-boundary deficit is zero.
For $w=\rho e^{iu}$, $\rho\ge1$, the equipotential $|W(z)|=\rho$ is
the confocal ellipse
\[
 \Gamma_\rho:
 z=\xi+(c\rho+d/\rho)\cos u+i(c\rho-d/\rho)\sin u.
\]
It equals $\Gamma$ at $\rho=1$.  If $z=\mathcal J(w)$ and
$\zeta=\mathcal J(v)\in\Gamma$, $|v|=1$, the exact factorization gives
\[
 |z-\zeta|=|w-v|\left|c-\frac{d}{wv}\right|
 \ge(\rho-1)(c-d/\rho)\ge B(\rho-1).
\]
Taking the infimum over $\zeta$ and using $\log\rho\le\rho-1$ proves
\[
 e(z)\le\operatorname{dist}(z,\Gamma)/B\le\delta/B
 \qquad(z\notin D,\ z\in T_\delta).
\]

For $x\in[0,2]$, set
$\chi(z)=1-(\Re z-\xi)^2/A^2-(\Im z)^2/B^2$.  On $D$,
$|\nabla\chi|\le2/B$, while $\chi(x)>1/5$.  The level curve
$\chi=1/8$ is more than $1/32$ from $[0,2]$, so the Green function there is
at most $\log160<6$.  The barrier $48\chi$ on the collar gives
\[
 g_D(z,x)\le96\delta/B\qquad(z\in D\cap T_\delta,\ x\in[0,2]).
\]
Since $A_0$ has source mass three and $\lambda\Re z$ is harmonic, the
interior field-sweep loss for $A_0-\phi_q$ is at most $288\delta/B$.
The term $-V$ is favorable in this estimate: $V$ is subharmonic, with
distributional Laplacian $2\pi$ times the source measure on
$[-1/2,0]$, so its harmonic sweep is at least its interior value and the
$-V$ contribution can be dropped from an upper bound.
For exterior $z=\mathcal J(w)$, the swept mean coordinate is
$\xi+c/\overline w+d/w$, within $(A+B)\delta/B$ of $z$.
For the base-$2$ terms the source bound is
$[1/2+(A+B)\log2]\delta/B<3\delta/B$: $A_0$ changes by exactly
$-3e(z)$, and the remaining coordinate displacement is at most
$(A+B)\delta/B$.  The additional affine term is at most
$\lambda(A+B)\delta/B<6\delta/B$ because
$\lambda=\log(q/2)<2$ and $A+B<3$.

The logarithmic source in $\phi_q$ needs a separate collar estimate.  Put
\[
 V(z)=\int_0^{1/2}\log|z+s|\,ds.
\]
For $0<\delta\le1/64$ and an exterior point with
$\operatorname{dist}(z,\Gamma)\le\delta$, the explicit primitive
\[
 V(z)=\Re\bigl((z+1/2)\operatorname{Log}(z+1/2)
              -z\operatorname{Log}z-1/2\bigr)
\]
on either side of the source segment gives
\[
 \nabla V(z)=\operatorname{Log}(z+1/2)-\operatorname{Log}z.
\]
For $z=\xi+A\cos u+iB\sin u$ on $\Gamma$, minimizing
$|z|^2$ gives its minimum at $u=\pi$, namely
$\xi-A=751/5000$; the derivative has no other critical point because
$A\xi>|A^2-B^2|$.  The disk of this radius about $0$ lies inside the ellipse,
so every exterior point satisfies $|z|\ge751/5000$.  For the radial
comparison path, $h\le\delta/B$ and
\[
 |\mathcal J((1+h)e^{iu})-\mathcal J(e^{iu})|
 \le(c+d)h=Ah\le\frac AB\delta
 =\frac{12783}{11450}\delta<1.117\delta.
\]
This is about $1.116\delta$ and is less than $0.02$ when
$\delta\le1/64$.  Thus all points used in the estimate have
$\Re z\ge-751/5000-(12783/11450)/64>-1/4$.  Therefore
\[
 \frac{|z+1/2|}{|z|}
 \le1+\frac1{2|z|}
 \le1+\frac{2500}{751}
 =\frac{3251}{751}<4.33,
 \qquad
 |\arg(z+1/2)-\arg z|\le\pi.
\]
Since $|z+1/2|^2-|z|^2=\Re z+1/4>0$, the ratio is at least one, and
the modulus term in the difference of logarithms is at most $\log(4.33)$.
The angle difference is at most $\pi$ because the two points lie on the
same horizontal line.  Also $e^{3/2}>4.33$ by the first six terms of its
positive series, so $\log(4.33)<3/2$.  Thus, on both sides of the cut,
\[
 |\nabla V(z)|
 \le\sqrt{(3/2)^2+(22/7)^2}<3.72.
\]
At the contact point $z_0=\xi-A=-751/5000$ the classical gradient is not
defined, but these are the bounded one-sided traces and they give the
required Lipschitz bound across the smooth ellipse and the source cut.  The
Joukowski boundary parametrization has speed at most $4/3$, so its
pullback is $6$-Lipschitz.  If the exterior pole is $w=\rho e^{iu}$, write
$\rho=1+h$.  The elementary estimate
\[
 \int_{-\pi}^{\pi}|t|P_{1+h}(t)\,\frac{dt}{2\pi}
 \le \frac{3\pi h}{8}\log(1+4/h^2)
\]
is used for $0<h\le1$ and follows from
$1-\cos t\ge2t^2/\pi^2$.  In the present collar
$h\le\delta/B\le25/1832<1$, so this is the required range.  Since
$\operatorname{dist}(\mathcal J(w),\Gamma)\ge Bh$ and $\delta\le1/64$,
these bounds imply
\[
 \left|\int_\Gamma V(\zeta)\,d\omega_z(\zeta)-V(z)\right|
 \le20\delta\log(1/\delta).
\]
Thus the exterior field loss is at most
$9\delta/B+20\delta\log(1/\delta)$.  Combining this with the interior
bound gives a uniform field loss at most
$288\delta/B+20\delta\log(1/\delta)$.

The three interaction terms change by at most $3b$, where
$b=\int e\,d\nu\le|\nu|\delta/B$.  Since $|\nu|\le1$, their total loss is
 at most $3\delta/B$.  Consequently the full transfer cost is at most
$291\delta/B+20\delta\log(1/\delta)
 <300\delta+20\delta\log(1/\delta)$ because $B=229/200>1$.
Thus
\[
 \mathcal E_q(\kappa+\nu_\delta,\mu)
       \ge\alpha(\lambda)-300\delta-20\delta\log(1/\delta),
       \qquad \frac{291}{B}<300.
\]
The dual problem is posed on the printed support interval \([1/2,20]\).
The frequency-cutoff argument below first replaces the full frequency sum by
the range \(R\subset\{g+1,\ldots,20m\}\), with a positive multiplicative
tail factor.  Hence every finite configuration used in the variational bound
has its frequency measure in this compact interval; the omitted tail is
accounted for separately in \eqref{eq:p0-frequency-cutoff}.

For completeness, the unbounded frequency sum is reduced to a finite one
before counting configurations.  For $H\ge g+1$, define
\[
 E_{m,q}(H)=\sum_{S\subset I_m}\frac1{d!}
 \int_{(m\Gamma)^d}
 \sum_{\substack{R\subset\{g+1,\ldots,H\}\\|R|=m}}
 Q_{m,q}(S,R,Z)\prod_{z\in Z}\frac{|dz|}{2\pi}.
\]
The Cauchy--Binet expansion~\cite{Gantmacher} gives
$|G_{m,q}|\le E_{m,q}(\infty)$.
For fixed $S,Z$, its frequency-dependent
factor is
\[
 f_q(R)=
 \frac{\prod_{r\in R}w_{r,q}\,\Delta(R)^2
       \prod_{r\in R,z\in Z}|r+z|}
      {\prod_{r\in R,i\in I_m}(r+i)
       \prod_{r\in R,j\in J_m\setminus S}(r+j)}.
\]
If $r=\max R>20m$, replace $r$ by the least missing
$s\in\{g+1,\ldots,g+m\}$, so $s\le3m/2$.  The exact ratio of
the frequency weights, both frequency Vandermonde factors, and all
contour-frequency factors gives
\[
\frac{f_q(R)}{f_q(R')}\le
      12500^m\frac{m^2}{4r^2}w_{r,2}.
\]
For the constant, at most two integers lie at each distance from $s$, so
\[
 \prod_{u\in R\setminus\{r\}}|s-u|
 \ge d_m:=
 \left\lfloor\frac{m-1}{2}\right\rfloor!
 \left\lceil\frac{m-1}{2}\right\rceil!
 \ge e^{-1}\left(\frac{m}{2e}\right)^{m-1}.
\]
Cancel the common factors in the two Cauchy determinants.  The two
frequency Vandermonde ratios contribute the squared products of distances
from $r$ and $s$; each contour-frequency factor is bounded by the
corresponding ratio involving $r+g$ and $s+g$, and the fixed-grid factors
are bounded by $5^m$.  The explicit denominator bounds
$s+i\le11m/4$, $s+j\le7m/2$, and ellipse bounds
$\Re z>-m/6$, $|z|<5m/2$ give
\[
 \frac{f_q(R)}{f_q(R')}
 \le \frac{w_{r,q}}{w_{s,q}}
       \left(\frac{77}{8}\right)^m
       \left(\frac{189}{16}\right)^{m-|S|}
       \frac{m^{2m}}{r^2d_m^2}.
\]
For $r>s$, $w_{r,q}/w_{s,q}\le w_{r,2}/w_{s,2}$ and
$w_{s,2}\ge2^{-3m/2}$ because $s\le3m/2$ and
$\binom{s-1}{g}\ge1$.  Substituting this and $e<3$ into the last
display gives the asserted $12500^m m^2w_{r,2}/(4r^2)$ bound.  This is the
full constant calculation; it does not require an estimate for $w_{s,q}$.
The exact base comparison is
\[
 \frac{w_{r,q}/w_{s,q}}{w_{r,2}/w_{s,2}}
       =\left(\frac2q\right)^{r-s}\le1.
\]
\begin{proof}[Frequency cutoff]\label{eq:p0-frequency-cutoff}
\[
 \vartheta_m=\frac{m}{1600}12500^m\sum_{r>20m}w_{r,2}.
\]
Indeed, the replacement map has at most $m$ preimages for each pair
$(R',r)$, and $m^2/(4r^2)<1/1600$ when $r>20m$.  For $t=15/8$,
\[
 \sum_{r>20m}w_{r,2}
 \le(8/15)^{20m}\sum_{r\ge g+1}w_{r,2}(15/8)^r
 =(8/15)^{20m}15^{g+1}.
\]
Since $g=m/2$ and
$12500\sqrt{15}<50000$, this gives
\[
 \vartheta_m\le\frac{15m}{1600}
       \left[50000(8/15)^{20}\right]^m
 <\frac3{320}2^{-m}<\frac12.
\]
The strict inequality uses $200000\cdot8^{20}<15^{20}$ and
$m2^{-m}\le1$.  If $Z_M$ denotes the finite frequency sum through $M$,
the replacement map gives $Z_M-Z_{20m}\le \vartheta_mZ_M$.  Hence
$Z_M\le(1-\vartheta_m)^{-1}Z_{20m}\le(1+2\vartheta_m)Z_{20m}$.
\[
 2\vartheta_m<\frac3{160}2^{-m}.
\]
Monotone convergence as $M\to\infty$ now proves
\[
 E_m(\infty)\le\left(1+\frac3{160}2^{-m}\right)E_m(20m).
\]
\end{proof}
There are $2^m$ pinned subsets and at most $\binom{39m/2}{m}$ retained
frequency sets.  The normalized contour length is at most $mA$, hence
$(mA)^d/d!\le e^{mA}$.  Using $c_*<49.979$ and
$c_*+\log39+1+A<56$, the finite configuration bound, the tube transfer
with $\delta=m^{-2}$, and the tail factor imply the following.

\begin{lemma}[Uniform smoothing bound]\label{lem:p0-smoothing-transfer}
For every integer $q\in\{2,\ldots,11\}$ and every $m\in4\mathbb N$,
 $m\ge8$,
\begin{equation}\label{eq:p0-q-uniform-det}
 \log|G_{m,q}(\Li_2(1/q))|
 \le-\alpha(\log(q/2))m^2+10m\log m+58m+301+40\log m.
\end{equation}
The frequency quantifier before replacement is the full set
$R\subset\{g+1,g+2,\ldots\}$ of cardinality $m$.  The constant $58$ is
uniform in $q$; the constant $301$ includes the tube cost $300$, the
real-cell remainder $1/3$, and the exponentially small frequency-tail
factor.  The additional $40\log m$ is the explicit logarithmic collar cost.
\end{lemma}

\begin{proof}
For each finite configuration the cell measures above are admissible in the
$m^{-2}$ tube problem.  Tube transfer bounds their energy below by
$\alpha(\lambda)-300m^{-2}-40m^{-2}\log m$.  Inserting this into the pointwise determinant
envelope gives a $300+40\log m$ additive loss.  Summing the $2^m$ pinned choices,
the $\binom{39m/2}{m}$ retained frequency choices, and the contour integral
 adds at most $(\log39+1+A)m$ beyond $c_*m$.  The first contribution is
strictly less than $56m$ by the displayed $c_*+\log39+1+A<56$; the extra
diagonal factor is strictly less than $e^{2m}$.  Hence the coefficient of
$m$ is strictly less than $56+2=58$.  The cell remainder is $1/3$.  For
$m\ge8$ the tail factor obeys
\[
 \log\left(1+\frac3{160}2^{-m}\right)
 \le\frac3{160}2^{-m}\le\frac3{40960},
\]
so $300+1/3+3/40960<301$.  Taking the
monotone limit of the finite frequency cutoffs proves the estimate for the
full half-line.
\end{proof}

\subsection{Endpoint certificates and the interval comparison}

The exact reference masses and thresholds are printed in
Tables~\ref{tab:q2data} and~\ref{tab:q11data}.  Their directed dual checks
prove
\begin{equation}\label{eq:p0-q-endpoint-bounds}
 \alpha(0)\ge\frac{248035}{125000},\qquad
 \alpha(\log(11/2))\ge\frac{77597}{100000}.
\end{equation}
These are lower bounds for the full continuum infimum, not empirical
minimizers.  In each certificate the contour inequality is checked on the
whole ellipse; the pinned positive-part integral is checked on its full
support; the frequency positive-part integral is checked through $81/10$.
The endpoint inequalities at $81/10$ are part of the directed integrals
recomputed from the printed measures: they give $W_\mu(81/10)>31/10$
for $q=2$ and $W_\mu(81/10)>56/5$ for $q=11$, while the respective thresholds
are $1427/500$ and $1217/200$.

The infinite tail has a direct bound that uses only the masses and supports
of these two references.  Both $\eta_0$ have mass one and support in
$E\cup\Gamma$, whose real part is at least $-751/5000$; both $\mu_0$ have
mass one and support in $[1/2,9/2]$.  Also $B_0'$ is the Cauchy transform
of a positive measure of mass three on $[0,2]$.  Differentiating the
definition of $W_\mu$ therefore gives, for $t\ge81/10$,
\begin{align*}
 W_\mu'(t)
 &=B_0'(t)+\log q-\log\frac{t}{t-1/2}
   -2\int\frac{d\mu_0(s)}{t-s}
   -\int\Re\frac{d\eta_0(z)}{t+z}\\
 &\ge H(t):=\frac3{t+2}+\log q-\log\frac{t}{t-1/2}
       -\frac2{t-9/2}-\frac1{t-751/5000}.
\end{align*}
For $t\ge81/10$, differentiation yields
\[
H'(t)=-\frac3{(t+2)^2}+\frac1{2t(t-1/2)}
       +\frac2{(t-9/2)^2}+\frac1{(t-751/5000)^2}>0,
\]
because the last two positive terms exceed $3/t^2$, while
$3/(t+2)^2<3/t^2$.  Thus $H(t)\ge H(81/10)$.  The elementary bounds
$\log q\ge\log2>2/3$, $\log(81/76)<1/13$, and
$5000/39749<13/100$ give
\[
W_\mu'(t)\ge H(81/10)>
\frac{30}{101}+\frac{69}{100}-\frac1{13}-\frac59-\frac{13}{100}
=\frac{66338}{295425}>\frac{31}{684}.
\]
Here we used the elementary rational-series bounds
$\log2>69/100$, $\log11>23/10$, $\log(81/76)<1/13$, and
$5000/39749<13/100$; the $q=11$ lower bound is larger than $1$.
Consequently $W_\mu'(t)>31/684>1/50$ for $q=2$ and is $>1$ for
$q=11$.  Hence $W_\mu(t)>h_\mu$ on the entire omitted ray, so $T$ has no
omitted tail.  The directed endpoint records are summarized in
Table~\ref{tab:p0-q-bounds}; the exact mass lists are the two preceding
reference tables.

For auditability, the scalar inputs and the outward rational penalties used
by the two records are listed here (the displayed finite decimals are
outward lower bounds for $Q_{\log}$):
\begin{table}[ht]
\centering
\caption{Scalar data in the two directed endpoint certificates.}
\label{tab:p0-q-bounds}
\renewcommand{\arraystretch}{1.12}
 {\normalsize
\begin{tabular}{@{}c l@{}}
\toprule
$q$ & exact thresholds and directed penalty bounds\\
\midrule
$2$ & $h_\eta=-3347002/1000000$, $h_\mu=1427/500$;\\
    & $Q_{\log}>-0.425508$,\quad
      $P\le0.087950$,\\
    & $T\le0.122962$;\\[2pt]
$11$ & $h_\eta=-3359/500$, $h_\mu=1217/200$;\\
    & $Q_{\log}>-0.635537$,\quad $P\le0.000001$,\\
    & $T\le1.069189$.\\
\bottomrule
\end{tabular}}
\end{table}

The arithmetic rate used for the positive bases is
\[
 \tau_+(\lambda)=K+\frac1{16}\log(q-1)-\frac{13}{16}\log q,
                    \qquad q=2e^\lambda,
\]
where the arithmetic estimate gives $K\le1271001/500000$.  Differentiation
with respect to $\lambda$ gives
\[
 \tau_+''(\lambda)=-\frac{q}{16(q-1)^2}<0.
\]
Its tangent at $q_0=63/25$ has exact slope
\[
 \tau_+'(\log(q_0/2))=-\frac{431}{608}.
\]
Concavity places $\tau_+$ below this tangent on the whole interval.  Directed
logarithm evaluation with the stated rational upper bound for $K$ gives the
endpoint comparison:
\begin{center}
\renewcommand{\arraystretch}{1.12}
 {\normalsize
\begin{tabular}{@{}c c c c@{}}
\toprule
$q$ & directed $\alpha$ lower bound & tangent upper bound for $\tau_+$
     & directed gap lower bound\\
\midrule
$2$ & $248035/125000$ & $1.981041880111$ & $>3/1000$\\
$11$ & $77597/100000$ & $0.772577360778$ & $>3/1000$\\
\bottomrule
\end{tabular}}
\end{center}
The displayed decimal upper bounds are outward roundings of the directed
intervals; the corresponding directed gaps are respectively
$0.00323811988930149\ldots$ and $0.00339263922279043\ldots$.

Let $L(\lambda)$ be the chord joining the two lower bounds in
\eqref{eq:p0-q-endpoint-bounds}.  By concavity of $\alpha$,
$\alpha(\lambda)\ge L(\lambda)$.  Let $T(\lambda)$ be the tangent to
$\tau_+$ at $q_0=63/25$.  The two directed endpoint checks show
$L-T>3/1000$ at both ends.  Since $L-T$ is affine in $\lambda$, it is
greater than $3/1000$ throughout
$0\le\lambda\le\log(11/2)$.  Also $\tau_+(\lambda)\le T(\lambda)$.
We have proved
\begin{equation}\label{eq:p0-q-positive-margin}
 \alpha(\lambda)-\tau_+(\lambda)>\frac3{1000}
                 \quad(0\le\lambda\le\log(11/2)).
\end{equation}
The determinant estimate \eqref{eq:p0-q-uniform-det} is a separate
smoothing estimate; the fixed-base integerizer and nonvanishing sequence are
applied separately at each integer $q$.  Thus this section uses
no interpolation of integerizers, seed primes, or endpoint reference
measures.


\section{Positive bases from 2 to 11}\label{sec:positive}
The fixed feasible set and affine fields give concavity in
\(\lambda=\log(a/2)\), as proved in Section~\ref{sec:p0-analytic}.  The
endpoint certificates use exact rational measures on
\[
 (\xi,A,B)=\left(\frac{11281}{10000},\frac{12783}{10000},\frac{229}{200}\right),
\]
with 64 contour nodes and Fej\'er degree 48.  They give
\[
 \alpha_2\ge\frac{248035}{125000},\qquad
 \alpha_{11}\ge\frac{77597}{100000}.
\]
The arithmetic comparison uses \(K<1271001/500000\), tangent point
\(q_0=63/25\), and slope \(-431/608\).  The directed endpoint gaps exceed
\(0.0032381198893\) and \(0.0033926392227\), respectively.  These checks and
the tangent majorant cover \(2\le a\le11\).  The integerization and
nonvanishing sequence are applied separately at each integer base.  The endpoint measures
are printed in Appendix~\ref{sec:positive-data}.

\section{Negative bases $-7\le a\le-2$}\label{sec:negative-middle}\label{sec:minus2}
The signed determinant construction uses the
\([C;B-TA;V]\) determinant.  At \(a=-2\), its directed analytic bound is
\(\alpha_2>4.34875\), while the fixed arithmetic rate is
\(\tau_2=A_{40}-(53/48)\log2=4.317643952609845\ldots\), giving margin
\(>0.03110\).  The denominator estimate, prime sum, analytic bound and
prime-power nonvanishing sequence includes \(q=2\) and is instantiated at each
remaining base.
For \(a=-q\), the field acquires the term \(-\pi\operatorname{Im}z\), so the
full signed contour in Section~\ref{sec:v-analytic} is required.  For each
\(q=2,\ldots,7\), the contour geometry, printed reference data, and directed
field bounds, together with the smoothing estimate in Section~\ref{sec:v-analytic}, give a
strict \(\alpha_q>\tau_q\), while Sections~\ref{sec:v-arithmetic} and
\ref{sec:nonvanishing} supply the matching integerizer and eventual
nonvanishing sequence.
Thus all six bases \(-7\le a\le-2\) are closed individually, without an
interpolation claim across \(q\).

\section{Arithmetic for the signed determinant}
\label{sec:v-arithmetic}

This section supplies the arithmetic interface for the negative
\(V\)-family.  The base is fixed before \(m\) tends to infinity.  Thus the
statement is made separately for \(q=2,3,4,5,6,7\), with \(a=-q\); it is
not a theorem uniform in an unbounded parameter \(q\).  The proof has four
parts: an exact polynomial divisor, residue-block valuations away from the
fixed primes, an outer pole assignment, and a weighted prime-number-theorem
sum.  All subsequent limits follow from the displayed residue valuations and
outer assignment bounds.

Put \(m=4n\), \(s=k=h=n\), and
\[
 I=[2n,6n),\qquad J=[n,10n),\qquad
 \Lambda=\{1,\ldots,n\},\qquad N=10n-1,\quad U=11n-1.
\]
For an integer parameter \(a\), define
\[
 t_j=\sum_{r=1}^{j}\frac{a^{j-r}}r,\qquad
 b_j=\sum_{r=1}^{j}\frac{a^{j-r}}{r^2},
 \qquad t_0=b_0=0,
\]
and, for \(i\in I\), \(j\in J\), \(\ell\in\Lambda\),
\[
 C_{ij}=\begin{cases}(a^i-a^j)/(j-i),&i\ne j,\\0,&i=j,\end{cases}
 \quad
 B_{ij}=\begin{cases}(t_i-t_j)/(j-i),&i\ne j,\\b_i,&i=j,\end{cases}
\]
\[
 A_{ij}=a^i\delta_{ij},\qquad V_{\ell j}=\frac1{j+\ell},\qquad
 G_m(T;a)=\det\begin{bmatrix}C(a)\\ B(a)-TA(a)\\ V\end{bmatrix}.
\]
The determinant has size \(2m+n\).  All valuations below are coefficientwise
in \(T\): \(v_p(F)=\min_r v_p([T^r]F)\).

\subsection{Denominator and the formal divisor}

Let \(d_r=\operatorname{lcm}(1,\ldots,r)\) and
\[
 \Delta_m=d_N^{3m}d_U^n,\qquad
 D_m=\left\lceil\frac{53m^2}{48}\right\rceil-m
      =\left\lceil\frac{53n^2}{3}\right\rceil-4n.
\]

\begin{lemma}[denominator]
\label{lem:v-denom}
For an indeterminate \(a\),
\[
 \Delta_mG_m(T;a)\in\mathbb Z[a,T].
 \]
\end{lemma}

\begin{proof}
Every denominator in a \(C\)-entry divides \(d_N\).  In an off-diagonal
\(B\)-entry one harmonic denominator and one divided-difference denominator
occur, while \(b_i\) has two harmonic denominators; hence \(d_N^2\) clears
each lower row.  The denominators of the \(V\)-rows divide \(d_U\).
Multiplying the rows and expanding the determinant proves the assertion.
\end{proof}

\begin{lemma}[formal $a$-divisor]
\label{lem:v-formal-divisor}
Coefficientwise in \(T\),
\[
 \Delta_mG_m(T;a)\in a^{D_m}(a-1)^m\mathbb Z[a,T].
\]
\end{lemma}

\begin{proof}
The following argument is coefficientwise in \(\mathbb Q[[a]][T]\), so
every use of an infinite series can be replaced by a finite truncation for a
fixed coefficient.  For \(u_t(r)=(r-t)^{-1}\) when \(r\ne t\), and
\(u_t(t)=0\), put \(S_1(a)=\sum_{v\ge1}a^v/v\) and
\(S_2(a)=\sum_{v\ge1}a^v/v^2\).  Direct coefficient comparison, including
the diagonal entries, gives
\[
 C_{ij}=a^i u_i(j)+a^j u_j(i),
\]
\[
 (B-S_1C)_{ij}
 =\sum_{t\ge0}a^tu_t(i)u_t(j)-\frac{a^i+a^j}{(i-j)^2}
 \quad(i\ne j),
\]
and
\[
 (B-TA)_{ii}=\sum_{t\ge0}a^tu_t(i)^2-a^i(T+S_2(a)).
\]
The last formula is the diagonal term and is the reason no off-diagonal
limiting convention is needed.

Expand first along the \(n\) rows of \(V\).  In each remaining \(2m\)-row
minor, use the preceding rank-one expansion.  A nonzero term selects
\(2m=8n\) distinct features.  If \(b,x,y,z\) count, respectively, the
series features, the upper boundary features, the lower boundary features,
and the column boundary features, then \(b+x+y+z=8n\).  Distinctness gives
the lower bound
\[
 \Psi\ge {b\choose2}+{x\choose2}+x\max(2n,b)
       +2ny+{y\choose2}+nz+{z\choose2}.
\]
Here the \(b\) series indices are distinct; an upper boundary index cannot
be one of them because the two corresponding right vectors coincide.  The
\(x\) and \(y\) indices are distinct members of \(I=[2n,6n)\), while the
\(z\) column indices are distinct members of \(J=[n,10n)\).  These are the
only features used in the count, so a collision can only delete a term and
cannot lower the displayed lower bound.
Writing \(B_0=b+x\), the right side is at least
\[
 Q(B_0,y)=B_0^2+B_0y+y^2-9nB_0-7ny+40n^2-4n
 =\frac{53n^2}{3}-4n+r^2+rs+s^2
\]
for suitable real \(r,s\).  Since the order is integral, every coefficient
of every \(2m\)-minor has order at least \(D_m\).  The \(V\)-minors are
rational constants independent of \(a,T\), so expansion along \(V\) gives
\(a^{D_m}\mid G_m(T;a)\) in \(\mathbb Q[a,T]\).  Combining this with the
denominator lemma and monic division gives
\(a^{D_m}\mid\Delta_mG_m\) in \(\mathbb Z[a,T]\).

Finally every \(C\)-entry is divisible by \(a-1\) in \(\mathbb Q[a]\),
including the zero diagonal convention.  There are \(m\) such rows, hence
\((a-1)^m\mid G_m\).  The two monic factors \(a^{D_m}\) and \((a-1)^m\)
are coprime, so their product divides the integral polynomial
\(\Delta_mG_m\).  This proves the lemma for every \(n\ge1\), not only for
the finite verification values.
\end{proof}

For \(a=-q\), the formal divisor implies the exact coefficientwise bound
\[
 v_p([T^r]G_m(T;-q))\ge v_p(q)D_m+v_p(q+1)m
       -3m\lfloor\log_pN\rfloor-n\lfloor\log_pU\rfloor .
\]
The two credited prime sets are disjoint because \(\gcd(q,q+1)=1\).  Notice
that this argument includes composite \(q\) and powers \(p^e\); no
square-free assumption is being made.
For reference, the finite exceptional sets in the five middle cases are
\[
\begin{array}{c|c@{\qquad}c|c}
q&\mathcal P_q=\{p:p\mid q(q+1)\}&q&\mathcal P_q\\ \hline
3&\{2,3\}&4&\{2,5\}\\
5&\{2,3,5\}&6&\{2,3,7\}\\
7&\{2,7\}&&
\end{array}
\]
The \(q=2\) row is \(\{2,3\}\), and is treated by the same rule.

\subsection{The $p$-unit residue blocks}

Fix an odd prime \(p\nmid q(q+1)\) with \(p^2>\max(N,U)\), and write
\(i=u+kp\), \(j=u+\ell p\), \(0\le u<p\).  Since \(a=-q\) and
\(a^p\equiv a\pmod p\), splitting the defining sums at multiples of \(p\)
gives
\begin{align*}
 p t_{u+kp}
 &=\sum_{\substack{1\le r\le u+kp\\p\nmid r}}
       \frac{p a^{u+kp-r}}r
   +\sum_{s=1}^{k}\frac{a^{u+p(k-s)}}s
   \equiv a^u t_k\pmod p,\\
 p^2 b_{u+kp}
 &=\sum_{\substack{1\le r\le u+kp\\p\nmid r}}
       \frac{p^2a^{u+kp-r}}{r^2}
   +\sum_{s=1}^{k}\frac{a^{u+p(k-s)}}{s^2}
   \equiv a^u b_k\pmod p,\qquad
 p^2TA_{ij}\equiv0\pmod p.
\end{align*}
Here \(k,\ell<p\) follows from \(i\le N\), \(j\le U\), and
\(p^2>\max(N,U)\), so each displayed denominator \(s\) is a \(p\)-unit.
Thus the scaled entries satisfy
\[
 pC_{u+kp,u+\ell p}\equiv a^uC_{k\ell}\pmod p,\qquad
 p^2B_{u+kp,u+\ell p}\equiv a^uB_{k\ell}\pmod p,
\]
whereas
\[
 pV_{r,u+\ell p}\equiv0\quad\text{unless }r+u=pw,\qquad
 pV_{r,u+\ell p}=\frac1{\ell+w}\quad(r+u=pw).
\]
The denominators in the last expression are units because \(p^2>U\).
Cross-residue \(C\) and \(B\) entries vanish after the displayed row
scalings.  The reduced residue-\(u\) block has \(2|K_u|+|W_u|\) rows and
\(|L_u|\) columns.  Therefore its rank over \(\mathbb F_p\) is at most
\[
 \min(2|K_u|+|W_u|,|L_u|)
\]
where
\[
 K_u=\{k:1/2\le u/m+kp/m<3/2\},\quad
 L_u=\{\ell:1/4\le u/m+\ell p/m<5/2\},
\]
\[
W_u=\{w\ge1:0<w p/m-u/m\le1/4\}.
\]
Let $\operatorname{rk}_u$ denote the actual rank of this reduced block.
After permuting rows and columns by residue class, the globally scaled square
matrix of order $2m+n$ is block diagonal with precisely these blocks.  Hence
its rank modulo $p$ equals $\sum_u\operatorname{rk}_u$, and its determinant
has valuation at least
\[
 (2m+n)-\sum_u\operatorname{rk}_u.
\]
The total row-scaling factor is $p^{3m+n}$: the $C$, $B$, and $V$ rows were
scaled by $p$, $p^2$, and $p$, respectively.  Undoing this factor gives
\[
 v_p(G_m(T;-q))\ge -\left(m+\sum_u\operatorname{rk}_u\right).
\]
These valuations are coefficientwise in $T$.  Indeed, the Gauss valuation on
$\mathbb Q_p(T)$ is discrete with residue field $\mathbb F_p(T)$; Smith
normal form over its valuation ring gives
$v_p(\det M)\ge (2m+n)-\operatorname{rank}(M\bmod p)$ for every integral
square matrix $M$ of order $2m+n$.  Thus only the actual block ranks
$\operatorname{rk}_u$ are credited,
not a count of non-divisible rows.  Summing the block ranks gives the
exponent used below.
The only rank saving deliberately credited in the valuation sum is the profile
\(K=\{0\},L=\{0,1\},W=\varnothing\).  After dividing its rows by the
unit \(a^u\), its quotient matrix is
\[
 \begin{pmatrix}0&1-a\\0&-1\end{pmatrix},
\]
which has rank one over every residue field.  This is a universal identity;
no accidental rank deficiency of another profile is credited.

For \(\rho=p/m\), define
\[
 K(x,\rho)=\{k\ge0:1/2\le x+k\rho<3/2\},\quad
 L(x,\rho)=\{\ell\ge0:1/4\le x+\ell\rho<5/2\},
\]
\[
 W(x,\rho)=\{w\ge1:0<w\rho-x\le1/4\},
\qquad 0\le x<\rho,
\]
and define
\[
\begin{aligned}
 \mathfrak r(x,\rho)
   &=\min(2|K(x,\rho)|+|W(x,\rho)|,|L(x,\rho)|)\\
   &\quad-\mathbf 1_{\{K(x,\rho)=\{0\},\,
       L(x,\rho)=\{0,1\},\,W(x,\rho)=\varnothing\}},\\
 \mathfrak c(\rho)&=1+\int_0^\rho\mathfrak r(x,\rho)\,dx .
\end{aligned}
\]
with the rank-one profile subtracted on its exact active cells.  The endpoint
conventions are left-closed/right-open for \(K,L\), and
\(w\rho-1/4\le x<w\rho\) for \(W\).  Every wall has the form
\(x=c+b\rho\) with \(4c,b\in\mathbb Z\).  Since \(m=4n\), the function is
constant on each lattice interval \([u/m,(u+1)/m)\), and hence
\[
 \sum_{u=0}^{p-1}\mathfrak r(u/m,p/m)
   =m\int_0^{p/m}\mathfrak r(x,p/m)\,dx.
\]
The residue-block lemma is therefore
\[
 v_p(G_m(T;-q))\ge -m\mathfrak c(p/m).
\]
The proof uses only the p-unit conditions and the index geometry, so it is
independent of the value of the fixed \(q\).

For the arrangement with cutoff $R=40$, all active walls on
\(1/40<\rho<5/4\) are enumerated by the finite affine rule above.  Define
the finite critical set explicitly by
\[
 \mathcal W_{40}=\{1/40,5/4\}\cup
 \left\{\frac{c-c'}{b'-b}:\ x=c+b\rho,\ x=c'+b'\rho
 \text{ are listed walls, }\frac1{40}<\frac{c-c'}{b'-b}<\frac54\right\}.
\]
At each open interval of this set, choose the least rational midpoint and
generate the $x$-cells and profiles by the endpoint inequalities above;
the profile is constant on the resulting cells.  Their
exact ordering gives 412 \(\rho\)-slabs, 2472 \(x\)-cells and 582 profiles;
the one intentionally deficient profile integrates to the saving \(5/48\).
More explicitly, retain on \(1/40\le\rho\le5/4\) and \(0\le x\le\rho\)
the affine walls
\[
\begin{gathered}
x=0,\quad x=\rho,\quad
x=\tfrac12-k\rho,\quad x=\tfrac32-k\rho,\\
x=\tfrac14-\ell\rho,\quad x=\tfrac52-\ell\rho,\quad
x=w\rho-\tfrac14,\quad x=w\rho .
\end{gathered}
\]
for all nonnegative \(k,\ell\) and positive \(w\) whose line meets the
closed strip.  Only finitely many such lines occur; sort their pairwise
rational intersections and evaluate the exact midpoint profile on each cell.
This finite rule gives the 412 \(\rho\)-slabs, 2472 \(x\)-cells and 582
profiles without floating-point geometry.
The corresponding exact finite comparisons use only the displayed residue
blocks and the stated rank formula.

\subsection{Outer poles and precedence}

For \(p>5n\), use row potentials \(\lambda_B(i),\lambda_V(\ell)\) and
column potentials \(\chi(j)\) as follows:
\[
\begin{array}{c|c|c|c}
 \text{range}&(\lambda_B,\lambda_V,\chi)&\text{cost}&\text{pole}\\ \hline
 5n<p<6n&(1+[i\ge p],1,[j\ge p+2n])&19n-2p&C,B,V\\
 6n\le p<8n&([i+p<10n],1,[j\ge p])&19n-2p&C,B,V\\
 8n\le p<10n&(0,1,[j\ge p])&11n-p&C,B,V\\
 10n\le p\le U&(0,[p-\ell\in J],0)&11n-p&V
\end{array}
\]
Here brackets are indicator functions.  A \(C\)-entry has at most a simple
pole at \(j=i+p\); a \(B\)-entry has the simple harmonic pole and, when the
residue agrees, the second pole (including the diagonal double pole); and a
\(V\)-entry has a simple pole at \(j+\ell=p\) or \(2p\).  The potentials in
the table dominate these orders entry by entry.  In particular, the two
families of \(V\)-poles are explicitly included, and no pole is assigned to
two ranges.  This proves
\[
 v_p(G_m(T;-q))\ge-E_p^{\rm out},\qquad
 E_p^{\rm out}=\begin{cases}19n-2p,&5n<p<8n,\\11n-p,&8n\le p\le U.
\end{cases}
\]
The \(C\)-row potential in the last row is \(0\); its absence there records
that no \(C\)-pole occurs once \(p>10n\).
The argument uses only p-unit cancellation and pole locations; it is unchanged
when \(-2\) is replaced by any fixed \(-q\) in the p-unit branch.

For completeness, fix the finite set \(\mathcal P_q=\{p:p\mid q(q+1)\}\).  Let
\(\mathcal W_{40}\) be the exact set of affine \(\rho\)-walls in this
arrangement, with \(1/40\) and \(5/4\) included as closed endpoints.  Define
the baseline exponent
\[
 e_{\rm nat}(p)=3m\left\lfloor\log_p N\right\rfloor
                 +n\left\lfloor\log_p U\right\rfloor,
\]
which is the coefficientwise exponent supplied by \(\Delta_m=d_N^{3m}d_U^n\),
and assign the exponent to each prime by the following ordered rule:
\[
 E_{p,q}=\begin{cases}
 e_{\rm nat}(p)-v_p(q)D_m-v_p(q+1)m,&p\in\mathcal P_q,\\[2pt]
 e_{\rm nat}(p),&p\notin\mathcal P_q,\ p\le m/40,\\
 e_{\rm nat}(p),&p\notin\mathcal P_q,\ p^2\le\max(N,U),\\
 e_{\rm nat}(p),&p\notin\mathcal P_q,\ p/m\in\mathcal W_{40},\\[2pt]
 m+\sum_u\mathfrak r(u/m,p/m),&
 \substack{p\notin\mathcal P_q,\ 1/40<p/m<5/4,\\ p/m\notin\mathcal W_{40}},\\[2pt]
 E_p^{\rm out},&5n<p\le U,\ p\notin\mathcal P_q,\\[2pt]
 0,&p>U.
 \end{cases}
\]
The first matching line always takes precedence; the cases are therefore an
ordered partition even when a small-prime condition overlaps a wall.
The outer line begins after the inner range.
The coefficientwise bounds above imply
\[
 D^*_{m,q}:=\prod_p p^{E_{p,q}},\qquad
 D^*_{m,q}G_m(T;-q)\in\mathbb Z[T].
\]
The exponent in the first line can be negative; this is allowed because the
formal divisor supplies the corresponding coefficient valuation.  No primitive
content or gcd equality is asserted.

\subsection{Prime summation and the fixed-base rate}

Let \(\vartheta(y)=\sum_{p\le y}\log p\) and
\[
 \eps_{\rm PNT}(X)=\sup_{y\ge X}\left|\frac{\vartheta(y)}y-1\right|.
\]
For every affine profile \(f\) on a rational \(\rho\)-cell,
Stieltjes summation gives
\[
 m\sum_{t_0m<p<t_1m}f(p/m)\log p
 =m^2\int_{t_0}^{t_1}f(\rho)\,d\rho
 +O\!\left(\eps_{\rm PNT}(m/40)m^2+m\log(3m)\right).
\]
The baseline strip contributes \(13/160\), the deliberately credited profile
contributes \(-5/48\), and the outer potentials contribute \(45/32\).  The
exact rational value for this arrangement (including all intervening cells) is
\[
 \begin{aligned}
 A_{40}&=N_{40}/D_{40}=5.082993964478118\ldots,\\
 N_{40}&=132263363803421925450338622847809742317124551497,\\
 D_{40}&=26020759561732371877438650035601516684254668800.
 \end{aligned}
\]
Higher prime powers contribute \(O(m^{3/2}\log(3m))\).  The exact sum over
the fixed exceptional set is
\[
 \sum_{p\in\mathcal P_q}E_{p,q}\log p
 =\sum_{p\in\mathcal P_q}e_{\rm nat}(p)\log p-D_m\log q-m\log(q+1).
\]
Only the first term on the right is the exceptional-prime cost; since
\(\mathcal P_q\) is fixed and \(N,U=O(m)\), it is \(O_q(m\log(3m))\).  The
\((q+1)^m\) formal factor contributes the separate linear term
\(-m\log(q+1)\).  As
\(D_m=53m^2/48-m+O(1)\), the other formal term contributes
\(-\frac{53}{48}m^2\log q+O_q(m)\).  Thus
\[
 \log D^*_{m,q}\le
 \left(A_{40}-\frac{53}{48}\log q\right)m^2+r_{q,m},
\]
where
\[
 r_{q,m}=O\!\left(\eps_{\rm PNT}(m/40)m^2
 +m^{3/2}\log(3m)+m\log(3m)+C_q^{\rm exc}m\log(3m)+B_qm\right),
\]
\(C_q^{\rm exc}=13|\mathcal P_q|/4\), and \(B_q\) is fixed after \(q\) is fixed.  Since
\(\eps_{\rm PNT}(X)\to0\), this is \(o_q(m^2)\).  The quadratic rates are
\[
\begin{array}{c|c}
q&A_{40}-(53/48)\log q\\\hline
3&3.869942895741\\
4&3.552293940742\\
5&3.305906269499\\
6&3.104592883873\\
7&2.934384841563
\end{array}
\]
For \(q=2\), the same formal rate is
\(4.317643952609845\ldots\); its analytic passage is the \(q=2\) instance of
Section~\ref{sec:v-analytic}.

\begin{proposition}[fixed-base arithmetic transfer]
\label{prop:v-fixed-arithmetic}
For each fixed \(q\in\{2,3,4,5,6,7\}\), the determinant has a rational
coefficientwise integerizer satisfying
\[
 \limsup_{m\to\infty,\ 4\mid m}m^{-2}\log D^*_{m,q}
 \le A_{40}-\frac{53}{48}\log q.
\]
For \(q=2,3,4,5,6,7\), the proof is the formal divisor, residue/outer bounds,
wall precedence and PNT argument above.  The finite comparisons instantiate
the displayed formulas; they do not replace the all-$m$ argument.
\end{proposition}

\section{Signed logarithmic energy and smoothing}
\label{sec:v-analytic}

For each fixed $q\in\{2,3,4,5,6,7\}$, directed interval checks on the printed
reference and the smoothing estimate below give an all-$m$ determinant
bound.  The interval checks certify the energy rate; the smoothing lemma
supplies the estimate for every $m=4n\ge64$.

Throughout, $q$ is a fixed integer with $2\le q\le7$, $a=-q$,
$m=4n$.  The three local shape parameters are set to $s=k=h=n$, so that
$L=m+h=5m/4$.  The three masses are
$(M_x,M_v,M_z)=(1,5/4,1)$.  Write
\[
 I(\sigma,\eta)=\iint\log|u-w|\,d\sigma(u)d\eta(w),\qquad
 D_q(\sigma,\eta)=\iint\log|q+uv|\,d\sigma(u)d\eta(v),
\]
and
\begin{equation}\label{eq:v-energy}
 \mathcal Q_q(\mu)=I(\mu_x,\mu_x)+I(\mu_v,\mu_v)+I(\mu_z,\mu_z)
       +I(\mu_v,\mu_z)-D_q(\mu_x,\mu_v).
\end{equation}
The supports used for the continuous problem are
\[
 \mathcal K_x=\mathcal K_v=[0,1],\qquad
 \mathcal K_z(q)=[-q,-r]\cup[r,1]\cup\{|z|=r\},\qquad r=3/5.
\]
For a reference triple $\nu=(\nu_x,\nu_v,\nu_z)$ of these masses, put
\begin{align*}
 U_i(w)&=\int\log|w-t|\,d\nu_i(t),\\
 D_v(w)&=\int\log|q+wt|\,d\nu_v(t),\\
 D_x(w)&=\int\log|q+tw|\,d\nu_x(t).
\end{align*}
and define the three fields
\begin{align}
 \widehat F_x(w)&=2U_x(w)-D_v(w),\notag\\
 \widehat F_v(w)&=2U_v(w)+U_z(w)-D_x(w),\label{eq:v-q-fields}\\
 \widehat F_z(w)&=2U_z(w)+U_v(w),\notag
\end{align}
For the full logarithmic energy we add the one-body terms separately:
\begin{align}
 F_x(w)&=\widehat F_x(w)+\tfrac34\log|w|,\notag\\
 F_v(w)&=\widehat F_v(w)+\tfrac14\log|w|,\label{eq:v-fields}\\
 F_z(w)&=\widehat F_z(w)-\tfrac54\log|w|.\notag
\end{align}
The constant field from the determinant prefactor is
\begin{equation}\label{eq:v-C0}
 C_0(q)=\tfrac32\log q.
\end{equation}

We first identify the determinant at the dilogarithm value. Put
\[
 M_{ij}(q)=\int_0^1\!\int_0^1\frac{x^i y^j}{q+xy}\,dx\,dy,
 \qquad \beta_q=\operatorname{Li}_2(-1/q),
 \qquad \lambda_q=\log(1+1/q).
\]
For $j\ge0$ set the two moment sequences
\[
 m_j=\int_0^1\frac{u^j}{q+u}\,du,
 \qquad n_j=\int_0^1\frac{u^j\log(1/u)}{q+u}\,du.
\]
Polynomial division gives, for $j\ge1$,
\[
 m_j=-q m_{j-1}+\frac1j,\qquad
 n_j=-q n_{j-1}+\frac1{j^2},
\]
while $m_0=\lambda_q$ and $n_0=-\beta_q$. Hence, with the
polynomials $t_j,b_j$ of Section~\ref{sec:v-arithmetic},
\[
 m_j=(-q)^j\lambda_q+t_j,\qquad n_j=b_j-(-q)^j\beta_q.
\]
The substitution $u=xy$ yields
\[
 M_{ij}(q)=\frac{m_i-m_j}{j-i}\quad(i\ne j),
 \qquad M_{ii}(q)=n_i,
\]
and therefore
\begin{equation}\label{eq:v-moment-identity}
 M(q)=B(-q)+\lambda_q C(-q)-\beta_q A(-q).
\end{equation}
Adding $\lambda_q$ times each $C$ row to the corresponding row of
$B(-q)-\beta_qA(-q)$ preserves the determinant. Thus the determinant of
Section~\ref{sec:v-arithmetic}, evaluated at $T=\beta_q$,
is exactly the determinant with middle block $M(q)$.

\subsection{The path and the signed energy}

Let $\gamma_+(t)=re^{it}$ for $0\le t\le\pi$ and
$\gamma_-(t)=re^{it}$ for $\pi\le t\le2\pi$.  Orient the two real
segments increasingly and set
\begin{equation}\label{eq:v-path}
 \Gamma_q=[-q,-r]+[r,1]+(c_q-1)\gamma_++c_q\gamma_-,
 \qquad c_q=\frac{\log q+i\pi}{2\pi i}.
\end{equation}
For every integer $d$,
\begin{equation}\label{eq:v-laurent}
 \int_{\Gamma_q}z^{d-1}\,dz=
 \begin{cases}(1-(-q)^d)/d,&d\ne0,\\0,&d=0.\end{cases}
\end{equation}
For $d\ne0$, integrate $z^{d-1}$ on each piece using the primitive
$z^d/d$; the two circular half-arcs contribute exactly the terms needed
to cancel the intermediate endpoints $\pm r$.  For $d=0$, the real pieces
give $-\log q$, while the circular pieces give
$i\pi(2c_q-1)=\log q$.  This also fixes the branch convention at the
zero Laurent moment.

The same Andreief identity~\cite{Andreief,Forrester} and Cauchy--Vandermonde
expansions~\cite{Krattenthaler} as for the determinant
in Section~\ref{sec:v-arithmetic}, now using \eqref{eq:v-laurent}, give
\begin{equation}\label{eq:v-integral}
 G_m(\operatorname{Li}_2(-1/q);-q)
 =\frac{q^{3m^2/2-m}}{(m!)^2L!}
   \int_{[0,1]^m}\int_{[0,1]^L}\int_{\Gamma_q^m}
       \Phi_q(x,v,z)\,dz\,dv\,dx,
\end{equation}
where
\[
 \Phi_q=\frac{\Delta(x)^2\Delta(v)^2\Delta(z)^2
       \prod_{d=1}^m\prod_{b=1}^L(z_d-v_b)
       \prod_{i=1}^m x_i^{3m/4}
       \prod_{b=1}^L v_b^{m/4}
       \prod_{d=1}^m z_d^{-5m/4}}
       {\prod_{i=1}^m\prod_{b=1}^L(q+x_iv_b)}.
\]

\begin{proof}[Proof of the integral representation]
By \eqref{eq:v-laurent},
\(C_{ij}=(-q)^i\int_{\Gamma_q}z^{j-i-1}\,dz\).  Put
\(f_i(v)=\int_0^1x^i/(q+xv)\,dx\).  Then the lower rational block and the
V block are respectively
\[
 M_{ij}=\int_0^1f_i(v)v^j\,dv,\qquad
 V_{\ell j}=\int_0^1v^{\ell-1}v^j\,dv.
\]
For integrable families, Andreief's identity~\cite{Andreief,Forrester} is
\[
 \det\left(\int f_i(t)g_j(t)\,dt\right)_{i,j=1}^N
 =\frac1{N!}\int
   \det(f_i(t_b))_{i,b=1}^N\det(g_j(t_b))_{b,j=1}^N
   \prod_{b=1}^Ndt_b.
\]
Apply it successively to the $L$ lower/V rows, the $m$ contour rows, and
the $m$ functions $f_i$.  These three applications account respectively
for $1/L!$, $1/m!$, and $1/m!$:
\[
 G_m=\frac1{(m!)^2L!}\int D_zD_{zv}D_xD_{xv}\,dx\,dv\,dz.
\]
Here $e_\Delta=m(m-1)/2$ and, with
\(\mathbf w=(z_1,\ldots,z_m,v_1,\ldots,v_L)\), the four determinants are
\begin{align*}
 D_z&=\det\bigl((-q)^iz_d^{-i-1}\bigr)_{i\in I,d}
   =(-1)^{e_\Delta}(-q)^{\sum_{i\in I}i}
     \Delta(z)\prod_dz_d^{-3m/2},\\
 D_{zv}&=\det(\mathbf w_b^j)_{b,j\in J}
   =(-1)^{mL}\Delta(z)\Delta(v)
     \prod_dz_d^{m/4}\prod_bv_b^{m/4}\prod_{d,b}(z_d-v_b),\\
 D_x&=\det(x_{i'}^i)_{i\in I,i'}
   =\Delta(x)\prod_{i'}x_{i'}^{m/2},\\
 D_{xv}&=\det\begin{pmatrix}
 ((q+x_{i'}v_b)^{-1})_{i',b}\\(v_b^{\ell-1})_{\ell,b}
 \end{pmatrix}
   =\frac{(-1)^{e_\Delta}q^{e_\Delta}\Delta(x)\Delta(v)\prod_{i'}x_{i'}^{m/4}}
          {\prod_{i',b}(q+x_{i'}v_b)}.
\end{align*}
The last equality is the Cauchy--Vandermonde identity~\cite{Krattenthaler}.
For completeness,
for $t_{i'}=q/x_{i'}$ it follows from
\[
 \det\begin{pmatrix}((t_{i'}+v_b)^{-1})_{i',b}\\
 (v_b^{\ell-1})_{\ell,b}\end{pmatrix}
 =\frac{\Delta(t)\Delta(v)}{\prod_{i',b}(t_{i'}+v_b)},
\]
after multiplying the first $m$ rows by $x_{i'}^{-1}$ and using
\(\Delta(t)=(-1)^{e_\Delta}q^{e_\Delta}\Delta(x)\prod_{i'}x_{i'}^{1-m}\).
Here is a proof of the identity, for every $m,h\ge0$ with $L=m+h$.
Treat the $t_i$ and $v_b$ as independent indeterminates.  The case $m=0$
is the Vandermonde determinant.  Assume the formula for $m-1$ and all
$h\ge0$, and regard both sides as rational functions of $v_L$.  Expanding
the determinant along its last column, its residue at $v_L=-t_i$ is
$(-1)^{i+L}$ times the complementary minor.  By induction this residue is
\[
 (-1)^{i+L}
 \frac{\Delta(\widehat t_i)\Delta(v_1,\ldots,v_{L-1})}
      {\prod_{j\ne i}\prod_{b<L}(t_j+v_b)}.
\]
The proposed quotient has the same residue: substitute $v_L=-t_i$,
cancel the factors $t_j-t_i$ and $t_i+v_b$, and use
$\Delta(t)=(-1)^{i-1}\Delta(\widehat t_i)
\prod_{j\ne i}(t_j-t_i)$.  Hence the difference has no finite poles.
After cancellation of the denominator $\prod_i(t_i+v_L)$, the difference
is a polynomial in $v_L$ of degree at most $h-1$.  It vanishes at each of
the $L-1$ values $v_L=v_b$ ($b<L$), since both sides have the factor
$\Delta(v)$.  For $m\ge1$, $L-1=m+h-1\ge h$, so these roots force the
polynomial to be zero.  When $h=0$, the difference tends to zero at
infinity instead.  This proves the Cauchy--Vandermonde identity.

Multiplying the four displayed factors gives precisely $\Phi_q$: the
$m/2$ from $D_x$ and the $m/4$ from $D_{xv}$ combine to the displayed
$3m/4$ power of each $x_{i'}$.  The sign is
$(-1)^{\sum_{i\in I}i+2e_\Delta+mL}=1$: here
$\sum_{i\in I}i=2n(8n-1)$ and $mL=20n^2$ are even.  The q-power is
\[
 \sum_{i\in I}i+e_\Delta
 =m(2n)+m(m-1)
 =\frac{3m^2}{2}-m,
 \qquad (m=4n,\ L=5n).
\]
The factors $q+x_iv_b$ are at least
$q$, while $|z|\ge r$ on $\Gamma_q$; all three integrations therefore
converge absolutely.  The $m!$, $L!$, and $m!$ factors are exactly the three
Andreief symmetrizations, proving \eqref{eq:v-integral}.
\end{proof}
For triples of positive measures of the stated masses, define
\begin{equation}\label{eq:v-full-energy}
 \mathcal E_q(\mu)=C_0(q)+\mathcal Q_q(\mu)
 +\tfrac34\int\log|x|\,d\mu_x
 +\tfrac14\int\log|v|\,d\mu_v
 -\tfrac54\int\log|z|\,d\mu_z.
\end{equation}

\begin{lemma}[signed-energy concavity for fixed $q$]
\label{lem:v-concavity}
Let $\mu,\nu$ be triples of positive measures of masses
$(1,5/4,1)$, supported in compact subsets of the indicated supports or
their sufficiently small tubular neighborhoods.  Assume that the pairings
in $\mathcal Q_q$ are finite real numbers.  No absolute logarithmic moment
assumption is imposed on a competitor's $x$- or $v$-component.  Then
\[
 \mathcal Q_q(\mu)\le-\mathcal Q_q(\nu)
       +\int \widehat F_x\,d\mu_x+\int \widehat F_v\,d\mu_v
       +\int \widehat F_z\,d\mu_z,
\]
where the hatted fields are formed from $\nu$.  Consequently the same
inequality for the full energy is obtained by adding the displayed one-body
terms to the hatted fields, i.e. by using $F_x,F_v,F_z$ in the field bounds.
\end{lemma}

\begin{proof}
By the shared quadratic-duality identity in
Section~\ref{sec:certificate-algorithm}, it suffices to prove that the
zero-mass quadratic remainder $\mathcal Q_q(\mu-\nu)$ is nonpositive.
First suppose that all $v$-supports avoid zero.  For a signed measure
$\sigma$ of mass zero, write $\|\sigma\|^2=-I(\sigma,\sigma)\ge0$;
this is the zero-mass logarithmic-energy positivity lemma.  Put
$R_q(v)=-q/v$.  The exact identities
\begin{align*}
 \log|q+xv|&=\log|x-R_q(v)|+\log|v|,\\
 \log|R_q(v)-R_q(w)|&=\log q+\log|v-w|-\log|v|-\log|w|
\end{align*}
Integrating against zero-mass measures cancels all one-variable terms and
gives
\begin{align*}
 D_q(\sigma_x,\sigma_v)&=I(\sigma_x,(R_q)_*\sigma_v),\\
 I((R_q)_*\sigma_v,(R_q)_*\sigma_v)&=I(\sigma_v,\sigma_v).
\end{align*}
Let $\sigma_i=\mu_i-\nu_i$, and set
$\sigma'_x=\sigma_x$, $\sigma'_v=(R_q)_*\sigma_v$,
$\sigma'_z=-(R_q)_*\sigma_z$.  Then
\[
 \mathcal Q_q(\sigma)
 =-\|\sigma'_x\|^2-\|\sigma'_v\|^2-\|\sigma'_z\|^2
   +\langle\sigma'_x,\sigma'_v\rangle
   +\langle\sigma'_v,\sigma'_z\rangle,
 \qquad \langle\alpha,\beta\rangle=-I(\alpha,\beta).
\]
The right side equals
\[
 -\tfrac12\bigl(\|\sigma'_x-\sigma'_v\|^2
              +\|\sigma'_v-\sigma'_z\|^2
              +\|\sigma'_x\|^2+\|\sigma'_z\|^2\bigr)\le0.
\]
The shared identity gives the asserted quadratic inequality.  For a
competitor $\mu_v$ whose support meets zero, fix a smooth probability
measure $\rho$ supported in $[3/4,7/8]$.  For $0<\delta<1/4$, let
$\lambda_\delta=\mu_v|_{\{|v|<\delta\}}$, let
$e_\delta=\lambda_\delta(\mathbb C)$, and set
\[
 \mu_{v,\delta}=\mu_v-\lambda_\delta+e_\delta\rho.
\]
Finite logarithmic self-energy excludes an atom at zero, so
$e_\delta\to0$ and $\mu_{v,\delta}\Rightarrow\mu_v$.  Absolute
integrability of the logarithmic pairings in $\mathcal Q_q$ implies that
the integrals over the deleted sets tend to zero.  The replacement terms
involving $e_\delta\rho$ tend to zero as well: $\rho$ is separated from
$\{|v|<\delta\}$, and its logarithmic potential is bounded on the compact
supports.  The kernel $\log|q+uv|$ is bounded and continuous on the stated
tubes because $q\ge2$ and $|u|,|v|\le1+o(1)$.  Hence
$\mathcal Q_q(\mu_{x},\mu_{v,\delta},\mu_z)\to\mathcal Q_q(\mu)$.
The reference fields $\widehat F_i$ are continuous and bounded on these
compact sets, so their integrals also converge.  Applying the already
proved inequality to $\mu_{v,\delta}$ and taking the limit proves the
quadratic inequality without requiring $\int|\log|v||\,d\mu_v<\infty$.

The one-body integrals in $\mathcal E_q$ are then interpreted as extended
real integrals.  On the compact $x$- and $v$-supports, $\log|w|$ is bounded
above and can have integral $-\infty$; on the $z$-support, $|z|\ge r-o(1)$,
so the $z$-term is finite.  Adding these linear terms to the quadratic
inequality is therefore valid in the extended reals.  In particular, when
the full energy is $-\infty$ the asserted upper bound is automatic.  This
argument uses no implication from finite self-energy to a finite absolute
logarithmic moment.
\end{proof}

Each printed reference has finitely many positive-length rational interval
rows, compact $z$ support in $[-q,-3/5]$, and Poisson centres with
$|c|<3/5$; the $x$ and $v$ chains start at the positive endpoints listed in
the data section and end at $1$.  Hence all six reference supports are
compact and avoid zero.  Interval densities are bounded by the finite
maximum of their rational mass/length ratios, while a Poisson row of mass
$w$ has density
\[
 \frac{w}{2\pi}\frac{r^2-|c|^2}{|re^{i\theta}-c|^2}
 \le\frac{w(r+|c|)}{2\pi r(r-|c|)},\qquad r=3/5.
\]
Thus the printed references have finite logarithmic moments directly from
their rational supports, without any general self-energy implication.

Combining the lemma with the one-body terms gives the explicit full-energy
upper bound used below:
\begin{equation}\label{eq:v-full-energy-dual}
 \mathcal E_q(\mu)
 \le C_0(q)-\mathcal Q_q(\nu)
      +\int F_x\,d\mu_x+\int F_v\,d\mu_v+\int F_z\,d\mu_z.
\end{equation}
Thus, if the full fields are bounded above by $K_{x,q},K_{v,q},K_{z,q}$
on their respective admissible supports, then
\begin{equation}\label{eq:v-energy-bound}
 \mathcal E_q(\mu)\le -\bigl(\mathcal Q_q(\nu)-C_0(q)-K_{x,q}
                         -\tfrac54K_{v,q}-K_{z,q}\bigr).
\end{equation}

The inversion is valid uniformly for the fixed bases considered here:
on $|v|\le1+3/4096$,
$|R_q'(v)|=q/|v|^2>1$ for $q\ge2$.  Thus a bounded arclength density
pushes forward to a bounded arclength density on each inverted segment or
circle.  No $q=2$ numerical identity is used in the concavity argument.

\subsection{Directed fixed-base certificates}

For a reference triple $\nu_q$, suppose its energy lower endpoint is
$Q_{q,-}$ and its support-wide field upper bounds are
$F_x\le K_{x,q}$, $F_v\le K_{v,q}$, $F_z\le K_{z,q}$.  Lemma
\ref{lem:v-concavity} and \eqref{eq:v-full-energy} give the quadratic
rate
\begin{equation}\label{eq:v-alpha}
 \alpha_q=Q_{q,-}-\tfrac32\log q-K_{x,q}
                         -\tfrac54K_{v,q}-K_{z,q}.
\end{equation}
The following are the directed interval-bound results.  The field thresholds
are exact terminating decimals.  The displayed $\alpha$ values are strict
lower bounds rounded down, and each $\tau$ is an outward-rounded upper bound.

\begin{table}[ht]
\centering\normalsize
\renewcommand{\arraystretch}{1.12}
\caption{Fixed-base directed certificates for $a=-q$.}
\label{tab:v-certificates}
\setlength{\tabcolsep}{2pt}
\begin{tabular}{@{}c c c c c c c c@{}}
\hline
$q$&rows&$K_x$&$K_v$&$K_z$&$\alpha_q$&$\tau_q$&margin\\
\hline
2&$41/48/62/149$&$-4.621$&$-4.475$&$-0.085$&$>4.34875$&$<4.31765$\footnotemark&$>0.03110$\\
3&$40/48/62/143$&$-5.076$&$-4.632$&$0.388$&$>3.96707$&$<3.86995$&$>0.09712$\\
4&$40/48/62/141$&$-5.408$&$-4.727$&$0.779$&$>3.68211$&$<3.55230$&$>0.12981$\\
5&$39/48/62/137$&$-5.670$&$-4.791$&$1.108002$&$>3.45240$&$<3.30591$&$>0.14649$\\
6&$41/47/61/135$&$-5.886$&$-4.839$&$1.391$&$>3.26092$&$<3.10460$&$>0.15632$\\
7&$39/47/62/131$&$-6.071$&$-4.876$&$1.639$&$>3.09651$&$<2.93439$&$>0.16212$\\
\hline
\end{tabular}
\end{table}
\footnotetext{The exact rate is
\(\tau_2=A_{40}-(53/48)\log2=4.317643952609845\ldots\); the displayed
decimal is a directed upper bound.}

The energy lower endpoints used in \eqref{eq:v-alpha} are printed separately
so that the energy and field slacks can be checked independently:
\begin{table}[ht]
\centering
\renewcommand{\arraystretch}{1.12}
\caption{Directed lower endpoints for the reference energies.}
\label{tab:v-energy-lower}
\begin{tabular}{@{}c c@{}}
\hline
$q$ & $Q_{q,-}$\\
\hline
2&$>-4.911272$\\
3&$>-4.863007$\\
4&$>-4.776198$\\
5&$>-4.684191$\\
6&$>-4.595191$\\
7&$>-4.511619$\\
\hline
\end{tabular}
\end{table}

For each $q=2,\ldots,7$, put $\delta=1/250$.  The boxes
$0\le x,v\le\delta$ are discharged by the endpoint bounds
\eqref{eq:v-endpoint}.  On $[\delta,1]$ and on the two real $z$ intervals,
including the complete negative interval $[-q,-3/5]$, the subdivision
recursively bisects each rational midpoint.  A cell is accepted exactly when
its directed upper field interval is finite and strictly below the displayed
threshold; otherwise both children are visited, in left-first
order, with maximum depth $30$.  Circle cells use the same rule in the
parameter $t\in[0,1]$, and the endpoints $t=0,1$ are checked directly.
This is the complete partition rule.  The six references contain $1733$ rows
in all, and their rows are printed in Section~\ref{sec:v-reference-data}.
The smoothing lemma below supplies the complex-tube transfer and the all-$m$
estimate.

\subsection{Reference densities and directed endpoint disks}
\label{subsec:v-allm-data}

All constants below are computed from the signed decimal tables in
Section~\ref{sec:v-reference-data}.  For
each interval row $[a,b]$ of mass $w$, put $H=w/(b-a)$.  Let $H_{x,q}$,
$H_{v,q}$ and $H_{z,q}$ be the largest such value in the $x$, $v$ and
negative-real $z$ chains.  For a Poisson row $(c,w)$ write
\[
 \rho_c=2^{-120}\left\lceil2^{120}|c|\right\rceil,
 \qquad
 P_q=\sum_{(c,w)}\frac{w(r+\rho_c)}{r(r-\rho_c)}.
\]
The ceiling is evaluated by integer square root on the printed rational
coordinates, and the construction checks $\rho_c<r$ row by row.  Since the
Poisson density on $|z|=r$ is at most
$w(r+\rho_c)/(2\pi r(r-\rho_c))$ per unit arclength, and a circle meets a
disk of radius $t$ in arclength at most $2\pi t$, the reference measures
satisfy $\nu_i(B(w,t))\le A_{i,q}t$, where
\[
 A_{x,q}=2H_{x,q},\qquad A_{v,q}=2H_{v,q},\qquad
 A_{z,q}=2H_{z,q}+P_q.
\]
Define
\[
 B_{i,q}=6A_{i,q}+2M_i,\qquad
 \Sigma_q=2B_{x,q}+\frac72B_{v,q}+\frac{13}{4}B_{z,q},
 \qquad \mathcal C_q=10+3\Sigma_q.
\]
In the table, $\Sigma_q$ is rounded upward to $10^{-3}$ and this rounded
rational value is used in $\mathcal C_q$.  The path variation is
\[
 \mathcal T_q=q+1-2r+r\sqrt{\pi^2+\log^2q}.
\]
The $E_{x,q}$ and $E_{v,q}$ columns are rounded upward to $10^{-4}$;
the $\mathcal T_q$ column is rounded upward to $10^{-6}$.

For $\delta=1/250$, put $P(t)=t(\log t-1)$ and define the exact real
integrals
\[
 J_{x,q}=\int\log(t+\delta)\,d\nu_x(t),\quad
 J_{v,q}=\int\log(t+\delta)\,d\nu_v(t),\quad
 J_{z,q}=\int\log(|z|+\delta)\,d\nu_z(z).
\]
An interval contributes $w(P(b+\delta)-P(a+\delta))/(b-a)$; for a
negative $z$ interval use endpoints $-a,-b$ in increasing order.  The
Poisson rows together contribute their total mass times $\log(r+\delta)$.
For $|w_0|\le\delta$, the triangle inequality and the masses give
\begin{align}
 F_x(w_0)&\le E_{x,q}:=2J_{x,q}-\frac54\log(q-\delta)
                      +\frac34\log\delta,\notag\\
 F_v(w_0)&\le E_{v,q}:=2J_{v,q}+J_{z,q}-\log(q-\delta)
                      +\frac14\log\delta.\label{eq:v-endpoint}
\end{align}
The endpoint bounds are directed upper endpoints rounded upward at
$10^{-4}$.  The same table gives $\Sigma_q$ rounded upward at $10^{-3}$ and
$\mathcal T_q$ rounded upward at $10^{-6}$; throughout the proof set
$\mathcal C_q=10+3\Sigma_q$ using the printed $\Sigma_q$.
\begin{table}[ht]
\centering\normalsize
\renewcommand{\arraystretch}{1.12}
\caption{Endpoint and all-$m$ transfer data recomputed from the printed tables.}
\label{tab:v-transfer-constants}
\begin{tabular}{@{}crrrr@{}}
\hline
$q$&$E_{x,q}$&$E_{v,q}$&$\Sigma_q$ (upper)&$\mathcal T_q$ (upper)\\
\hline
2&$-6.2821$&$-4.4935$&$14092.327$&$3.730292$\\
3&$-6.7656$&$-4.6587$&$13219.499$&$4.796889$\\
4&$-7.1126$&$-4.7555$&$12702.122$&$5.860319$\\
5&$-7.3835$&$-4.8208$&$12351.983$&$6.917916$\\
6&$-7.6061$&$-4.8688$&$12095.881$&$7.969979$\\
7&$-7.7947$&$-4.9057$&$11898.583$&$9.017257$\\
\hline
\end{tabular}
\end{table}

The endpoint entries satisfy $E_{x,q}<K_{x,q}$ and $E_{v,q}<K_{v,q}$,
where the field bounds are in Table~\ref{tab:v-certificates}.  The displayed
$\Sigma_q$ values
give the explicit coefficients $\mathcal C_q=10+3\Sigma_q$; no unprinted density or
endpoint constant is used below.

\subsection{The all-$m$ smoothing estimate}
\label{subsec:v-allm-proof}

\begin{lemma}[logarithmic-potential modulus]
\label{lem:v-potential-modulus}
If a compactly supported positive measure $\eta$ of mass $M$ satisfies
$\eta(B(w,t))\le At$ for every disk, then for $d=|w-w'|\le1/3$,
\[
 |U_\eta(w)-U_\eta(w')|
 \le(6A+2M)d\bigl(1+\log(1/d)\bigr).
\]
\end{lemma}
\begin{proof}
Let $N=\{z:|w-z|\le2d\}$.  On $N$ both distances are at most $3d$.
Approximating the closed disk from outside gives $\eta(N)\le2Ad$.  For
either centre $w_*=w,w'$, layer cake bounds
the integral on $N$ by
\[
 \int_N\log\frac1{|w_*-z|}\,d\eta(z)
 \le\int_0^\infty\min(2Ad,Ae^{-s})\,ds
 =2Ad\bigl(1+\log(1/(2d))\bigr).
\]
Off $N$, the logarithmic difference is at most $2d/|w-z|$, while another
layer-cake integration gives
$\int_{|w-z|>2d}|w-z|^{-1}d\eta(z)\le M+A\log(1/(2d))$.
Adding these bounds gives at most
$6Ad\log(1/d)+2Md$, which is bounded by the displayed expression.
\end{proof}

For $q=2,\ldots,7$, let $\alpha_q$ be the strict lower endpoints in
Table~\ref{tab:v-certificates}.  In particular, the printed table verifies
$\alpha_q>\tau_q$ in each of the six cases.  Then for every
$m=4n\ge64$,
\begin{equation}\label{eq:v-allm-bound}
 \log|G_m(\operatorname{Li}_2(-1/q);-q)|
 \le-\alpha_qm^2+\mathcal C_qm\log m,
 \qquad \mathcal C_q=10+3\Sigma_q.
\end{equation}
\begin{proof}
Fix $q$ and a nonzero integrand point in \eqref{eq:v-integral}.  Put
$\epsilon=m^{-2}$.  Replace each atom $u$ of its empirical measures by
the uniform probability measure on
$u+2i\epsilon+\epsilon e^{i\vartheta}$.  Every point of this circle is
within $3\epsilon$ of $u$.  The common translation preserves pairwise
centres; subharmonicity of $\log|z|$ shows that every off-diagonal positive
logarithmic pair energy does not decrease.  Each diagonal circle has
self-energy $\log\epsilon$, and there are $m+5m/4+m=13m/4$ atoms.  Thus,
writing $\mu^\epsilon$ for the smoothed empirical triple and
$\mathcal Q_{q,\rm disc}$ for the three self-energies with their diagonals
removed,
\[
 \mathcal Q_{q,\rm disc}(\mu)
 \le\mathcal Q_q(\mu^\epsilon)
       +\frac{13}{4m}\log(1/\epsilon)+C_{D,q}\epsilon,
 \quad C_{D,q}=8L_{D,q},
 \quad L_{D,q}=\frac{1+3/4096}{q-(1+3/4096)^2}.
\]
Indeed, on the two smearing disks $|u|,|v|\le1+3\epsilon$ and
$|uv-xv_0|\le6\epsilon+9\epsilon^2$.  The denominator stays at distance
at least $q-(1+3/4096)^2$ from zero.  The mean-value bound for
$\log|q+uv|$, multiplied by the cross mass $5/4$, is at most
\[
\left(\frac{15}{2}+\frac{45}{4}\epsilon\right)L_{D,q}\epsilon
\le C_{D,q}\epsilon,
\]
since $\epsilon\le1/4096$.

The smoothed measures are tame and their $v$-support avoids zero, so
Lemma~\ref{lem:v-concavity} applies.  The modulus lemma and the constants
in Table~\ref{tab:v-transfer-constants} bound the changes in all reference
potentials by
\[
 \Sigma_q(3\epsilon)\bigl(1+\log(1/(3\epsilon))\bigr).
\]
For the two denominator potentials in the dual fields, each derivative in
the field variable is at most $L_{D,q}$ on the tube.  Their total mass
coefficient is $5/2$, so their transfer error is at most
$C_{D,q}\epsilon$.

For the positive one-body terms, if a real centre $t\ge\delta/2$, the
circle mean of $\log|t+2i\epsilon+\epsilon e^{i\vartheta}|$ is
$\log|t+2i\epsilon|\le\log t+8\epsilon^2/\delta^2$.  The combined
weighted loss for $x$ and $v$ is therefore at most
$(17/2)\epsilon^2/\delta^2$.  If $t<\delta/2$, the whole smearing circle
lies in $|w|\le\delta$, so \eqref{eq:v-endpoint} gives the endpoint bound directly.
For the negative $z$ one-body term, $|z|\ge r$,
$|z^\epsilon-z|\le3\epsilon$, and $3\epsilon<r$ give
\[
 \left|\log|z^\epsilon|-\log|z|\right|
 \le\frac{3\epsilon}{r-3\epsilon}.
\]
Consequently,
\[
 -\frac54\log|z^\epsilon|
 \le-\frac54\log|z|+\frac{15\epsilon}{4(r-3\epsilon)}.
\]
The centres lie on the certified supports; combining the last three
estimates with the endpoint bounds and the directed field bounds yields
\begin{align*}
 \frac{\log|\Phi_q|}{m^2}
 \le{}&-\mathcal Q_q(\nu_q)+K_{x,q}+\frac54K_{v,q}+K_{z,q}
 +\frac{13}{4m}\log(1/\epsilon)\\
 &+3\Sigma_q\epsilon\bigl(1+\log(1/(3\epsilon))\bigr)
 +\left(2C_{D,q}+\frac{15}{4(r-3/4096)}\right)\epsilon
 +\frac{17}{2}\frac{\epsilon^2}{\delta^2}.
\end{align*}
Here the coefficient $2C_{D,q}$ accounts separately for the denominator
in the smoothed energy and in the transferred dual fields.  For all
$q\ge2$, the rational comparison
$100(1+3/4096)<101(2-(1+3/4096)^2)$ gives
$L_{D,q}<101/100$, hence $C_{D,q}<8.08$.  Also
$15/(4(r-3/4096))<7$, and
\[
 \frac{17}{2\delta^2m^2}\le\frac{531250}{4096}<130
 \qquad(m\ge64).
\]
The total constant after multiplying by $m^2$ is therefore less than
$16.2+7+130<200$.

The path length is exactly $\mathcal T_q$.  Since $q\le7$ and $e>8/3$, we have
$\log q<2$.  Also $\pi<22/7$, so
$\sqrt{\pi^2+\log^2q}<4$, whence
$\mathcal T_q/q<21/10<e$.  The prefactor in \eqref{eq:v-integral} contributes
$(3m^2/2-m)\log q$, and the path contributes $m\log\mathcal T_q$; their linear
part is $m\log(\mathcal T_q/q)\le m$.  Add these terms and drop the favorable
$-\log((m!)^2L!)$.  With $\epsilon=m^{-2}$, the diagonal term is
$(13/2)m\log m$.  Also
\[
 3\Sigma_q\bigl(1+\log(m^2/3)\bigr)\le3\Sigma_qm\log m,
 \quad 200\le m\log m,
 \quad m\le\tfrac14m\log m
\]
for $m\ge64$.  The resulting coefficient is at most
$13/2+3\Sigma_q+1+1/4<10+3\Sigma_q$.  Finally the definition of $\alpha_q$ in
\eqref{eq:v-alpha} gives \eqref{eq:v-allm-bound}.  If the determinant
integrand vanishes, the bound holds trivially; otherwise all empirical
measures above are defined and the argument applies.  This proves the
claim for every $m=4n\ge64$, not just for the finite certificate values.
\end{proof}

Thus the directed margins and the all-$m$ analytic estimate are separate
inputs: the former supplies $\alpha_q>\tau_q$, while
\eqref{eq:v-allm-bound} supplies the full quantifier with an explicit
finite coefficient for each fixed $q=2,\ldots,7$.

\section{Seed determinants and eventual nonvanishing}\label{sec:nonvanishing}

This section proves the nonvanishing input for the determinants
used in each of the two constructions.  The P0 family has a base-dependent
auxiliary prime; the signed construction uses one explicitly listed prime for
each base $a=-r$, $2\le r\le7$.  The prime-power variable below is always
separate from the base parameter.  Finite seed calculations are printed, but
nonvanishing at all sufficiently high powers follows from the residue-block
identities proved here, not from extrapolating a finite table.

\subsection{The P0 seed and its base-dependent prime}

Write $G^{(0)}_{m,a}(T)$ for the P0 determinant of
\eqref{eq:p0det}.  At $m=4$ its row and column sets are
$I_4=\{1,2,3,4\}$ and $J_4=\{0,1,\ldots,7\}$.  Exact elimination in this
the $8\times8$ determinant gives
\begin{equation}\label{eq:p0-seed-factor}
 G^{(0)}_{4,a}(0)=\frac{a^9(a-1)^2H(a)}{D},\qquad
 D=2090188800000=2^{17}3^6 5^5 7.
\end{equation}
For an exact coefficient record, write $H(a)=\sum_{j=0}^{17}h_ja^j$.
In descending degree the vector $(h_{17},\ldots,h_0)$ is
\[
\left(\begin{array}{r}
10369949184000000\\
-79502943744000000\\
273672012096000000\\
-559106235072000000\\
754609551358320000\\
-707442931414800000\\
469847663340300000\\
-220367388160080000\\
71443538332097400\\
-15753667204602600\\
2634445947232050\\
-486626431113900\\
91676122191935\\
-9473737791290\\
424034049270\\
4759882606\\
-44546479\\
1296
\end{array}
\right).
\]
Reduction of this vector gives
\begin{equation}\label{eq:p0-seed-mod5}
 H(a)\equiv a^2+a+1\pmod 5.
\end{equation}
The values of the right side at $a=0,1,2,3,4$ are respectively
$1,3,2,3,1$.  Thus $H(a)\ne0$ for every integer $a$.

Fix any integer $a$ with $|a|\ge2$, and put
\[
 N_a=10!\,|a(a-1)H(a)|+1.
\]
This integer is greater than one.  Choose any prime divisor $p$ of $N_a$.
Since $N_a\equiv1\pmod\ell$ for every prime $\ell\le10$, we have $p>10$.
Also $p$ divides none of $a$, $a-1$, and $H(a)$, and $p\nmid D$ by the
factorization in \eqref{eq:p0-seed-factor}.  Consequently
\begin{equation}\label{eq:p0-seed-unit}
 a,\quad a-1,\quad G^{(0)}_{4,a}(0)\quad\text{are all units in }\mathbb Z_{(p)}.
\end{equation}
The choice is allowed to depend on $a$; no fixed prime is asserted to work
for every base.

\subsection{P0 Frobenius blocks}

Let $Q=p^k$ and $m=4Q$.  Write each index uniquely as $i=Qu+r$ with
$0\le r<Q$.  The P0 row indices are $i=Q,\ldots,5Q-1$, and its column
indices are $j=0,\ldots,8Q-1$.  The following congruence is coefficientwise
in $\mathbb Z_{(p)}[X]$:
\begin{equation}\label{eq:p0-frobenius}
 Q^{12Q}G^{(0)}_{4Q,a}(X/Q^2)\equiv
 G^{(0)}_{4,a}(X)^Q\pmod p.
\end{equation}

We give the entry argument, including the denominator and permutation
checks.  In the P0 determinant, all summation denominators, index
 differences, and diagonal denominators occurring at level $4Q$ are nonzero
 integers whose absolute values are less than $10Q$.  As $p>10$, each has $p$-adic valuation at most
$k$; the only other denominators are powers of $a-1$, which are $p$-units by
\eqref{eq:p0-seed-unit}.  Therefore multiplying each $C$ row by $Q$ and
each $B-TA$ row by $Q^2$ produces a matrix over $\mathbb Z_{(p)}[X]$.
Its determinant is exactly the left side of \eqref{eq:p0-frobenius}.

Lucas' congruence and $a^Q\equiv a$ in $\mathbb F_p$ give
\[
 \binom{Q(u+2)+r}{2Q}a^{Qu+r}
 \equiv a^r\binom{u+2}{2}a^u\pmod p.
\]
In the finite sums defining the P0 rows, multiplication by $Q$ (or $Q^2$
on a second-order denominator) kills every summand whose denominator is not
divisible by $Q$.  The surviving summands, after division of their indices by
$Q$, are exactly the corresponding level-$4$ sums.  For the correction terms
this uses $(a/(a-1))^Q\equiv a/(a-1)\pmod p$.  Thus, with $u,v$ denoting
the quotient indices and $r,s$ the residues,
\begin{align}
 Q C_{Qu+r,Qv+s}&\equiv \mathbf 1_{r=s}a^r C_{u,v}\pmod p,\\
 Q^2 B_{Qu+r,Qv+s}&\equiv \mathbf 1_{r=s}a^r B_{u,v}\pmod p,\\
 A_{Qu+r,Qv+s}&\equiv \mathbf 1_{r=s}a^r A_{u,v}\pmod p.
 \label{eq:p0-entry-blocks}
\end{align}
For the off-diagonal entries with unequal residues, the first two congruences
also follow because $v_p(j-i)<k$, so $Q/(j-i)$ is zero modulo $p$ after
localization.  When the residues agree, $j-i=Q(v-u)$; the reduced nonzero
differences satisfy $|v-u|\le7<p$.  The diagonal identities follow by the
same denominator filtering in the displayed harmonic sums.  This proves
\eqref{eq:p0-entry-blocks} for every entry, including the diagonals.

After grouping rows and columns by residues, the scaled matrix is block
diagonal modulo $p$, with $Q$ blocks
\[
 a^r\begin{bmatrix}C_4(a)\\ B_4(a)-XA_4(a)\end{bmatrix},
 \qquad 0\le r<Q.
\]
There are eight rows in each block, so their scalar product is
$a^{8\sum_{r=0}^{Q-1}r}=a^{4Q(Q-1)}$.  The column regrouping sign is
$(-1)^{\binom Q2\binom82}=1$, and interleaving the two four-row families
has sign $(-1)^{4^2\binom Q2}=1$.  Since $a^{Q-1}\equiv1\pmod p$, the
scalar product is also $1$ modulo $p$.  Taking determinants proves
\eqref{eq:p0-frobenius}.

\begin{lemma}[P0 eventual nonvanishing]\label{lem:p0-eventual-nonzero}
For every fixed integer $a$ with $|a|\ge2$ and every fixed rational number
$\xi$, there is $k_0=k_0(a,\xi)$ such that
\[
 G^{(0)}_{4p^k,a}(\xi)\ne0\qquad(k\ge k_0),
\]
where $p$ is any prime selected as above.
\end{lemma}
\begin{proof}
Set $W=Q^2\xi$ in the scaled determinant of
\eqref{eq:p0-frobenius}.  For all sufficiently large $k$,
$v_p(W)=2k+v_p(\xi)>0$ (with the assertion immediate for $\xi=0$).  Its
left-side scaled matrix is over $\mathbb Z_{(p)}$, while its determinant
reduces modulo $p$ to
\[
 G^{(0)}_{4,a}(0)^Q,
\]
which is a unit by \eqref{eq:p0-seed-unit}.  The scaled determinant is
therefore nonzero, and its scaling factor is nonzero; hence
$G^{(0)}_{4Q,a}(\xi)\ne0$.
\end{proof}

\subsection{Negative V seeds}

For $a=-r$, $2\le r\le7$, write $G^{(V)}_{4n,a}(T)$ for the
column-$V$ determinant
\[
 G^{(V)}_{4n,a}(T)=\det\begin{bmatrix}C(a)\\ B(a)-TA(a)\\ V\end{bmatrix},
\]
with
\[
 I=[2n,6n),\qquad J=[n,10n),\qquad
 \Lambda=\{1,\ldots,n\}.
\]
\[
 C_{ij}=\frac{a^i-a^j}{j-i}\quad(i\ne j),\qquad C_{ii}=0,
 \qquad A_{ij}=a^i\delta_{ij}.
\]
\[
 t_j=\sum_{s=1}^j\frac{a^{j-s}}s,\qquad
 b_j=\sum_{s=1}^j\frac{a^{j-s}}{s^2},\qquad
 \]
 \[
 B_{ij}=\begin{cases}(t_i-t_j)/(j-i),&i\ne j,\\ b_i,&i=j,\end{cases}
 \qquad V_{\ell j}=\frac1{j+\ell}.
\]
Here $I,J,\Lambda$ are integer intervals, with increasing order.  At the seed
level $n=1$, the determinant is $9\times9$, with
$I=\{2,3,4,5\}$, $J=\{1,\ldots,9\}$, and $\Lambda=\{1\}$.
The exact source formulas give the following reductions.  Coefficient vectors
are in ascending powers of $T$; thus each row specifies the entire seed
polynomial modulo its indicated prime.
\begin{center}\normalsize\renewcommand{\arraystretch}{1.12}
\begin{tabular}{c c c l c}
\hline
$r$&$a$&$p_r$&$G^{(V)}_{4,-r}(T)\bmod p_r$&$G^{(V)}_{4,-r}(0)\bmod p_r$\\
\hline
2&$-2$&11&$(4,2,7,3,9)$&4\\
3&$-3$&11&$(1,4,6,7,1)$&1\\
4&$-4$&23&$(5,8,16,21,20)$&5\\
5&$-5$&11&$(5,2,5,8,7)$&5\\
6&$-6$&17&$(5,12,4,7,4)$&5\\
7&$-7$&17&$(8,0,8,16,5)$&8\\
\hline
\end{tabular}
\end{center}
For example, the first vector means
$4+2T+7T^2+3T^3+9T^4$ in $\mathbb F_{11}[T]$.  All six constant entries
are nonzero.  Each $p_r$ exceeds $7$.  The individual seed entries have
denominators dividing $2520^2$, and hence every coefficient denominator of
the seed determinant is supported on the primes $2,3,5,7$ (cancellation may
make it smaller, but need not make it divide $2520^2$).  Thus $p_r$ is prime
to every seed denominator, to $a$, and to $a-1$.  The seed constant is
consequently a $p_r$-adic unit and all seed coefficients are defined in
$\mathbb Z_{(p_r)}$.

\subsection{Signed Frobenius blocks}

Fix $r$ and abbreviate $p=p_r$.  Put $Q=p^k$ and $m=4Q$.  At this level the
three shape parameters of the V source are all set equal to \(Q\); they are
distinct from the exponent \(k\) in \(Q=p^k\).  Then the row and column
intervals are
\[
 I=[2Q,6Q),\qquad J=[Q,10Q),\qquad \Lambda=\{1,\ldots,Q\}.
\]
Define
\begin{equation}\label{eq:v-scaled-polynomial}
 P_{r,k}(X)=Q^{3m+Q}G^{(V)}_{m,-r}(X/Q^2)
              =Q^{13Q}G^{(V)}_{4Q,-r}(X/Q^2).
\end{equation}
The determinant on the left is formed by scaling its $C$ and $V$ rows by
$Q$ and its $B-TA$ rows by $Q^2$; the middle rows become $Q^2B-XA$.
Every nonzero denominator in a harmonic sum, index difference, or V-entry is
an integer whose absolute value is less than $11Q$.  Since $p\ge11$, its $p$-adic valuation is
at most $k$.  Hence the scaled matrix has entries in
$\mathbb Z_{(p)}[X]$, and so does $P_{r,k}(X)$.

Write $i=Qi_0+u$ and $j=Qj_0+v$, with residues in $\{0,\ldots,Q-1\}$.
As in the P0 calculation, only summands with denominator divisible by $Q$
survive after row scaling.  In particular,
\begin{equation}\label{eq:v-harmonic-filter}
 Q t_{Qj_0+u}\equiv a^u t_{j_0},\qquad
 Q^2 b_{Qj_0+u}\equiv a^u b_{j_0}\pmod p,
\end{equation}
because $a^Q\equiv a\pmod p$.  For $C$ and the off-diagonal $B$ entries,
unequal residues vanish: if $u\ne v$, then $v_p(j-i)<k$.  With equal
residues and unequal quotient indices, division by $j-i=Q(j_0-i_0)$ gives
the seed $C$ and $B$ entries with row factor $a^u$; here
$|j_0-i_0|\le7<p$.  The diagonal $B$ entry follows from
\eqref{eq:v-harmonic-filter}.  The diagonal $A$ entries satisfy
$A_{Qi_0+u,Qi_0+u}\equiv a^uA_{i_0i_0}\pmod p$.

For the V rows, fix a column residue $u$.  There is exactly one
$\ell\in\{1,\ldots,Q\}$ with $\ell+u\equiv0\pmod Q$, namely
$\ell=Q-u$ when $u>0$ and $\ell=Q$ when $u=0$.  All other scaled entries
vanish modulo $p$, and for this row
\[
 \frac{Q}{j+\ell}=\frac1{j_0+1},
\]
the seed V-entry.  Thus residue regrouping gives $Q$ diagonal blocks
\[
 \begin{bmatrix}a^u C_4\\ a^u(B_4-XA_4)\\ V_4\end{bmatrix},
 \qquad 0\le u<Q.
\]
Each block has eight rows carrying $a^u$ and one V row.  Row and column
regrouping contributes a sign $\varepsilon_{r,k}\in\{1,-1\}$, which is a
$p$-adic unit and need not be evaluated.  Therefore, coefficientwise,
\begin{equation}\label{eq:v-frobenius}
 P_{r,k}(X)\equiv
 \varepsilon_{r,k}a^{4Q(Q-1)}G^{(V)}_{4,-r}(X)^Q\pmod p.
\end{equation}
The exponent is $8\sum_{u=0}^{Q-1}u=4Q(Q-1)$.  The table above shows
$p\nmid G^{(V)}_{4,-r}(0)$, and $p\nmid a$, so the constant coefficient
on the right of \eqref{eq:v-frobenius} is nonzero modulo $p$.

\begin{lemma}[negative-base eventual nonvanishing]\label{lem:v-eventual-nonzero}
For each $r\in\{2,3,4,5,6,7\}$ and every fixed rational number $\xi$,
there is $k_0=k_0(r,\xi)$ such that
\[
 G^{(V)}_{4p_r^k,-r}(\xi)\ne0\qquad(k\ge k_0).
\]
\end{lemma}
\begin{proof}
Use the convention $v_{p_r}(0)=+\infty$.  For all sufficiently large $k$,
$v_{p_r}(Q^2\xi)=2k+v_{p_r}(\xi)>0$ (and this is immediate when
$\xi=0$).
Since $P_{r,k}\in\mathbb Z_{(p_r)}[X]$, its value at $X=Q^2\xi$ is
congruent modulo $p_r$ to its constant coefficient.  By
\eqref{eq:v-frobenius}, that coefficient is a unit.  Thus
$P_{r,k}(Q^2\xi)\ne0$, and the nonzero scaling in
\eqref{eq:v-scaled-polynomial} yields the claim.
\end{proof}

\subsection{Use in the determinant criterion}

For a fixed rational target $\xi$, the preceding lemmas give nonvanishing at
every sufficiently high index of the selected prime-power sequence, not
merely at one computed level.  Once $k\ge k_0$, the determinants at indices
$k$ and $k+1$ are both nonzero.  This assertion concerns the
selected prime-power subsequence only; no nonvanishing assertion is made for
every intervening integer degree $m$.

For each fixed integer base, the irrationality criterion of
Section~\ref{sec:criterion} must now be applied to one family at a time:
P0 with its base-dependent prime, or V with the corresponding $p_r$.
The integerizer, analytic estimate, and nonzero determinant must all belong
to that same source.  The prime-power sequences proved here supply exactly the
nonzero term required by the criterion; they do not by themselves imply a
common analytic or arithmetic rate across the two families.

\section{Hermite--Pad\'e tails}\label{sec:hp}
This section proves the two tails from the displayed construction, without
using an irrationality theorem as input.  The base is fixed throughout each
limit; constants in contour and remainder estimates may depend on that base.

\begin{theorem}\label{thm:hp-tail}
For every fixed integer $a\ge12$ or $a\le-8$, the determinant construction
below proves that $\operatorname{Li}_2(1/a)$ is irrational.
\end{theorem}

\subsection{Residues, tails, and exact integerization}
For $n\ge0$ put
\begin{equation}\label{eq:hp-phi}
 \Phi_n(w)=\frac{\prod_{r=n+1}^{3n+1}(w-r)}
                   {\prod_{j=0}^{n}(w-j)^2}
 =\sum_{j=0}^{n}\left(\frac{u_{n,j}}{(w-j)^2}
                         +\frac{v_{n,j}}{w-j}\right).
\end{equation}
Writing $H_k=\sum_{r=1}^k1/r$ and $H_0=0$, direct residue calculation gives
\begin{equation}\label{eq:hp-residues}
 u_{n,j}=-\frac{(3n+1-j)!}{(j!)^2((n-j)!)^3},\qquad
 v_{n,j}=u_{n,j}\bigl(3H_{n-j}-2H_j-H_{3n+1-j}\bigr).
\end{equation}
Set
\[
 B_n(z)=\sum_{j=0}^n v_{n,j}z^j,\qquad
 C_n(z)=\sum_{j=0}^n u_{n,j}z^j,
\]
and
\[
 A_n(z)=-\sum_{k=1}^{n}z^k\sum_{j=0}^{k-1}
       \left(\frac{v_{n,j}}{k-j}+\frac{u_{n,j}}{(k-j)^2}\right).
\]
The coefficient of $z^k$ in $B_n(-\log(1-z))+C_n\operatorname{Li}_2(z)$ is
the partial-fraction sum $\sum_{j<k}(v_{n,j}/(k-j)+u_{n,j}/(k-j)^2)$.
For $1\le k\le n$ it is canceled by $A_n$; for $n<k\le3n+1$ it is zero
because $\Phi_n(k)=0$; for $k\ge3n+2$ it equals $\Phi_n(k)$.  Therefore
\begin{equation}\label{eq:hp-tail-series}
 R_n(z):=A_n(z)+B_n(z)(-\log(1-z))+C_n(z)\operatorname{Li}_2(z)
       =\sum_{k\ge3n+2}\Phi_n(k)z^k.
\end{equation}
The adjacent Casoratian
\begin{equation}\label{eq:hp-casoratian}
 E_n(z):=B_{n+1}(z)R_n(z)-B_n(z)R_{n+1}(z)
       =W_n(z)+V_n(z)\operatorname{Li}_2(z)
\end{equation}
has
\[
 W_n=B_{n+1}A_n-B_nA_{n+1},\qquad
 V_n=B_{n+1}C_n-B_nC_{n+1}.
\]
The logarithm cancels identically.  Both $W_n,V_n$ have degree at most
$N=2n+1$.

Here is an all-degree proof of the coefficient formula needed for
integerization.  Let $b_n=\binom{2n+1}{n}$ and
\[
 J_n(t)=\sum_{k=0}^{N}(-1)^k\binom Nk\binom{n+k}{k}^{\!2}t^k.
\]
For $\Re(s)>-1$ the integral is absolutely convergent, and direct
termwise integration gives
\begin{equation}\label{eq:hp-mellin}
 \int_0^1t^sJ_n(t)\,dt=
 \frac{(2n+1)!}{(n!)^2}
 \frac{[s(s-1)\cdots(s-n+1)]^2}
      {(s+1)(s+2)\cdots(s+2n+2)}.
\end{equation}
The right side then supplies the meromorphic continuation in $s$; no
continuation is used before the stated half-plane of convergence.
Write $V_n(z)=\sum_{i=0}^Nv_i z^i$ and
$\Pi_n(t)=t^NV_n(1/t)=\sum_{j=0}^N\pi_jt^j$, so $\pi_j=v_{N-j}$.
The polynomial part of $E_n$ has degree at most $N$.  For
$k=N+\ell$, $1\le\ell\le n$, its coefficient is therefore
\[
 0=\sum_{i=0}^N\frac{v_i}{(k-i)^2}
   =\sum_{j=0}^N\frac{\pi_j}{(\ell+j)^2},
\]
which is $-\int_0^1t^{\ell-1}\Pi_n(t)\log t\,dt=0$.  The second
elimination, $C_{n+1}R_n-C_nR_{n+1}$, cancels $\operatorname{Li}_2$;
its nonpolynomial part is $V_n(z)\log(1-z)$.  The same coefficient range
gives
\[
 0=\sum_{i=0}^N\frac{v_i}{k-i}
   =\sum_{j=0}^N\frac{\pi_j}{\ell+j}
   =\int_0^1t^{\ell-1}\Pi_n(t)\,dt.
\]
Thus the two exact moment families are
\begin{equation}\label{eq:hp-moments}
 \int_0^1t^j\Pi_n(t)\,dt=0,\qquad
 \int_0^1t^j\Pi_n(t)\log t\,dt=0\quad(0\le j<n).
\end{equation}
The Mellin transform of any
polynomial of degree at most $N$ has denominator
$\prod_{r=1}^{N+1}(s+r)$ and numerator degree at most $N$.  The first family
in \eqref{eq:hp-moments} makes this numerator vanish at $s=0,\ldots,n-1$;
the second makes each zero double.  Hence the numerator is divisible by
$[s(s-1)\cdots(s-n+1)]^2$, leaving a factor of degree at most one.  The
solution space of \eqref{eq:hp-moments} consequently has dimension at most
two.

Equation~\eqref{eq:hp-mellin} shows that $J_n$ is one solution.  It also
shows $J_n(1)=0$, since the right side is $O(s^{-2})$ as $s\to+\infty$.
Integration by parts then shows that $tJ_n'(t)$ is a second solution: the
boundary terms vanish at both endpoints, and the two moment families for
$tJ'_n$ reduce to those for $J_n$.  The two solutions are independent by
their top two coefficients: if $a_N=[t^N]J_n$ and
$a_{N-1}=[t^{N-1}]J_n$, the determinant formed from the top two coefficients
of $J_n$ and $tJ_n'$ is
$a_N(N-1)a_{N-1}-Na_Na_{N-1}=-a_Na_{N-1}\ne0$.  Thus every solution is
their linear combination.  Direct substitution of the residues
\eqref{eq:hp-residues}
gives the constant and linear coefficients of $\Pi_n$; matching those two
coefficients in this two-dimensional space yields
\begin{equation}\label{eq:hp-vpoly}
 V_n(z)=b_n^2z^NP_n(1/z),
\end{equation}
where
\begin{equation}\label{eq:hp-Pn}
 P_n(t)=-\sum_{k=0}^{N}(-1)^k\binom Nk\binom{n+k}{k}
 \left((12n+17)\binom{n+k}{k}
       +(21n+29)\binom{n+k}{k-1}\right)t^k.
\end{equation}
The convention is $\binom{n+k}{k-1}=0$ at $k=0$.  In particular, the
second binomial is $\binom{n+k}{k-1}$, not
$\binom{n+k-1}{k-1}$.

The first two coefficients follow from the endpoint residues. Put
$\eta_m=-2H_m-H_{2m+1}$, $\rho_m=2m^2(m+1)$, and
$\kappa_m=3+2/m-1/(2m+2)$.  Their ratios are
\[
 \frac{v_{m,m}}{u_{m,m}}=\eta_m,\qquad
 \frac{u_{m,m-1}}{u_{m,m}}=\rho_m,\qquad
 \frac{v_{m,m-1}}{u_{m,m-1}}=\eta_m+\kappa_m.
\]
For $n\ge1$, expanding $V_n=B_{n+1}C_n-B_nC_{n+1}$ at its top two
degrees yields, with $\pi_j=[t^j]\Pi_n$,
\begin{align*}
 \pi_0&=u_{n,n}u_{n+1,n+1}(\eta_{n+1}-\eta_n),\\
 \pi_1&=u_{n,n}u_{n+1,n+1}
 \bigl((\rho_n+\rho_{n+1})(\eta_{n+1}-\eta_n)
       -\rho_n\kappa_n+\rho_{n+1}\kappa_{n+1}\bigr).
\end{align*}
Here $u_{n,n}=-(n+1)b_n$, $u_{n+1,n+1}=-2(2n+3)b_n$, and
$\eta_{n+1}-\eta_n=-5/(2(n+1))-1/(2n+3)$.  Substitution yields
$\pi_0=-b_n^2(12n+17)$ and
\[
 \pi_1=b_n^2(2n+1)(n+1)
       \bigl((n+1)(12n+17)+(21n+29)\bigr).
\]
For $n=0$, the residues give the same formulas directly: $\pi_0=-17$ and
$\pi_1=46$.  For the Casoratian calculation below, put
$h_m=3H_m-H_{3m+1}$ and
\[
 \delta_m=h_{m+1}-h_m=
 \frac{2(27m^2+54m+23)}{3(m+1)(3m+2)(3m+4)},
\]
The values of $\pi_0,\pi_1$ are the first two coefficients of
$-(12n+17)J_n-(21n+29)tJ'_n/(n+1)$.  The identity
$\frac{k}{n+1}\binom{n+k}{k}=\binom{n+k}{k-1}$ gives
\eqref{eq:hp-Pn}.  The dimension bound proves \eqref{eq:hp-vpoly} for every
$n$, with no finite-range restriction.

For reference, the boundary residues used in this matching can be checked
directly from \eqref{eq:hp-residues}.  For $m\ge1$,
\begin{align}
 u_{m,m-1}&=-\frac{(2m+2)!}{((m-1)!)^2},\qquad
 u_{m,m}=-\frac{(2m+1)!}{(m!)^2},\\
 \frac{u_{m,m-1}}{u_{m,m}}&=(2m+2)m^2=2m^2(m+1),\\
 \frac{v_{m,m-1}}{u_{m,m-1}}-\frac{v_{m,m}}{u_{m,m}}
   &=3+\frac2m-\frac1{2m+2},\\
 \frac{u_{n,n}u_{n+1,n+1}}{b_n^2}&=(2n+2)(2n+3).
\end{align}
These identities also cover the boundary-residue normalization in the
three-form extraction below; the case $n=0$ is read directly from
\eqref{eq:hp-residues}.

Let $d_N=\operatorname{lcm}(1,\ldots,N)$ and
$Q_n(z)=V_n(z)/b_n^2\in\mathbb Z[z]$.  Since $E_n=W_n+V_n\operatorname{Li}_2$
vanishes through degree $N$, its polynomial part satisfies
\[
 \frac{W_n(z)}{b_n^2}=-[Q_n(z)\operatorname{Li}_2(z)]_{\le N}.
\]
Each coefficient on the right has denominators among $r^2$ with
$1\le r\le N$, and $r\mid d_N$.  Therefore
\begin{equation}\label{eq:hp-integer-poly}
 \frac{d_N^2}{b_n^2}W_n,\quad
 \frac{d_N^2}{b_n^2}V_n\in\mathbb Z[z].
\end{equation}
For every fixed integer $a$ with $|a|\ge2$, put
\begin{equation}\label{eq:hp-integer-form}
 L_n(a)=a^N\frac{d_N^2}{b_n^2}E_n(1/a)
       =X_n(a)+Y_n(a)\operatorname{Li}_2(1/a),
       \qquad X_n(a),Y_n(a)\in\mathbb Z.
\end{equation}
The factor $a^N$ clears evaluation denominators because both polynomials in
\eqref{eq:hp-integer-poly} have degree at most $N$.

\subsection{Remainder bound for both signs}
For $-1<z<1$, beta integration of \eqref{eq:hp-tail-series} gives the exact
identity
\begin{equation}\label{eq:hp-beta}
 R_n(z)=\frac{(2n+1)!}{(n!)^2}z^{3n+2}
 \int_0^1\!\int_0^1\frac{tu}{(1-ztu)^2}
 \left(\frac{t^2u^2(1-t)(1-u)}{(1-ztu)^2}\right)^n dt\,du.
\end{equation}
To maximize the bracket, fix $p=tu$.  Since $t+u\ge2\sqrt p$,
$(1-t)(1-u)=1-(t+u)+p\le(1-\sqrt p)^2$, with equality only at
$t=u=\sqrt p$.  Write $s=\sqrt p$.  The bracket is at most
\[
 g_z(s)=\frac{s^4(1-s)^2}{(1-zs^2)^2},\qquad 0\le s\le1.
\]
For $|z|<1$, its unique interior maximum satisfies
$zs^3-3s+2=0$: differentiating $\log g_z$ gives
$2-3s+zs^3=0$, and the left side is strictly decreasing on $[0,1]$.
At that root, $1-zs^2=2(1-s)/s$, so $g_z(s)=s^6/4$.  The prefactor
$tu/(1-ztu)^2$ in \eqref{eq:hp-beta} is bounded for fixed $z$ with
$|z|<1$.  Stirling's formula~\cite{WhittakerWatson} and the unit area of the square therefore give
\begin{equation}\label{eq:hp-R-rate}
 \limsup_{n\to\infty}|R_n(z)|^{1/n}\le |z|^3s(z)^6.
\end{equation}
No cancellation estimate is used.

\subsection{Positive coefficient bound}
For $z=1/q$, $q>1$, all $u_{n,j}$ have the same (negative) sign.  Put
$x=j/n$ and
\[
 f_q(x)=(3-x)\log(3-x)-2x\log x-3(1-x)\log(1-x)-x\log q,
 \quad 0\le x\le1,
\]
where $0\log0=0$.  Uniform Stirling estimates~\cite{WhittakerWatson} in
\eqref{eq:hp-residues} give
\[
 \log\bigl(|u_{n,j}|q^{-j}\bigr)=nf_q(j/n)+O(\log(n+1))
\]
uniformly for $0\le j\le n$.  Indeed, the coefficients of $n\log n$ and
$n$ cancel, and all factorial contributions are $O(\log(n+1))$,
including at endpoints.  Now
\[
 f_q'(x)=\log\frac{(1-x)^3}{q x^2(3-x)},\qquad
 f_q''(x)=\frac1{3-x}-\frac2x-\frac3{1-x}<0.
\]
There is a unique maximizer $t(q)\in(0,1)$, determined by
$(1-t)^3=qt^2(3-t)$.  Substitution of this equation gives
\[
 \exp(f_q(t))=\left(\frac{3-t}{1-t}\right)^3.
\]
Since the sum has only $n+1$ terms,
\[
 \limsup_{n\to\infty}|C_n(1/q)|^{1/n}
 \le \left(\frac{3-t(q)}{1-t(q)}\right)^3.
\]
Also $|3H_{n-j}-2H_j-H_{3n+1-j}|\le6(1+\log(3n+1))$, so the same bound
holds for $B_n(1/q)$.  This upper-bound form of the discrete Laplace
principle is all that is needed; no unproved ratio limit is used.

\subsection{Negative coefficient bound: contour and endpoint estimate}
Fix $q>1$, put $z=-1/q$, and define the two logarithmic branches
\[
 F_n^\pm(z)=B_n(z)+C_n(z)(-\log q\pm i\pi).
\]
Residues of $\Phi_n(w)e^{(-\log q\pm i\pi)w}$ give
\begin{equation}\label{eq:hp-contour-res}
 F_n^\pm(z)=\frac1{2\pi i}\oint_{\Gamma_n^\pm}
          \Phi_n(w)e^{(-\log q\pm i\pi)w}\,dw,
\end{equation}
where the rectangle encloses the poles $0,1,\ldots,n$.  Set $\beta_q>0$
equal to the absolute imaginary part of the lower nonreal root of
\begin{equation}\label{eq:hp-saddle-cubic}
 p_q(w)=(q+1)w^3-3(q+1)w^2+3w-1=0.
\end{equation}
For all sufficiently large $n$, let $\Gamma_n^+$ be the positively oriented
rectangle with vertical sides $\Re w=-1/2,n+1/2$ and horizontal sides
$\Im w=\pm n\beta_q$.  It encloses precisely the poles $0,1,\ldots,n$.
Its distance from each enclosed pole is at least
$\min\{1/2,n\beta_q\}$, hence at least $1/2$ once
$n\beta_q\ge1/2$.  All numerator zeros lie outside the rectangle at distance
at least $1/2$.  The rate proof uses only these sufficiently large indices.
For $F_n^-$ use the conjugate exponential branch on the same positively
oriented rectangle.

We give the uniform sum-to-integral estimate needed on these contours.
If $\xi=w/n$ lies on the scaled rectangle, then
\begin{equation}\label{eq:hp-log-em}
 \log|\Phi_n(n\xi)|=n\left(
 \int_1^3\log|\xi-r|\,dr-2\int_0^1\log|\xi-r|\,dr\right)
 +O_q(\log(n+1)),
\end{equation}
uniformly along all four sides.  On a horizontal side $|\Im\xi|=\beta_q$,
the functions $r\mapsto\log|\xi-r|$ have uniformly bounded derivative, so
the ordinary left/right Riemann-sum error is $O_q(1)$.  On the left
vertical side $\Re\xi=-1/(2n)$, the real part lies outside both integration
intervals by at least $1/(2n)$.  On the right side
$\Re\xi=1+1/(2n)$ it lies outside $[0,1]$ by the same amount.  For $[1,3]$
there, the only possible near-singularity lies in the first mesh cell:
its real part is exactly halfway between the endpoint and the first mesh
point.  If $|\Im\xi|\le1/n$, rescaling this one cell by $r=1+u/n$ shows
that both its mesh contribution and $n$ times its integral are
$O(\log(n+1))$; the integral of $\log\sqrt{u^2+v^2}$ over a fixed
interval is bounded uniformly for $0\le v\le1$.  If
$|\Im\xi|>1/n$, or on the remaining cells, the total variation is bounded
by $C_q+\int_{1/n}^{O_q(1)}dr/r=O_q(\log(n+1))$.  Endpoint shifts of one
mesh point in the numerator product contribute at most another
$O_q(\log(n+1))$.  This proves \eqref{eq:hp-log-em}; in particular the
logarithmic singularity is handled explicitly rather than hidden in a
smooth-function Euler--Maclaurin assertion.

For completeness, the estimate in \eqref{eq:hp-log-em} is a bounded-variation
Riemann-sum estimate, including the endpoint cell.  Put
$f_\xi(r)=\log|\xi-r|$.  The numerator sum is the right Riemann sum on
$[1,3]$, with the extra point $f_\xi(3+1/n)$, and the denominator sum is the
right Riemann sum on $[0,1]$, with $f_\xi(0)$.  On a horizontal edge the
total variation is $O_q(1)$.  On a vertical edge the only possible singular
cell is the first half-mesh cell; after $r=1+u/n$ (or $r=u/n$) its mesh term
and its $n$-scaled integral are $O(\log(n+1))$.  The remaining variation is
$O_q(1)+\int_{1/n}^{O_q(1)}dr/r=O_q(\log(n+1))$.  Thus the estimate is
uniform on all four edges and does not invoke a smooth Euler--Maclaurin
formula across a logarithmic endpoint.

Combining \eqref{eq:hp-log-em} with the exact modulus of the exponential
factor gives
\[
 \log|\Phi_n(n\xi)e^{(-\log q+i\pi)n\xi}|
       =nH_q(\xi)+O_q(\log(n+1)),
\]
where
\begin{equation}\label{eq:hp-H}
 H_q(\xi)=\int_1^3\log|\xi-r|\,dr-2\int_0^1\log|\xi-r|\,dr
          -\log q\,\Re\xi-\pi\Im\xi.
\end{equation}
The discriminant of \eqref{eq:hp-saddle-cubic} is
$-108q^2(q+1)<0$, and
$p_q(2)=1-4q<0<p_q(3)=8$.  Thus the real root is in $(2,3)$; the
conjugate roots have real part $u=(3-R)/2\in(0,1/2)$.  Write the lower
root as $\xi_0=u-i\beta_q$.  Differentiating the analytic logarithmic
phase underlying \eqref{eq:hp-H} shows $\partial_xH_q(\xi_0)=0$ by the
saddle equation $z(\xi-1)^3=\xi^2(\xi-3)$.  Along the lower side,
\[
 \frac{d^2}{dx^2}H_q(x-i\beta_q)
 =-6\Re\frac1{(x-i\beta_q)(x-1-i\beta_q)(x-3-i\beta_q)}<0
 \quad(0\le x\le1),
\]
because the real part of the denominator is
$x(x-1)(x-3)+\beta_q^2(4-3x)>0$.  Hence the unique lower-side maximum is
$\xi_0$.  On the vertical sides, parameterized downward,
\[
 \frac d{dy}H_q(-iy)=3\arctan y-\arctan(y/3)>0,
\]
\[
 \frac d{dy}H_q(1-iy)=\frac\pi2+2\arctan y-\arctan(y/2)>0.
\]
On the lower halves of the limiting vertical sides, parameterized as
$\xi=-iy$ and $\xi=1-iy$ with $0<y\le\beta_q$, the displayed derivatives
are strictly positive.  On the upper halves use the exact reflection identity
$H_q(x+iy)=H_q(x-iy)-2\pi y$ for $y>0$.  Hence neither upper vertical half
can exceed its reflected lower half, and the maxima on both lower vertical
halves occur at their lower endpoints.  The strictly concave lower horizontal
phase has its unique maximum at the interior saddle $\xi_0$, strictly above
both endpoints.

The logarithmic potentials defining $H_q$ are continuous when $\xi$ approaches
the real integration intervals because their logarithmic singularities are
integrable.  They are therefore uniformly continuous on a compact
neighborhood of the limiting rectangle.  The maxima over the boundaries of
the scaled finite rectangles converge to the maximum over the limiting
boundary.  The original rectangles have perimeter $O_q(n)$, while
\eqref{eq:hp-log-em} contributes only a polynomial factor in $n$ to the
modulus bound.  Taking $n$-th roots now gives
the claimed rate bound.
\begin{equation}\label{eq:hp-F-rate}
 \limsup_{n\to\infty}|F_n^\pm(-1/q)|^{1/n}\le e^{H_q(\xi_0)}.
\end{equation}

To evaluate this number, choose logarithms continuously from the lower
contour.  An antiderivative gives the analytic phase
\[
 \Psi_q(\xi)=3(\xi-1)\operatorname{Log}(\xi-1)- (\xi-3)\operatorname{Log}(\xi-3)
             -2\xi\operatorname{Log}\xi+(-\log q+i\pi)\xi.
\]
At the lower saddle choose the principal logarithms of
$\xi_0,\xi_0-1,\xi_0-3$, whose arguments lie in $(-\pi,0)$ (with the
negative-real limiting value approached from below).  The algebraic saddle
equation implies that $\Psi_q'(\xi_0)$ is an integer multiple of $2\pi i$.
As $q\to\infty$, the lower root satisfies
$\xi_0=-i/\sqrt{3q}+O(q^{-1})$; the three arguments tend respectively to
$-\pi/2,-\pi,-\pi$.  The corresponding imaginary-part combination is zero,
so continuity in $q>1$ forces this integer multiple to vanish identically.
Substituting the saddle equation into $\Psi_q$ gives
\[
 e^{H_q(\xi_0)}=\left|\frac{\xi_0-3}{\xi_0-1}\right|^3.
\]
For $y=(\xi-3)/(\xi-1)$ the saddle equation becomes
$y(y-3)^2=4z$.  Its real root is $y_0=-s_-^2/q$, where
$s_-(q)\in(0,2/3)$ is the unique solution of
\begin{equation}\label{eq:hp-sminus}
 -\frac{s_-^3}{q}-3s_-+2=0.
\end{equation}
The product of the three roots of $y(y-3)^2=4z$ is $4z=-4/q$.
Consequently the squared modulus of either nonreal root is
$4/s_-^2$, and
\begin{equation}\label{eq:hp-negative-coeff-rate}
 e^{H_q(\xi_0)}=\frac8{s_-(q)^3}.
\end{equation}
Finally
\[
 C_n(-1/q)=\frac{F_n^+-F_n^-}{2\pi i},\qquad
 B_n(-1/q)=\frac{F_n^++F_n^-}{2}+\log q\,C_n(-1/q).
\]
Equations~\eqref{eq:hp-F-rate}--\eqref{eq:hp-negative-coeff-rate} thus give
the same upper exponential rate $8/s_-^3$ for $B_n(-1/q)$.

\subsection{Integer-form rates and exact cutoffs}
The prime number theorem gives
$\log d_N=N+o(N)$, while Stirling~\cite{WhittakerWatson} gives
$n^{-1}\log b_n\to2\log2$.  With $N=2n+1$,
\begin{equation}\label{eq:hp-den-rate}
 \lim_{n\to\infty}\left(\frac{d_N^2}{b_n^2}\right)^{1/n}
       =\frac{e^4}{16}.
\end{equation}
No finite global inequality for $d_N$ is asserted.  Combining
\eqref{eq:hp-R-rate}, the coefficient bounds, and
\eqref{eq:hp-casoratian} yields
\begin{align}
 \limsup |L_n(q)|^{1/n}&\le
 r_+(q):=\frac{e^4}{16q}
       \left(\frac{3-t(q)}{1-t(q)}\right)^3s_+(q)^6,
       &&a=q>1,\\
 \limsup |L_n(-q)|^{1/n}&\le
 r_-(q):=\frac{e^4s_-(q)^3}{2q},
       &&a=-q<-1,
\end{align}
where $s_+$ is the unique root in $(0,1)$ of
$s_+^3/q-3s_++2=0$.  The factors are as follows: $q^2$ from $|a|^N$,
$e^4/16$ from \eqref{eq:hp-den-rate}, $q^{-3}s^6$ from the remainder,
and the relevant coefficient growth rate.  Thus no estimate for an
uneliminated logarithmic form is being used.

For the positive rate set $z=1/q$.  Implicit differentiation of the two
root equations gives
\[
 \frac{d\log r_+}{d\log z}
       =1+t+\frac{2zs_+^2}{1-zs_+^2}>0,
\]
so $r_+(q)$ decreases with $q$.  On the negative side,
$q=s_-^3/(2-3s_-)$ is increasing in $s_-$ and
$r_-(q)=e^4(2-3s_-)/2$ is decreasing.  It remains to check four rational
endpoints.  The Taylor sum for $e^4$ through degree 16 and a geometric tail
bound give
\[
 \frac{6972316451}{127702575}<e^4<\frac{377139364267}{6907548375}.
\]
Opposite signs of the defining cubics at the following rational endpoints
isolate the roots.  The derivatives are
$-3(1-t)^2-q(6t-3t^2)<0$ for $t(q)$,
$3s^2/q-3<0$ for $s_+(q)$, and $-3s^2/q-3<0$ for $s_-(q)$:
\[
\begin{array}{c|c}
 t(11)&(14179/100000,14180/100000)\\
 s_+(11)&(67602/100000,67604/100000)\\
 t(12)&(13681/100000,13683/100000)\\
 s_+(12)&(67521/100000,67523/100000)\\
 s_-(7)&(65338/100000,65339/100000)\\
 s_-(8)&(65496/100000,65497/100000)
\end{array}
\]
Exact rational interval substitution gives
\[
\begin{array}{c|c}
 r_+(11)&(10937/10000,10941/10000)\\
 r_+(12)&(9834/10000,9837/10000)\\
 r_-(7)&(10877/10000,10879/10000)\\
 r_-(8)&(9587/10000,9589/10000)
\end{array}
\]
and hence $r_+(12)<1$ and $r_-(8)<1$.  Monotonicity proves
$r_+(q)<1$ for every integer $q\ge12$ and $r_-(q)<1$ for every integer
$q\ge8$.

\subsection{All-degree nonvanishing and conclusion}
It remains to show that the constructed integer forms do not vanish
eventually.
Let $U_m=(3m+1)!/(m!)^3$,
$K_n=((2n+1)!)^3/((3n+2)!)^2$, and
$h_m=3H_m-H_{3m+1}$.  Direct extraction from
\eqref{eq:hp-tail-series} gives
\[
 \Phi_n(3n+2)=K_n,\qquad
 \frac{\Phi_n(3n+3)}{K_n}=\frac{8(n+1)}9.
\]
For $m\ge0$, put $M_m=B_mC_{m+1}-C_mB_{m+1}$; in particular
$M_0=B_0C_1-C_0B_1$ is defined by the displayed residue polynomials.
For $m\ge1$, put $r_m=m^3/(3m+1)$ and $c_m=2+3/m-1/(3m+1)$.
The lower endpoint residues give
\[
 u_{m,0}=-U_m,\quad u_{m,1}=r_m u_{m,0},\quad
 v_{m,0}=h_m u_{m,0},\quad
 v_{m,1}=r_m(h_m-c_m)u_{m,0}.
\]
Expanding $B_mC_{m+1}-C_mB_{m+1}$ at degrees zero and one gives
\[
 M_m(0)=-U_mU_{m+1}\delta_m,
\]
\[
 [z]M_m=U_mU_{m+1}
 \{-\delta_m(r_m+r_{m+1})+r_{m+1}c_{m+1}-r_mc_m\},
\]
The three-form determinant
\[
 \Delta_n(z)=\det\begin{bmatrix}
 A_n&B_n&C_n\\ A_{n+1}&B_{n+1}&C_{n+1}\\ A_{n+2}&B_{n+2}&C_{n+2}
 \end{bmatrix}
\]
has degree at most $3n+3$.  Replacing its first column by the corresponding
$R$-column and using \eqref{eq:hp-tail-series} shows its order is at least
$3n+2$.  Expanding along that column gives
\[
 \Delta_n=R_nM_{n+1}
 -R_{n+1}(B_nC_{n+2}-C_nB_{n+2})+R_{n+2}M_n.
\]
The last two terms start at degree $3n+5$.  Also
\[
 R_n(z)=K_nz^{3n+2}\left(1+\frac{8(n+1)}9z+O(z^2)\right).
\]
It follows that
\[
 [z^{3n+2}]\Delta_n=K_nM_{n+1}(0),\qquad
 \frac{[z^{3n+3}]\Delta_n}{[z^{3n+2}]\Delta_n}
 =\frac{[z]M_{n+1}}{M_{n+1}(0)}+\frac{8(n+1)}9.
\]
Substitution of the displayed coefficients gives the ratio $-\theta_n$, with
\[
 \theta_n=
 \frac{-\delta_{n+1}(r_{n+1}+r_{n+2})
       +r_{n+2}c_{n+2}-r_{n+1}c_{n+1}}{\delta_{n+1}}
       -\frac{8(n+1)}9.
\]
Substituting $r_m,c_m,\delta_m$ and clearing denominators gives
\begin{equation}\label{eq:hp-theta}
 \theta_n=\frac{1728n^4+11520n^3+28140n^2+29965n+11770}
 {18(3n+4)^2(27n^2+108n+104)}.
\end{equation}
Therefore
\begin{equation}\label{eq:hp-delta}
 \Delta_n(z)=-K_nU_{n+1}U_{n+2}\delta_{n+1}z^{3n+2}(1-\theta_nz),
\end{equation}
The numerator is positive for $n\ge0$; denominator minus numerator is
\[
 2646n^4+17640n^3+43140n^2+46067n+18182>0.
\]
Also $\delta_{n+1}>0$.  Therefore $\Delta_n(z)\ne0$ for every real
$0<|z|<1$.

For such $z$, the beta integral \eqref{eq:hp-beta} gives $R_m(z)\ne0$.
If both $E_n(z)$ and $E_{n+1}(z)$ vanished, the two adjacent minors of the
columns $A+\operatorname{Li}_2(z)C$ and $B$ would vanish.  Since
$R_n(z)$ and $R_{n+1}(z)$ are nonzero, this makes the three columns in
$\Delta_n(z)$ linearly dependent (the conclusion also holds when one of the
corresponding $B$ entries is zero, by using the other adjacent equation).
Hence $\Delta_n(z)=0$, contradicting \eqref{eq:hp-delta}.  Thus no two
consecutive $E_n(z)$ vanish.

Suppose now that $\operatorname{Li}_2(1/a)=r/s$ were rational, in lowest
terms.  By \eqref{eq:hp-integer-form}, $sL_n(a)$ is an integer for every
$n$.  The strict rate bounds imply $L_n(a)\to0$, so $sL_n(a)=0$ for all
sufficiently large $n$.  This contradicts the fact that no two consecutive
$E_n(1/a)$, and hence no two consecutive $L_n(a)$, vanish.  This proves
Theorem~\ref{thm:hp-tail}.


\section{Assembly of the proof}
The positive construction supplies the criterion for $2\le n\le11$, the
signed construction supplies it for $-7\le n\le-2$, and the Hermite--Pad\'e
construction supplies it for $n\ge12$ and $n\le-8$.  Together with the two
classical evaluations at $n=\pm1$, these ranges prove Theorem~\ref{thm:main}.

\section*{Use of AI tools}
The author devised the method.  Codex (OpenAI, GPT-5.6 Sol) assisted with
mathematical exploration and implementation.  Claude Code (Anthropic,
Claude Opus 5.5) assisted with review of the implementation, parameters, and
printed data; that review is not a proof of the mathematical claims.  The
directed numerical bounds are computer-assisted:
they were obtained with outward-rounded interval arithmetic as specified in
Appendix~\ref{sec:certificate-algorithm}.  The author is responsible for the
mathematical statements and their interpretation.

\section*{Declarations}
\textbf{Competing interests.} The author has no competing interests.

\textbf{Funding.} No funding was received for this work.

\textbf{Data availability.} The ancillary files contain the printed-data
replay and its logs. The source-range script checks text and table structure
only; it does not certify the mathematical estimates.

\begin{thebibliography}{99}

\bibitem{Andreief}
C.~Andr\'eief,
Note sur une relation entre les int\'egrales d\'efinies des produits des
fonctions,
\emph{M\'em. Soc. Sci. Phys. Nat. Bordeaux} (3) \textbf{2} (1886), 1--14.

\bibitem{Davenport}
H.~Davenport,
\emph{Multiplicative Number Theory}, 3rd ed., revised by H.~L.~Montgomery,
Graduate Texts in Mathematics~74, Springer, New York, 2000.

\bibitem{Forrester}
P.~J.~Forrester,
Meet Andr\'eief, Bordeaux 1886, and Andreev, Kharkov 1882--1883,
\emph{Random Matrices Theory Appl.} \textbf{8} (2019), no.~2, 1930001.

\bibitem{Gantmacher}
F.~R.~Gantmacher,
\emph{The Theory of Matrices}, Vol.~1,
Chelsea, New York, 1959.

\bibitem{Hata}
M.~Hata,
Rational approximations to the dilogarithm,
\emph{Trans. Amer. Math. Soc.} \textbf{336} (1993), no.~1, 363--387.

\bibitem{Krattenthaler}
C.~Krattenthaler,
Advanced determinant calculus,
\emph{S\'em. Lothar. Combin.} \textbf{42} (1999), Art.~B42q, 67~pp.

\bibitem{Kummer}
E.~E.~Kummer,
\"Uber die Erg\"anzungss\"atze zu den allgemeinen Reciprocit\"atsgesetzen,
\emph{J. Reine Angew. Math.} \textbf{44} (1852), 93--146.

\bibitem{Legendre}
A.-M.~Legendre,
\emph{\'El\'ements de g\'eom\'etrie}, Firmin Didot, Paris, 1794, Note~IV.

\bibitem{Lucas}
E.~Lucas,
Th\'eorie des fonctions num\'eriques simplement p\'eriodiques,
\emph{Amer. J. Math.} \textbf{1} (1878), 184--240 and 289--321.

\bibitem{Prellberg}
T.~Prellberg,
Three irrationality results for the dilogarithm: Rhin--Viola and
Viola--Zudilin constructions at $-1/4$, $1/5$, and $-1/3$,
preprint, arXiv:2609.06895, 2026.

\bibitem{RhinViola}
G.~Rhin and C.~Viola,
The permutation group method for the dilogarithm,
\emph{Ann. Sc. Norm. Super. Pisa Cl. Sci.} (5) \textbf{4} (2005), no.~3,
389--437.

\bibitem{SaffTotik}
E.~B.~Saff and V.~Totik,
\emph{Logarithmic Potentials with External Fields},
Grundlehren der mathematischen Wissenschaften~316, Springer, Berlin, 1997.

\bibitem{WhittakerWatson}
E.~T.~Whittaker and G.~N.~Watson,
\emph{A Course of Modern Analysis}, 4th ed.,
Cambridge University Press, Cambridge, 1927.

\end{thebibliography}

\appendix
\section{Positive endpoint data}\label{sec:positive-data}
The tables use ordinary decimal integers.  Divide each mass list by the
displayed denominator; the contour list gives the first half of the
conjugate-symmetric Fej\'er weights, with zero filling at the remaining nodes.

\renewcommand{\arraystretch}{1.12}
{\normalsize
\begin{longtable}{@{}p{0.29\linewidth}p{0.66\linewidth}@{}}
\caption{$q=2$ endpoint (27 independent masses).}\label{tab:q2data}\\
\toprule
item & exact value \\
\midrule
denominator & $5000000000000$\\
pinned edges & $8,12,16,20,24,28,32,36,40$ divided by $32$\\
pinned masses & $16526887295,294754833065,416674605400,489006355115,$\\
& $568607559425,625000000000,625000000000,625000000000$\\
frequency edges & $1/2,9/10,13/10,17/10,21/10,5/2,29/10,33/10,37/10,41/10,9/2$\\
frequency masses & $1839928019547,925598240328,604396185629,$\\
& $446793550571,346483328360,273707087420,$\\
& $215955707916,166118856466,118641146544,62377877219$\\
contour pair masses & $103692576010,101009799410,96050653090,$\\
& $89324463765,81159146555,71580231315,60226824760,$\\
& $46279839230,20391345715$\\
\bottomrule
\end{longtable}
}

{\normalsize
\begin{longtable}{@{}p{0.29\linewidth}p{0.66\linewidth}@{}}
\caption{$q=11$ endpoint (20 independent masses).}\label{tab:q11data}\\
\toprule
item & exact value \\
\midrule
denominator & $3000000000000$\\
pinned edges & $8,12,16,20,24,28,32,36,40$ divided by $32$\\
pinned masses & $28,0,107702656325,206939134529,260854191850,$\\
& $321658806124,374999999996,374999999996$\\
frequency edges & $1/2,3/2,5/2,7/2,9/2$\\
frequency masses & $2334546252656,652078675810,13375071534,0$\\
contour pair masses & $105521636098,103136532801,98660112985,92490547592,$\\
& $84344850292,75717735060,60802883424,55748307324$\\
\bottomrule
\end{longtable}
}

\section{Exact signed reference data}\label{sec:v-reference-data}
Every entry below is an ordinary signed base-ten integer. For real chains,
divide each length integer by $10^8$ and each mass increment by $10^8$. A
single listed correction $(j,r)$ adds $r/(5\cdot10^{13})$ to length $j$;
$j=0$ means no correction. Start each chain at its printed endpoint (units
$1/(5\cdot10^{13})$), then accumulate lengths from left to right. Start
each mass at zero and accumulate $\Delta m_i$. For Poisson rows divide all
increments by $10^8$, accumulate coordinates and masses from zero, then append
the reversed conjugates of all but the final (real) prefix row. Entries are
read left to right and then top to bottom. These conventions, together with
the displayed masses and supports, determine every reference measure uniquely.
\begin{VChainTable}{2}{x}{41}{3349365000000}{0}{0}
1727789 & 904764 & 1905814 & 711734 & 2075677 & 286937 & 4624699 & 2447859\\
2529220 & -2005912 & 2659559 & 74649 & 2778512 & 56594 & 2885565 & 43329\\
2980263 & 33330 & 3062199 & 25792 & 3131021 & 19868 & 19467262 & 13206641\\
3228206 & -13157141 & 3186436 & 336 & 3131021 & -423 & 3062199 & -1053\\
2980263 & -1592 & 2885565 & -1972 & 2778512 & -2210 & 2659559 & -2400\\
2529220 & -2556 & 2388048 & -2629 & 2236651 & -2632 & 2075677 & -2616\\
1905814 & -2625 & 1727789 & -2659 & 1542367 & -2693 & 1350340 & -2849\\
1152529 & -3266 & 949785 & -4562 & 742973 & -5487 & 532979 & -24509\\
320703 & 20464 & 4075 & -2499124 & 11603 & 73573 & 17365 & 82752\\
20484 & 76534 & 20484 & 66626 & 17365 & 54610 & 11604 & 21958\\
4074 & 54225 &  &  &  &  &  & \\
\end{VChainTable}
\begin{VChainTable}{2}{y}{48}{480367989919}{1}{10081}
742973 & 1385843 & 949785 & 680059 & 1152529 & 209735 & 1350340 & 128408\\
1542367 & 75603 & 1727789 & 50189 & 1905814 & 35139 & 2075677 & 26034\\
2236651 & 20335 & 2388048 & 16778 & 2529220 & 14604 & 2659559 & 13425\\
2778512 & 12980 & 2885565 & 13054 & 2980263 & 13531 & 3062199 & 14320\\
3131021 & 15294 & 16211108 & 11164507 & 3256154 & -11054943 & 3228206 & 16310\\
3186436 & 19197 & 3131021 & 24778 & 3062199 & 19543 & 2980263 & 18296\\
2885565 & 16770 & 2778512 & 15360 & 2659559 & 13983 & 2529220 & 12669\\
2388048 & 11386 & 2236651 & 10131 & 2075677 & 8937 & 1905814 & 7774\\
1727789 & 6583 & 1542367 & 5377 & 1350340 & 4050 & 1152529 & 2463\\
949785 & -131 & 742973 & -2258 & 532979 & -25599 & 320703 & 26042\\
4075 & -2927729 & 11603 & 86195 & 17365 & 96969 & 20484 & 89696\\
20484 & 78091 & 17365 & 64014 & 11604 & 25747 & 4074 & 63559\\
\end{VChainTable}
\begin{VChainTable}{2}{z\in[-q,-3/5]}{62}{-100000000000000}{0}{0}
5704 & 244590 & 16245 & -24961 & 24311 & -10122 & 28678 & -25136\\
28677 & -30658 & 24312 & -35202 & 16244 & -38047 & 5704 & -33807\\
448985 & 1148248 & 746170 & -11922 & 1040162 & 7481 & 1329699 & -2526\\
1613542 & -4230 & 1890475 & -6119 & 2159314 & -7630 & 2418905 & -9040\\
2668139 & -10408 & 2905947 & -11807 & 3131312 & -13194 & 3343268 & -14608\\
3540907 & -16105 & 3723384 & -17629 & 3889916 & -19173 & 4039792 & -20794\\
4172368 & -22487 & 4287077 & -24221 & 4383430 & -26007 & 4461010 & -27859\\
4519490 & -29766 & 4558614 & -31723 & 4578219 & -33734 & 4578219 & -35765\\
4558614 & -37788 & 4519490 & -39788 & 4461010 & -41748 & 4383430 & -43637\\
4287077 & -45399 & 4172368 & -46965 & 4039792 & -48254 & 3889916 & -49200\\
3723384 & -49730 & 3540907 & -49748 & 3343268 & -49123 & 3131312 & -47768\\
2905947 & -45635 & 2668139 & -42681 & 2418905 & -38860 & 2159314 & -34222\\
1890475 & -28926 & 1613542 & -23184 & 1329699 & -17235 & 1040162 & -11410\\
746170 & -5884 & 448985 & -2342 & 5704 & -5971 & 16244 & 172\\
24312 & 144 & 28677 & 104 & 28678 & 57 & 24311 & 23\\
16245 & -24 & 5704 & 28 &  &  &  & \\
\end{VChainTable}
\begin{VPoissonTable}{2}{149}{75}
44433733 & 38967338 & 63870 & -1298766 & 1432967 & 213587 & -1344956 & 1389706 & -12589\\
-1389706 & 1344956 & 48627 & -1432967 & 1298766 & 23616 & -1474695 & 1251185 & 24942\\
-1514843 & 1202264 & 20256 & -1553368 & 1152057 & 18160 & -1590231 & 1100615 & 15953\\
-1625391 & 1047995 & 14237 & -1658810 & 994252 & 12744 & -1690453 & 939446 & 11446\\
-1720285 & 883632 & 10308 & -1748277 & 826873 & 9288 & -1774394 & 769228 & 8392\\
-1798613 & 710760 & 7582 & -1820906 & 651531 & 6854 & -1841248 & 591603 & 6179\\
-1859618 & 531042 & 5554 & -1875999 & 469913 & 4983 & -1890369 & 408281 & 4450\\
-1902716 & 346211 & 3961 & -1913024 & 283770 & 3494 & -1921285 & 221026 & 3056\\
-1927488 & 158045 & 2636 & -1931627 & 94895 & 2239 & -1933698 & 31643 & 1862\\
-1933698 & -31643 & 1497 & -1931627 & -94895 & 1146 & -1927488 & -158045 & 799\\
-1921285 & -221026 & 469 & -1913024 & -283770 & 138 & -1902716 & -346211 & -177\\
-1890369 & -408281 & -496 & -1875999 & -469913 & -804 & -1859618 & -531042 & -1117\\
-1841248 & -591603 & -1425 & -1820906 & -651531 & -1735 & -1798613 & -710760 & -2046\\
-1774394 & -769228 & -2354 & -1748277 & -826873 & -2663 & -1720285 & -883632 & -2973\\
-1690453 & -939446 & -3290 & -1658810 & -994252 & -3613 & -1625391 & -1047995 & -3941\\
-1590231 & -1100615 & -4270 & -1553368 & -1152057 & -4601 & -1514843 & -1202264 & -4933\\
-1474695 & -1251185 & -5276 & -1432967 & -1298766 & -5623 & -1389706 & -1344956 & -5977\\
-1344956 & -1389706 & -6332 & -1298766 & -1432967 & -6685 & -1251185 & -1474695 & -7038\\
-1202264 & -1514843 & -7399 & -1152057 & -1553368 & -7767 & -1100615 & -1590231 & -8140\\
-1047995 & -1625391 & -8506 & -994252 & -1658810 & -8866 & -939446 & -1690453 & -9218\\
-883632 & -1720285 & -9571 & -826873 & -1748277 & -9924 & -769228 & -1774394 & -10273\\
-710760 & -1798613 & -10606 & -651531 & -1820906 & -10923 & -591603 & -1841248 & -11241\\
-531042 & -1859618 & -11546 & -469913 & -1875999 & -11823 & -408281 & -1890369 & -12225\\
-346211 & -1902716 & -12281 & -283770 & -1913024 & -13037 & -221026 & -1921285 & -12682\\
-158045 & -1927488 & -14794 & -94895 & -1931627 & -10897 & -31643 & -1933698 & -46425\\
\end{VPoissonTable}
\begin{VChainTable}{3}{x}{40}{3349364905389}{1}{94611}
1727789 & 702301 & 1905814 & 840952 & 2075677 & 292770 & 7153919 & 4717567\\
2659559 & -4169045 & 2778512 & 61508 & 2885565 & 47893 & 2980263 & 37685\\
3062199 & 29875 & 3131021 & 23786 & 19467262 & 13212594 & 3228206 & -13139713\\
3186436 & 3120 & 3131021 & 2175 & 3062199 & 1404 & 2980263 & 759\\
2885565 & 234 & 2778512 & -176 & 2659559 & -503 & 2529220 & -774\\
2388048 & -980 & 2236651 & -1137 & 2075677 & -1274 & 1905814 & -1398\\
1727789 & -1531 & 1542367 & -1720 & 1350340 & -2032 & 1152529 & -2570\\
949785 & -4019 & 742973 & -5056 & 532979 & -24489 & 320703 & 21004\\
4075 & -2536305 & 11603 & 74667 & 17365 & 83984 & 20484 & 77668\\
20484 & 67628 & 17365 & 55434 & 11604 & 22288 & 4074 & 55035\\
\end{VChainTable}
\begin{VChainTable}{3}{y}{48}{480367989919}{1}{10081}
742973 & 1317749 & 949785 & 719438 & 1152529 & 211244 & 1350340 & 132611\\
1542367 & 78595 & 1727789 & 53138 & 1905814 & 38062 & 2075677 & 29011\\
2236651 & 23361 & 2388048 & 19821 & 2529220 & 17664 & 2659559 & 16447\\
2778512 & 15893 & 2885565 & 15820 & 2980263 & 16082 & 3062199 & 16563\\
3131021 & 17168 & 16211108 & 11230618 & 3256154 & -11118909 & 3228206 & 16001\\
3186436 & 18042 & 3131021 & 22282 & 3062199 & 17911 & 2980263 & 16743\\
2885565 & 15368 & 2778512 & 14113 & 2659559 & 12888 & 2529220 & 11695\\
2388048 & 10543 & 2236651 & 9427 & 2075677 & 8335 & 1905814 & 7258\\
1727789 & 6174 & 1542367 & 5040 & 1350340 & 3771 & 1152529 & 2252\\
949785 & -308 & 742973 & -2384 & 532979 & -25690 & 320703 & 25964\\
4075 & -2926999 & 11603 & 86175 & 17365 & 96942 & 20484 & 89662\\
20484 & 78075 & 17365 & 64003 & 11604 & 25740 & 4074 & 63540\\
\end{VChainTable}
\begin{VChainTable}{3}{z\in[-q,-3/5]}{62}{-150000000000000}{0}{0}
9779 & 292281 & 27847 & -29820 & 41677 & -12082 & 49162 & -30034\\
49161 & -36633 & 41677 & -42073 & 27847 & -45478 & 9779 & -40415\\
769688 & 1372610 & 1279149 & -13478 & 1783135 & 10100 & 2279483 & -1453\\
2766072 & -3165 & 3240815 & -4980 & 3701681 & -6421 & 4146695 & -7731\\
4573952 & -8963 & 4981624 & -10270 & 5367963 & -11546 & 5731316 & -12864\\
6070127 & -14287 & 6382943 & -15728 & 6668428 & -17236 & 6925357 & -18871\\
7152631 & -20584 & 7349276 & -22367 & 7514451 & -24276 & 7647446 & -26320\\
7747696 & -28470 & 7814767 & -30740 & 7848376 & -33156 & 7848376 & -35715\\
7814767 & -38394 & 7747696 & -41196 & 7647446 & -44117 & 7514451 & -47142\\
7349276 & -50225 & 7152631 & -53312 & 6925357 & -56337 & 6668428 & -59222\\
6382943 & -61852 & 6070127 & -64064 & 5731316 & -65661 & 5367963 & -66430\\
4981624 & -66146 & 4573952 & -64554 & 4146695 & -61406 & 3701681 & -56550\\
3240815 & -49983 & 2766072 & -41820 & 2279483 & -32497 & 1783135 & -22656\\
1279149 & -12806 & 769688 & -3864 & 9779 & -9271 & 27847 & 259\\
41677 & 267 & 49161 & 216 & 49162 & 143 & 41677 & 66\\
27847 & -38 & 9779 & 70 &  &  &  & \\
\end{VChainTable}
\begin{VPoissonTable}{3}{143}{72}
40400305 & 43134967 & 250203 & -1432967 & 1298766 & -21785 & -1474695 & 1251185 & 48302\\
-1514843 & 1202264 & 21180 & -1553368 & 1152057 & 23333 & -1590231 & 1100615 & 18518\\
-1625391 & 1047995 & 16603 & -1658810 & 994252 & 14511 & -1690453 & 939446 & 12850\\
-1720285 & 883632 & 11458 & -1748277 & 826873 & 10202 & -1774394 & 769228 & 9118\\
-1798613 & 710760 & 8165 & -1820906 & 651531 & 7310 & -1841248 & 591603 & 6538\\
-1859618 & 531042 & 5829 & -1875999 & 469913 & 5175 & -1890369 & 408281 & 4582\\
-1902716 & 346211 & 4037 & -1913024 & 283770 & 3532 & -1921285 & 221026 & 3054\\
-1927488 & 158045 & 2599 & -1931627 & 94895 & 2172 & -1933698 & 31643 & 1772\\
-1933698 & -31643 & 1396 & -1931627 & -94895 & 1034 & -1927488 & -158045 & 681\\
-1921285 & -221026 & 341 & -1913024 & -283770 & 14 & -1902716 & -346211 & -296\\
-1890369 & -408281 & -598 & -1875999 & -469913 & -896 & -1859618 & -531042 & -1188\\
-1841248 & -591603 & -1477 & -1820906 & -651531 & -1756 & -1798613 & -710760 & -2027\\
-1774394 & -769228 & -2293 & -1748277 & -826873 & -2557 & -1720285 & -883632 & -2820\\
-1690453 & -939446 & -3082 & -1658810 & -994252 & -3340 & -1625391 & -1047995 & -3593\\
-1590231 & -1100615 & -3842 & -1553368 & -1152057 & -4087 & -1514843 & -1202264 & -4334\\
-1474695 & -1251185 & -4579 & -1432967 & -1298766 & -4823 & -1389706 & -1344956 & -5064\\
-1344956 & -1389706 & -5299 & -1298766 & -1432967 & -5529 & -1251185 & -1474695 & -5755\\
-1202264 & -1514843 & -5979 & -1152057 & -1553368 & -6203 & -1100615 & -1590231 & -6422\\
-1047995 & -1625391 & -6637 & -994252 & -1658810 & -6844 & -939446 & -1690453 & -7047\\
-883632 & -1720285 & -7244 & -826873 & -1748277 & -7435 & -769228 & -1774394 & -7621\\
-710760 & -1798613 & -7803 & -651531 & -1820906 & -7974 & -591603 & -1841248 & -8148\\
-531042 & -1859618 & -8323 & -469913 & -1875999 & -8461 & -408281 & -1890369 & -8693\\
-346211 & -1902716 & -8754 & -283770 & -1913024 & -9194 & -221026 & -1921285 & -8885\\
-158045 & -1927488 & -11037 & -94895 & -1931627 & -4863 & -31643 & -1933698 & -49270\\
\end{VPoissonTable}
\begin{VChainTable}{4}{x}{40}{3349364905389}{1}{94611}
1727789 & 595630 & 1905814 & 908738 & 2075677 & 295753 & 7153919 & 4676244\\
2659559 & -4111317 & 2778512 & 64043 & 2885565 & 50252 & 2980263 & 39935\\
3062199 & 32026 & 6317457 & 2623978 & 16280826 & 8015569 & 3228206 & -10528099\\
3186436 & 4690 & 3131021 & 3670 & 3062199 & 2817 & 2980263 & 2109\\
2885565 & 1520 & 2778512 & 1028 & 2659559 & 618 & 2529220 & 284\\
2388048 & 1 & 2236651 & -241 & 2075677 & -447 & 1905814 & -646\\
1727789 & -864 & 1542367 & -1122 & 1350340 & -1514 & 1152529 & -2140\\
949785 & -3666 & 742973 & -4797 & 532979 & -24466 & 320703 & 21337\\
4075 & -2557168 & 11603 & 75281 & 17365 & 84677 & 20484 & 78311\\
20484 & 68186 & 17365 & 55892 & 11604 & 22471 & 4074 & 55491\\
\end{VChainTable}
\begin{VChainTable}{4}{y}{48}{480367989919}{1}{10081}
742973 & 1297307 & 949785 & 731656 & 1152529 & 212004 & 1350340 & 134240\\
1542367 & 79909 & 1727789 & 54472 & 1905814 & 39425 & 2075677 & 30411\\
2236651 & 24786 & 2388048 & 21259 & 2529220 & 19082 & 2659559 & 17816\\
2778512 & 17179 & 2885565 & 16977 & 2980263 & 17068 & 3062199 & 17341\\
3131021 & 17699 & 16211108 & 11268944 & 3256154 & -11159498 & 3228206 & 15339\\
3186436 & 16907 & 3131021 & 20352 & 3062199 & 16539 & 2980263 & 15453\\
2885565 & 14206 & 2778512 & 13060 & 2659559 & 11943 & 2529220 & 10864\\
2388048 & 9813 & 2236651 & 8788 & 2075677 & 7786 & 1905814 & 6791\\
1727789 & 5772 & 1542367 & 4708 & 1350340 & 3496 & 1152529 & 2023\\
949785 & -476 & 742973 & -2516 & 532979 & -25732 & 320703 & 25849\\
4075 & -2920506 & 11603 & 85984 & 17365 & 96728 & 20484 & 89461\\
20484 & 77901 & 17365 & 63859 & 11604 & 25679 & 4074 & 63399\\
\end{VChainTable}
\begin{VChainTable}{4}{z\in[-q,-3/5]}{62}{-200000000000000}{0}{0}
13853 & 318383 & 39451 & -32482 & 59042 & -13160 & 69646 & -32714\\
69645 & -39907 & 59042 & -45825 & 39451 & -49544 & 13853 & -44029\\
1090391 & 1495445 & 1812128 & -14257 & 2526108 & 11672 & 3229268 & -732\\
3918601 & -2319 & 4591155 & -4113 & 5244047 & -5442 & 5874485 & -6637\\
6479766 & -7753 & 7057301 & -8919 & 7604614 & -10088 & 8119364 & -11284\\
8599346 & -12563 & 9042503 & -13903 & 9446940 & -15303 & 9810923 & -16812\\
10132893 & -18430 & 10411475 & -20146 & 10645471 & -21989 & 10833883 & -23979\\
10975902 & -26118 & 11070921 & -28412 & 11118532 & -30883 & 11118532 & -33553\\
11070921 & -36424 & 10975902 & -39500 & 10833883 & -42788 & 10645471 & -46299\\
10411475 & -50015 & 10132893 & -53909 & 9810923 & -57943 & 9446940 & -62053\\
9042503 & -66147 & 8599346 & -70081 & 8119364 & -73647 & 7604614 & -76580\\
7057301 & -78542 & 6479766 & -79118 & 5874485 & -77858 & 5244047 & -74310\\
4591155 & -68112 & 3918601 & -59129 & 3229268 & -47655 & 2526108 & -34426\\
1812128 & -20681 & 1090391 & -6612 & 13853 & -11878 & 39451 & 348\\
59042 & 388 & 69645 & 335 & 69646 & 249 & 59042 & 139\\
39451 & -29 & 13853 & 118 &  &  &  & \\
\end{VChainTable}
\begin{VPoissonTable}{4}{141}{71}
38967338 & 44433733 & 23084 & -1474695 & 1251185 & 211005 & -1514843 & 1202264 & -20189\\
-1553368 & 1152057 & 43666 & -1590231 & 1100615 & 18891 & -1625391 & 1047995 & 20970\\
-1658810 & 994252 & 16621 & -1690453 & 939446 & 14911 & -1720285 & 883632 & 12991\\
-1748277 & 826873 & 11496 & -1774394 & 769228 & 10220 & -1798613 & 710760 & 9078\\
-1820906 & 651531 & 8088 & -1841248 & 591603 & 7202 & -1859618 & 531042 & 6411\\
-1875999 & 469913 & 5700 & -1890369 & 408281 & 5051 & -1902716 & 346211 & 4451\\
-1913024 & 283770 & 3898 & -1921285 & 221026 & 3388 & -1927488 & 158045 & 2917\\
-1931627 & 94895 & 2476 & -1933698 & 31643 & 2057 & -1933698 & -31643 & 1662\\
-1931627 & -94895 & 1289 & -1927488 & -158045 & 940 & -1921285 & -221026 & 608\\
-1913024 & -283770 & 286 & -1902716 & -346211 & -26 & -1890369 & -408281 & -325\\
-1875999 & -469913 & -609 & -1859618 & -531042 & -881 & -1841248 & -591603 & -1144\\
-1820906 & -651531 & -1405 & -1798613 & -710760 & -1659 & -1774394 & -769228 & -1905\\
-1748277 & -826873 & -2142 & -1720285 & -883632 & -2370 & -1690453 & -939446 & -2594\\
-1658810 & -994252 & -2818 & -1625391 & -1047995 & -3037 & -1590231 & -1100615 & -3252\\
-1553368 & -1152057 & -3460 & -1514843 & -1202264 & -3660 & -1474695 & -1251185 & -3857\\
-1432967 & -1298766 & -4052 & -1389706 & -1344956 & -4243 & -1344956 & -1389706 & -4430\\
-1298766 & -1432967 & -4611 & -1251185 & -1474695 & -4790 & -1202264 & -1514843 & -4964\\
-1152057 & -1553368 & -5134 & -1100615 & -1590231 & -5295 & -1047995 & -1625391 & -5452\\
-994252 & -1658810 & -5606 & -939446 & -1690453 & -5760 & -883632 & -1720285 & -5912\\
-826873 & -1748277 & -6053 & -769228 & -1774394 & -6186 & -710760 & -1798613 & -6315\\
-651531 & -1820906 & -6439 & -591603 & -1841248 & -6573 & -531042 & -1859618 & -6706\\
-469913 & -1875999 & -6807 & -408281 & -1890369 & -6997 & -346211 & -1902716 & -7014\\
-283770 & -1913024 & -7453 & -221026 & -1921285 & -6970 & -158045 & -1927488 & -9502\\
-94895 & -1931627 & -1362 & -31643 & -1933698 & -54028 &  &  & \\
\end{VPoissonTable}
\begin{VChainTable}{5}{x}{39}{3349364905389}{1}{94611}
1727789 & 529775 & 1905814 & 950484 & 2075677 & 297565 & 7153919 & 4650255\\
2659559 & -4075339 & 2778512 & 65583 & 2885565 & 51687 & 2980263 & 41306\\
3062199 & 33350 & 22598283 & 15808801 & 3228206 & -15688094 & 3186436 & 5690\\
3131021 & 4629 & 3062199 & 3737 & 2980263 & 2986 & 2885565 & 2352\\
2778512 & 1820 & 2659559 & 1366 & 2529220 & 978 & 2388048 & 649\\
2236651 & 365 & 2075677 & 99 & 1905814 & -146 & 1727789 & -407\\
1542367 & -729 & 1350340 & -1165 & 1152529 & -1838 & 949785 & -3439\\
742973 & -4608 & 532979 & -24443 & 320703 & 21547 & 4075 & -2570516\\
11603 & 75674 & 17365 & 85119 & 20484 & 78722 & 20484 & 68542\\
17365 & 56189 & 11604 & 22591 & 4074 & 55779 &  & \\
\end{VChainTable}
\begin{VChainTable}{5}{y}{48}{480367989919}{1}{10081}
742973 & 1293462 & 949785 & 734356 & 1152529 & 212459 & 1350340 & 134919\\
1542367 & 80571 & 1727789 & 55179 & 1905814 & 40164 & 2075677 & 31175\\
2236651 & 25561 & 2388048 & 22033 & 2529220 & 19832 & 2659559 & 18510\\
2778512 & 17795 & 2885565 & 17490 & 2980263 & 17449 & 3062199 & 17561\\
3131021 & 17739 & 16211108 & 11293776 & 3256154 & -11187610 & 3228206 & 14672\\
3186436 & 15935 & 3131021 & 18854 & 3062199 & 15441 & 2980263 & 14424\\
2885565 & 13264 & 2778512 & 12207 & 2659559 & 11182 & 2529220 & 10177\\
2388048 & 9206 & 2236651 & 8262 & 2075677 & 7318 & 1905814 & 6389\\
1727789 & 5438 & 1542367 & 4408 & 1350340 & 3256 & 1152529 & 1832\\
949785 & -645 & 742973 & -2615 & 532979 & -25760 & 320703 & 25723\\
4075 & -2913154 & 11603 & 85771 & 17365 & 96482 & 20484 & 89231\\
20484 & 77704 & 17365 & 63701 & 11604 & 25616 & 4074 & 63235\\
\end{VChainTable}
\begin{VChainTable}{5}{z\in[-q,-3/5]}{62}{-250000000000000}{0}{0}
17928 & 334975 & 51054 & -34169 & 76407 & -13848 & 90129 & -34416\\
90129 & -41985 & 76408 & -48219 & 51054 & -52124 & 17928 & -46322\\
1411093 & 1573532 & 2345108 & -14716 & 3269080 & 12728 & 4179054 & -223\\
5071131 & -1687 & 5941494 & -3489 & 6786414 & -4692 & 7602274 & -5783\\
8385580 & -6855 & 9132977 & -7878 & 9841266 & -8937 & 10507413 & -10055\\
11128565 & -11205 & 11702063 & -12425 & 12225451 & -13736 & 12696488 & -15135\\
13113157 & -16626 & 13473673 & -18243 & 13776492 & -19996 & 14020319 & -21891\\
14204109 & -23934 & 14327074 & -26166 & 14388688 & -28606 & 14388688 & -31250\\
14327074 & -34124 & 14204109 & -37263 & 14020319 & -40672 & 13776492 & -44372\\
13473673 & -48364 & 13113157 & -52655 & 12696488 & -57239 & 12225451 & -62081\\
11702063 & -67096 & 11128565 & -72195 & 10507413 & -77209 & 9841266 & -81865\\
9132977 & -85820 & 8385580 & -88586 & 7602274 & -89520 & 6786414 & -87936\\
5941494 & -83107 & 5071131 & -74470 & 4179054 & -61970 & 3269080 & -46202\\
2345108 & -28901 & 1411093 & -10183 & 17928 & -14170 & 51054 & 432\\
76408 & 502 & 90129 & 464 & 90129 & 375 & 76407 & 239\\
51054 & -3 & 17928 & 179 &  &  &  & \\
\end{VChainTable}
\begin{VPoissonTable}{5}{137}{69}
35977800 & 46887182 & 112739 & -1553368 & 1152057 & 95323 & -1590231 & 1100615 & 9508\\
-1625391 & 1047995 & 30038 & -1658810 & 994252 & 18683 & -1690453 & 939446 & 17726\\
-1720285 & 883632 & 14825 & -1748277 & 826873 & 13099 & -1774394 & 769228 & 11485\\
-1798613 & 710760 & 10160 & -1820906 & 651531 & 9018 & -1841248 & 591603 & 8002\\
-1859618 & 531042 & 7122 & -1875999 & 469913 & 6321 & -1890369 & 408281 & 5610\\
-1902716 & 346211 & 4961 & -1913024 & 283770 & 4371 & -1921285 & 221026 & 3826\\
-1927488 & 158045 & 3330 & -1931627 & 94895 & 2859 & -1933698 & 31643 & 2432\\
-1933698 & -31643 & 2018 & -1931627 & -94895 & 1644 & -1927488 & -158045 & 1280\\
-1921285 & -221026 & 945 & -1913024 & -283770 & 622 & -1902716 & -346211 & 316\\
-1890369 & -408281 & 22 & -1875999 & -469913 & -257 & -1859618 & -531042 & -525\\
-1841248 & -591603 & -782 & -1820906 & -651531 & -1027 & -1798613 & -710760 & -1268\\
-1774394 & -769228 & -1496 & -1748277 & -826873 & -1723 & -1720285 & -883632 & -1932\\
-1690453 & -939446 & -2147 & -1658810 & -994252 & -2344 & -1625391 & -1047995 & -2548\\
-1590231 & -1100615 & -2735 & -1553368 & -1152057 & -2925 & -1514843 & -1202264 & -3102\\
-1474695 & -1251185 & -3278 & -1432967 & -1298766 & -3448 & -1389706 & -1344956 & -3613\\
-1344956 & -1389706 & -3778 & -1298766 & -1432967 & -3932 & -1251185 & -1474695 & -4087\\
-1202264 & -1514843 & -4232 & -1152057 & -1553368 & -4376 & -1100615 & -1590231 & -4514\\
-1047995 & -1625391 & -4649 & -994252 & -1658810 & -4780 & -939446 & -1690453 & -4904\\
-883632 & -1720285 & -5027 & -826873 & -1748277 & -5144 & -769228 & -1774394 & -5260\\
-710760 & -1798613 & -5371 & -651531 & -1820906 & -5478 & -591603 & -1841248 & -5587\\
-531042 & -1859618 & -5696 & -469913 & -1875999 & -5788 & -408281 & -1890369 & -5954\\
-346211 & -1902716 & -5957 & -283770 & -1913024 & -6390 & -221026 & -1921285 & -5782\\
-158045 & -1927488 & -8662 & -94895 & -1931627 & 1153 & -31643 & -1933698 & -58970\\
\end{VPoissonTable}
\begin{VChainTable}{6}{x}{41}{3349364905389}{1}{94611}
1727789 & 485088 & 1905814 & 978778 & 2075677 & 298775 & 7153919 & 4632391\\
2659559 & -4050763 & 2778512 & 66616 & 2885565 & 52657 & 2980263 & 42225\\
3062199 & 34237 & 6317457 & 2619041 & 13024672 & 5366428 & 3256154 & -7866090\\
3228206 & 7669 & 3186436 & 6386 & 3131021 & 5296 & 3062199 & 4377\\
2980263 & 3607 & 2885565 & 2942 & 2778512 & 2370 & 2659559 & 1892\\
2529220 & 1477 & 2388048 & 1106 & 2236651 & 791 & 2075677 & 500\\
1905814 & 208 & 1727789 & -86 & 1542367 & -433 & 1350340 & -918\\
1152529 & -1628 & 949785 & -3267 & 742973 & -4482 & 532979 & -24417\\
320703 & 21698 & 4075 & -2579794 & 11603 & 75952 & 17365 & 85421\\
20484 & 79011 & 20484 & 68789 & 17365 & 56392 & 11604 & 22674\\
4074 & 55984 &  &  &  &  &  & \\
\end{VChainTable}
\begin{VChainTable}{6}{y}{47}{480367989919}{1}{10081}
742973 & 1296301 & 949785 & 733184 & 1152529 & 212765 & 1350340 & 135179\\
1542367 & 80931 & 1727789 & 55581 & 1905814 & 40597 & 2075677 & 31626\\
2236651 & 26018 & 2388048 & 22480 & 2529220 & 20254 & 2659559 & 18884\\
2778512 & 18096 & 2885565 & 17699 & 2980263 & 17554 & 3062199 & 17549\\
3131021 & 17593 & 19467262 & 14176287 & 3228206 & -14059423 & 3186436 & 15114\\
3131021 & 17661 & 3062199 & 14540 & 2980263 & 13587 & 2885565 & 12500\\
2778512 & 11504 & 2659559 & 10551 & 2529220 & 9617 & 2388048 & 8700\\
2236651 & 7814 & 2075677 & 6932 & 1905814 & 6044 & 1727789 & 5138\\
1542367 & 4164 & 1350340 & 3049 & 1152529 & 1654 & 949785 & -776\\
742973 & -2713 & 532979 & -25781 & 320703 & 25619 & 4075 & -2906127\\
11603 & 85567 & 17365 & 96244 & 20484 & 89020 & 20484 & 77512\\
17365 & 63545 & 11604 & 25555 & 4074 & 63083 &  & \\
\end{VChainTable}
\begin{VChainTable}{6}{z\in[-q,-3/5]}{61}{-300000000000000}{0}{0}
22002 & 346514 & 62657 & -35341 & 93773 & -14323 & 110613 & -35605\\
110613 & -43431 & 93773 & -49880 & 62657 & -53916 & 22003 & -47931\\
1731796 & 1627865 & 2878087 & -15016 & 4012053 & 13462 & 5128838 & 188\\
6223661 & -1245 & 7291834 & -2974 & 8328781 & -4127 & 9330064 & -5154\\
10291393 & -6126 & 11208654 & -7081 & 12077917 & -8049 & 12895461 & -9068\\
13657784 & -10143 & 14361623 & -11261 & 15003963 & -12471 & 15582054 & -13783\\
16093419 & -15174 & 16535871 & -16684 & 16907514 & -18342 & 17206755 & -20137\\
17432315 & -22084 & 17583227 & -24231 & 17658845 & -26586 & 17658845 & -29163\\
17583227 & -31986 & 17432315 & -35103 & 17206755 & -38522 & 16907514 & -42269\\
16535871 & -46374 & 16093419 & -50853 & 15582054 & -55732 & 15003963 & -60981\\
14361623 & -66571 & 13657784 & -72445 & 12895461 & -78479 & 12077917 & -84431\\
11208654 & -89984 & 10291393 & -94640 & 9330064 & -97688 & 8328781 & -98254\\
7291834 & -95287 & 6223661 & -87780 & 5128838 & -75149 & 4012053 & -57661\\
2878087 & -37220 & 1731796 & -14240 & 22003 & -16337 & 156430 & 2342\\
110613 & -636 & 110613 & 514 & 93773 & 360 & 62657 & 43\\
22002 & 248 &  &  &  &  &  & \\
\end{VChainTable}
\begin{VPoissonTable}{6}{135}{68}
34424432 & 48039239 & 28503 & -1590231 & 1100615 & 177007 & -1625391 & 1047995 & -15336\\
-1658810 & 994252 & 37497 & -1690453 & 939446 & 16709 & -1720285 & 883632 & 18252\\
-1748277 & 826873 & 14532 & -1774394 & 769228 & 12997 & -1798613 & 710760 & 11310\\
-1820906 & 651531 & 10007 & -1841248 & 591603 & 8861 & -1859618 & 531042 & 7858\\
-1875999 & 469913 & 6979 & -1890369 & 408281 & 6194 & -1902716 & 346211 & 5498\\
-1913024 & 283770 & 4861 & -1921285 & 221026 & 4284 & -1927488 & 158045 & 3746\\
-1931627 & 94895 & 3258 & -1933698 & 31643 & 2807 & -1933698 & -31643 & 2392\\
-1931627 & -94895 & 2001 & -1927488 & -158045 & 1631 & -1921285 & -221026 & 1282\\
-1913024 & -283770 & 955 & -1902716 & -346211 & 650 & -1890369 & -408281 & 359\\
-1875999 & -469913 & 82 & -1859618 & -531042 & -186 & -1841248 & -591603 & -439\\
-1820906 & -651531 & -681 & -1798613 & -710760 & -911 & -1774394 & -769228 & -1133\\
-1748277 & -826873 & -1349 & -1720285 & -883632 & -1553 & -1690453 & -939446 & -1756\\
-1658810 & -994252 & -1945 & -1625391 & -1047995 & -2130 & -1590231 & -1100615 & -2310\\
-1553368 & -1152057 & -2483 & -1514843 & -1202264 & -2652 & -1474695 & -1251185 & -2815\\
-1432967 & -1298766 & -2971 & -1389706 & -1344956 & -3123 & -1344956 & -1389706 & -3266\\
-1298766 & -1432967 & -3412 & -1251185 & -1474695 & -3549 & -1202264 & -1514843 & -3683\\
-1152057 & -1553368 & -3812 & -1100615 & -1590231 & -3934 & -1047995 & -1625391 & -4059\\
-994252 & -1658810 & -4170 & -939446 & -1690453 & -4289 & -883632 & -1720285 & -4390\\
-826873 & -1748277 & -4502 & -769228 & -1774394 & -4598 & -710760 & -1798613 & -4702\\
-651531 & -1820906 & -4798 & -591603 & -1841248 & -4891 & -531042 & -1859618 & -5000\\
-469913 & -1875999 & -5065 & -408281 & -1890369 & -5245 & -346211 & -1902716 & -5207\\
-283770 & -1913024 & -5685 & -221026 & -1921285 & -4933 & -158045 & -1927488 & -8166\\
-94895 & -1931627 & 3172 & -31643 & -1933698 & -63714 &  &  & \\
\end{VPoissonTable}
\begin{VChainTable}{7}{x}{39}{3349364905389}{1}{94611}
1727789 & 452787 & 1905814 & 999210 & 2075677 & 299640 & 7153919 & 4619366\\
2659559 & -4032921 & 2778512 & 67361 & 2885565 & 53357 & 2980263 & 42885\\
3062199 & 34872 & 22598283 & 15799919 & 3228206 & -15668266 & 3186436 & 6897\\
3131021 & 5792 & 3062199 & 4853 & 2980263 & 4060 & 2885565 & 3377\\
2778512 & 2781 & 2659559 & 2282 & 2529220 & 1849 & 2388048 & 1454\\
2236651 & 1107 & 2075677 & 797 & 1905814 & 480 & 1727789 & 151\\
1542367 & -218 & 1350340 & -728 & 1152529 & -1469 & 949785 & -3133\\
742973 & -4382 & 532979 & -24397 & 320703 & 21806 & 4075 & -2586615\\
11603 & 76150 & 17365 & 85655 & 20484 & 79215 & 20484 & 68976\\
17365 & 56540 & 11604 & 22733 & 4074 & 56134 &  & \\
\end{VChainTable}
\begin{VChainTable}{7}{y}{47}{480367989919}{1}{10081}
742973 & 1302156 & 949785 & 730240 & 1152529 & 212980 & 1350340 & 135224\\
1542367 & 81118 & 1727789 & 55813 & 1905814 & 40862 & 2075677 & 31904\\
2236651 & 26297 & 2388048 & 22743 & 2529220 & 20485 & 2659559 & 19072\\
2778512 & 18224 & 2885565 & 17752 & 2980263 & 17520 & 3062199 & 17419\\
3131021 & 17363 & 19467262 & 14190705 & 3228206 & -14077572 & 3186436 & 14420\\
3131021 & 16690 & 3062199 & 13797 & 2980263 & 12896 & 2885565 & 11868\\
2778512 & 10919 & 2659559 & 10023 & 2529220 & 9147 & 2388048 & 8273\\
2236651 & 7435 & 2075677 & 6607 & 1905814 & 5755 & 1727789 & 4889\\
1542367 & 3946 & 1350340 & 2863 & 1152529 & 1514 & 949785 & -901\\
742973 & -2800 & 532979 & -25782 & 320703 & 25499 & 4075 & -2899676\\
11603 & 85372 & 17365 & 96033 & 20484 & 88821 & 20484 & 77342\\
17365 & 63401 & 11604 & 25497 & 4074 & 62944 &  & \\
\end{VChainTable}
\begin{VChainTable}{7}{z\in[-q,-3/5]}{62}{-350000000000000}{0}{0}
26077 & 355036 & 74260 & -36213 & 111139 & -14672 & 131096 & -36481\\
131097 & -44500 & 111138 & -51099 & 74261 & -55251 & 26077 & -49110\\
2052499 & 1667978 & 3411066 & -15218 & 4755026 & 14007 & 6078623 & 506\\
7376191 & -899 & 8642173 & -2580 & 9871148 & -3682 & 11057853 & -4644\\
12197207 & -5553 & 13284331 & -6438 & 14314568 & -7344 & 15283509 & -8291\\
16187004 & -9277 & 17021182 & -10328 & 17782475 & -11458 & 18467619 & -12673\\
19073683 & -13983 & 19598069 & -15410 & 20038535 & -16971 & 20393191 & -18668\\
20660521 & -20533 & 20839381 & -22592 & 20929001 & -24854 & 20929001 & -27337\\
20839381 & -30104 & 20660521 & -33160 & 20393191 & -36528 & 20038535 & -40266\\
19598069 & -44397 & 19073683 & -48966 & 18467619 & -53977 & 17782475 & -59461\\
17021182 & -65413 & 16187004 & -71803 & 15283509 & -78539 & 14314568 & -85451\\
13284331 & -92242 & 12197207 & -98450 & 11057853 & -103377 & 9871148 & -106019\\
8642173 & -105092 & 7376191 & -99170 & 6078623 & -87102 & 4755026 & -68605\\
3411066 & -45520 & 2052499 & -18587 & 26077 & -18457 & 74261 & 581\\
111138 & 702 & 131097 & 714 & 131096 & 659 & 111139 & 497\\
74260 & 106 & 26077 & 327 &  &  &  & \\
\end{VChainTable}
\begin{VPoissonTable}{7}{131}{66}
31208810 & 50187849 & 185535 & -1658810 & 994252 & -8205 & -1690453 & 939446 & 34749\\
-1720285 & 883632 & 17141 & -1748277 & 826873 & 17765 & -1774394 & 769228 & 14285\\
-1798613 & 710760 & 12707 & -1820906 & 651531 & 11051 & -1841248 & 591603 & 9767\\
-1859618 & 531042 & 8635 & -1875999 & 469913 & 7659 & -1890369 & 408281 & 6797\\
-1902716 & 346211 & 6035 & -1913024 & 283770 & 5358 & -1921285 & 221026 & 4734\\
-1927488 & 158045 & 4175 & -1931627 & 94895 & 3651 & -1933698 & 31643 & 3180\\
-1933698 & -31643 & 2747 & -1931627 & -94895 & 2342 & -1927488 & -158045 & 1969\\
-1921285 & -221026 & 1607 & -1913024 & -283770 & 1277 & -1902716 & -346211 & 960\\
-1890369 & -408281 & 670 & -1875999 & -469913 & 392 & -1859618 & -531042 & 127\\
-1841248 & -591603 & -125 & -1820906 & -651531 & -369 & -1798613 & -710760 & -597\\
-1774394 & -769228 & -814 & -1748277 & -826873 & -1023 & -1720285 & -883632 & -1223\\
-1690453 & -939446 & -1415 & -1658810 & -994252 & -1604 & -1625391 & -1047995 & -1781\\
-1590231 & -1100615 & -1955 & -1553368 & -1152057 & -2117 & -1514843 & -1202264 & -2277\\
-1474695 & -1251185 & -2433 & -1432967 & -1298766 & -2581 & -1389706 & -1344956 & -2728\\
-1344956 & -1389706 & -2861 & -1298766 & -1432967 & -2997 & -1251185 & -1474695 & -3120\\
-1202264 & -1514843 & -3250 & -1152057 & -1553368 & -3367 & -1100615 & -1590231 & -3487\\
-1047995 & -1625391 & -3595 & -994252 & -1658810 & -3705 & -939446 & -1690453 & -3810\\
-883632 & -1720285 & -3909 & -826873 & -1748277 & -4007 & -769228 & -1774394 & -4099\\
-710760 & -1798613 & -4195 & -651531 & -1820906 & -4282 & -591603 & -1841248 & -4376\\
-531042 & -1859618 & -4469 & -469913 & -1875999 & -4536 & -408281 & -1890369 & -4698\\
-346211 & -1902716 & -4664 & -283770 & -1913024 & -5147 & -221026 & -1921285 & -4311\\
-158045 & -1927488 & -7829 & -94895 & -1931627 & 4850 & -31643 & -1933698 & -68165\\
\end{VPoissonTable}

\section{Exact certificate algorithm}
\label{sec:certificate-algorithm}
All numerical bounds in this appendix are computer-assisted.  The code and
logs supplied as ancillary files use outward-rounded directed interval
arithmetic; the rules below specify the calculation independently of those
files.  All inputs are the printed tables. Terminating decimals and fractions are
parsed as exact rationals; a missing or duplicate field, invalid mass, failed
strict inequality, or unresolved cell is fatal. No unprinted numerical data
are part of the argument.

\paragraph{Shared quadratic duality.}
Let $Q(\boldsymbol\mu)=B(\boldsymbol\mu,\boldsymbol\mu)$ be a quadratic
form with symmetric polarization $B$, and let
$\mathcal E(\boldsymbol\mu)=C+Q(\boldsymbol\mu)+L(\boldsymbol\mu)$, where
$L$ is linear. For a fixed reference $\boldsymbol\nu$, let $F_i^\nu$ be
the component potentials representing
$L(\boldsymbol\mu)+2B(\boldsymbol\nu,\boldsymbol\mu)$, and assume
$\mu_i$ and $\nu_i$ have the same prescribed mass in each component. With
$\boldsymbol\sigma=\boldsymbol\mu-\boldsymbol\nu$, polarization gives the
exact identity
\[
 \mathcal E(\boldsymbol\mu)=C-Q(\boldsymbol\nu)
       +\sum_i\int F_i^\nu\,d\mu_i+Q(\boldsymbol\sigma).
\]
Thus a nonnegative (respectively nonpositive) zero-mass remainder gives the
lower (respectively upper) dual bound. The P0 and V proofs below differ only
in the sign proof for this remainder and in their explicit transfer constants.

\paragraph{P0 endpoint measures.}
Divide each printed integer mass by its table denominator. On every pinned or
frequency interval $[b_{i-1},b_i]$, use the uniform measure of that mass.
The contour has $N=64$ nodes
\[
 \theta_j=\pi(2j+1)/64,\quad c_{63-j}=c_j,
 \qquad F_{48}(\theta)=\frac1{49}\left|\sum_{k=0}^{48}e^{ik\theta}\right|^2.
\]
The nine (respectively eight) printed pair masses at $q=2$ (respectively
$q=11$) occupy the first nodes; zero-fill the rest of the first half, then
reflect. The contour density is $(2\pi)^{-1}\sum_jc_jF_{48}(\theta-\theta_j)$.
The algorithm verifies exact total masses, nonnegative weights, interval
density bounds, and the printed ellipse $(\xi,A,B)$. For an interval row use
the primitive $H(u)=u(\log u-1)$ with $H(0)=0$; for the contour use the
finite Fourier/Joukowski formulas below.  Put
\[
 c=\frac{A+B}{2},\qquad d=\frac{A-B}{2},\qquad \rho=\frac dc,\qquad
 z(u)=\xi+ce^{iu}+de^{-iu},
\]
and
\[
 m_n=\frac{49-n}{49}\sum_{j=0}^{63}c_j\cos(n\theta_j),\qquad
 \gamma_n=\frac{(1+\rho^n)m_n}{n}\quad(1\le n\le48).
\]
These are exact algebraic numbers in the real subfield generated by
$e^{\pi i/64}$ and are enclosed by directed algebraic (or Arb) intervals;
they are not assumed to be rational.  On the contour,
\[
 U_\nu(z(u))=M\log c-\sum_{n=1}^{48}\gamma_n\cos(nu).
\]
For a real point inside the ellipse set \(w=(t-\xi)/c\) and
\(P_0(w)=2,\ P_1(w)=w,\ P_n(w)=wP_{n-1}(w)-\rho P_{n-2}(w)\).  Then
\[
 U_\nu(t)=M\log c-\sum_{n=1}^{48}\frac{m_n}{n}P_n(w).
\]
For a real point $-t$ on the negative axis to the left of the ellipse, let
\[
 u=t+\xi,\qquad
 \Delta_t=\sqrt{u^2-4cd},\qquad
 s_t=-\frac{2c}{u+\Delta_t},
\]
where the positive real square root is chosen and \(|s_t|<1\).  The exterior
formula is
\[
 U_\nu(-t)=M\log\frac{u+\Delta_t}{2}
          -\sum_{n=1}^{48}\gamma_n s_t^n.
\]
The same formulas are used for interval boxes, with square roots and
logarithms evaluated by outward rounding.  For a complex contour point and a
real interval, use
\[
 \mathcal I(z;a,b)=\Re\{\mathcal L(z-a)-\mathcal L(z-b)\},\qquad
 \mathcal L(y)=y(\log y-1),
\]
with the logarithm continued from the upper half-plane; the lower half is
obtained by conjugation.  On a real crossing, split at the crossing and use
$H(u)=u(\log|u|-1)$, $H(0)=0$.  Thus all mixed contour--interval
potentials, their primitives, and the contour self-energy
\[
 I(\nu,\nu)=M^2\log c-\sum_{n=1}^{48}\frac{(1+\rho^n)m_n^2}{n}
\]
are finite expressions in the printed masses and geometry.  No unprinted
recurrence or rationality assertion is used.

\paragraph{\textbf{V reference tables.}}
The V data in Section~\ref{sec:v-reference-data} are printed as signed decimal
integers in ordinary tables. In each real chain, the header gives the initial
endpoint in units of $1/(5\cdot10^{13})$; each displayed length is in units of
$10^{-8}$, except for the single residual recorded by $(j,r)$ in the header,
which adds $r/(5\cdot10^{13})$ to length $j$ (the pair $(0,0)$ means no
correction). Mass increments are in units of $10^{-8}$. Starting from zero
mass, accumulate each pair left to right and top to bottom. In each Poisson
prefix, all three columns are successive
increments of the real coordinate, imaginary coordinate, and mass, in units
of $10^{-8}$. The table header gives the prefix and full row counts; these must
agree with Table~\ref{tab:v-certificates}. Complete the Poisson list by adjoining the
reversed conjugates of all prefix rows except its final real row. The verification
checks every count, increment, support, total mass, and conjugate pair. Real
rows are uniform. A Poisson row
is harmonic measure on $|z|=r=3/5$, with density
$w(r^2-|c|^2)/(2\pi|re^{i\theta}-c|^2)$ and potential
$w\log|r^2-z\bar c|-w\log r$ inside, $w\log|z-c|$ outside, and their common
boundary value. It requires the exact strict test $|c|^2<r^2$.

\paragraph{\textbf{Directed interval evaluation.}}
For a real interval row the potential is evaluated by $H$; pair integrals use
$G(u)=\frac12u^2\log|u|-\frac34u^2$, $G(0)=0$. The V denominator field uses
\[
 \int_a^b\log(q+xt)\,dt=
 \frac{(q+xb)(\log(q+xb)-1)-(q+xa)(\log(q+xa)-1)}x,
\]
with value $(b-a)\log q$ at $x=0$. Its energy series for $\log(q+xv)$ ends
at degree $420$; the absolute tail is at most
$2M_xM_v/(421q^{421})$, since $0\le xv/q\le1/2$.

\begin{center}
\begin{tabular}{@{}>{\raggedright\arraybackslash}p{0.18\linewidth}p{0.33\linewidth}p{0.41\linewidth}@{}}
\hline
check & initial cover and evaluator & deterministic acceptance rule\\
\hline
P0 contour & $u/\pi\in[j/256,(j+1)/256]$, $0\le j<256$; Arb, 45 digits; upper half only, lower by exact conjugation & lower field endpoint $\ge h_\eta$; otherwise rational midpoint bisection, left first, depth at most 10\\
P0 integrals & 8192 equal cells on $[1/4,5/4]$ and 16384 on $[1/2,81/10]$ & each lower field bound gives $|I|(h-W_-)_{+}$; the printed tail test uses $W_\mu'>1/50$\\
V endpoint boxes & $0\le x,v\le\delta$, $\delta=1/250$; use the printed $E_{x,q},E_{v,q}$ & accept the two boxes when $E_{x,q}<K_{x,q}$ and $E_{v,q}<K_{v,q}$\\
V real supports & $x,v\in[\delta,1]$, $z\in[-q,-3/5]$ and $[3/5,1]$; Arb, 90 digits & accept iff the directed upper endpoint is strictly below its printed threshold; otherwise bisect the rational midpoint, left first, to depth 30\\
V circle bounds & each half-circle parameter $z=(3/5)e^{i\pi t}$, $t\in[0,1]$; Arb, 90 digits & same strict threshold test and bisection; circle endpoints are checked directly\\
\hline
\end{tabular}
\end{center}
The V supports cover the entire admissible set.  The lower circle follows by
conjugation.  The two real endpoint boxes are discharged by the displayed
endpoint inequalities rather than by evaluating $\log 0$.  Nonfinite balls,
unvisited endpoints, or unresolved cells are rejected.
Finally recompute $Q_q$ and require the directed lower endpoint of
\[
 Q_{q,-}-\frac32\log q-K_{x,q}-\frac54K_{v,q}-K_{z,q}
\]
to exceed the printed rational $\alpha_q$ for every $q=2,\ldots,7$.

All quantities needed for these evaluations are printed in the preceding
tables or defined by the displayed formulas.  The leaf intervals generated by
the bisection rule are auxiliary computation, not additional data.

\end{document}
